
\section{From a charged material to a certified finite interface}
\label{sec:charged-material}

A material action can retain the charged fields rather than replace
them by neutral oscillators. The operator boson--fermion correspondence
gives a convenient construction on a circle of circumference $L$.
For one Dirac flavor with antiperiodic spatial fermions, its current
oscillators and integer chiral charges have free energy
\[
 H_0=\frac\pi L(N_R^2+N_L^2)
       +\sum_{k\ge1}\frac{2\pi k}{L}(m_{R,k}+m_{L,k}).
\]
The current descendants of the charge vacua form a complete
orthonormal basis of the fermion Hilbert space
\cite[Sections~2.5--2.6]{LangmannMoosavi2015}.
The correspondence includes Klein factors and charged odd operators;
it is stronger than equality of partition functions.

For a scalar $\phi$ of period $\sqrt\pi$, put
$w=N_R+N_L$ and $n=N_R-N_L$. Its zero-mode description is
\[
 \bigoplus_{w\in\mathbb Z}
 L^2\!\left([0,\sqrt\pi],\,
       \psi(\phi_0+\sqrt\pi)=(-1)^w\psi(\phi_0)\right)
       \otimes\mathcal F_{\rm osc}.
\]
Here $p_0=\sqrt\pi n$, $n\equiv w\pmod2$, and the charge
energy is $\pi(w^2+n^2)/(2L)$. The state $w=1,n=0$ would have
$N_R=N_L=1/2$ and is excluded. An ordinary compact scalar with
periodic zero-mode wavefunctions has different sectors. Fixed-spin
bosonization requires this additional spin information
\cite[Section~3.5]{KarchTongTurner2019}.

Take two such flavors and real material fields $Z,\theta,y$.
The local classical coframe is
$d\phi_a+\lambda_a\theta f(Z)dZ$, together with
$m_ZdZ,\sqrt I\,d\theta,\sqrt J\,dy$, where
$m_Z^2,I,J>0$. Its positive canonical energy is
\begin{align*}
 H_{\rm mat}={}&\frac12\int\!\left[
 \sum_a\Pi_a^2+
 \frac{(\Pi_Z-\sum_a\lambda_a\theta f(Z)\Pi_a)^2}{m_Z^2}
 +\frac{\Pi_\theta^2}{I}+\frac{\Pi_y^2}{J}\right]dx\\
 &+\frac12\int\!\left[
 \sum_a(\phi_a'+\lambda_a\theta f(Z)Z')^2
 +m_Z^2Z'^2+I\theta'^2+Jy'^2+2V\right]dx .
\end{align*}
Assume smooth $f$ and a smooth nonnegative $V$ independent of
the scalar zero modes. A Fourier regulator retains modes
$0<|k|\le K$ of the interaction and material fields, all winding
and charge sectors, and the complete free oscillator tail above $K$.
The real Fourier coordinates are not separately wrapped angles.

\begin{proposition}[Positive charged material at finite resolution]
\label{prop:charged-material-regulator}
For each finite $K$, the displayed ordered sum of squares has a
nonnegative self-adjoint Friedrichs realization. With the free tail
retained, the known boson--fermion unitary transports it to the full
four-chiral-field CAR Hilbert space tensor the material space.
It strongly commutes with all four canonical chiral charges.
No small-coupling bound on the finite $\lambda_a$ is needed.
\end{proposition}
\begin{proof}
Work first in each winding sector on smooth zero-mode line-bundle
sections and compactly supported smooth wavefunctions of the finite
real coordinates. The kinetic factor
\[
 D_Z(x)=\Pi_Z(x)-\sum_a\lambda_a\theta(x)f(Z(x))\Pi_a(x)
\]
is symmetric: its extra coefficient is independent of every $\phi_a$
coordinate, so its formal divergence vanishes. The other momentum
factors are symmetric and the spatial factors are real multiplication
operators. Their integrated squares and $V$ define a densely defined
nonnegative symmetric operator on the common core. Its Friedrichs
extension exists. Finite sums of winding sectors give a dense core
for their direct sum. Add the positive free-tail Hamiltonian and use
the complete current/charge basis to transport this operator to the
fermion space. The smeared fermion operators themselves are transported,
not replaced by cutoff vertex exponentials, so their CAR are retained.
Translations of each $\phi_{a,0}$ leave the core and closed form
invariant, as do winding phases. The associated charge unitaries
therefore commute with the self-adjoint realization.
\end{proof}

Low-mode canonical squares retain their common zero-point constant;
the free tail uses one fixed normal ordering. No coupling-dependent
vacuum subtraction is made. The material momenta, including the
shift of $\Pi_Z$, remain in the operator. Fourier resolution acts
on the coupling, not on the carrier's charges or occupation numbers.
It is spatially nonlocal at finite $K$. Removing this cutoff while
controlling propagation, renormalized stress and interacting thermal
preparations is a separate limit.

The local current also fixes which coupling is intended. With signature
$(-,+)$ and $\epsilon^{01}=1$, let
$j^\mu=\epsilon^{\mu\nu}\partial_\nu\phi/\sqrt\pi$ and
$a_\mu=\lambda\theta f(Z)\partial_\mu Z$. Then
\[
 -a_\mu\partial^\mu\phi=A_\mu j^\mu,
 \qquad A_\mu=\sqrt\pi\,\epsilon_{\mu\nu}a^\nu.
\]
For a static texture, $A_0=-\sqrt\pi\lambda\theta f(Z)Z'$ and
$A_1=0$, giving the required density potential. The coframe's
$-a_\mu a^\mu/2$ term and momentum shifts accompany it.
The local dictionary does not by itself fix global holonomies or
spin weights~\cite[Eq.~(9) and Section~III]{YaoFukusumi2019};
the construction above uses the specified operator correspondence.

Rotating $(\phi_1,\phi_2)$ to active and spectator Gaussian modes
must also rotate their charge lattice. For instance $(w_1,w_2)=(1,0)$
becomes $(1/\sqrt2,1/\sqrt2)$, not two independent integer windings.
This scalar normal-mode change is different from rotating the actual
fermion flavors within each direction. A Gaussian diagonalization
cannot discard their spin constraint or their charged operators.

This material also admits finite autonomous record-controlled operation.
Extend its Hamiltonian $H_K$ trivially over finite degenerate record
registers. For a bounded self-adjoint permitted operator $C$ and
real $h\in C_c^\infty$, set
\[
 A=\int h(s)e^{isH_K}Ce^{-isH_K}ds,\qquad V=e^{i\ell A}.
\]
Integration by parts gives
$\|[H_K,A]\|\le\|h'\|_1\|C\|$ and
$\|[H_K,[H_K,A]]\|\le\|h''\|_1\|C\|$.
The energy-change estimates preceding
Proposition~\ref{prop:exterior-source-gate} now apply with the
\emph{interacting material} $H_K$, so finite products meet
Proposition~\ref{prop:energy-regular-return}.
Record-controlled $C$ are finite block sums of the permitted operators;
all previous reference systems remain in its input purification.

For $0\le g\in C^\infty(S^1)$, weighting the regulated coupled-sector
energy density above by $g$ gives a positive form $Q_g$ with
$0\le Q_g\le\|g\|_\infty H_K$. It is a closable sum of the same
weighted differential and multiplication squares. Its actual
expectation during autonomous execution obeys
Eq.~\eqref{eq:energy-regular-form}, with
$a=\|g\|_\infty$ and $d=0$.
This supplies a state and energy certificate on the charged material
itself, without a hard physical energy cutoff. It concerns the
regulated energy test: a renormalized covariant stress, spatially
local history clock and interacting thermal preparation need their
own estimates. Gate permission and the finite-rms initial state are
specified inputs, not products of this smoothing argument.

\subsection{A finite material mode and its remainder}
\label{sec:finite-material-mode}

The finite mechanical coordinates need not arise from fields uniform
over the whole line. A local binding potential supplies normalizable
angle and storage modes, with the discarded modes still present.
Here we specify a fixed-texture Gaussian realization; its state,
source and thermal estimates are not an equivalence theorem for the
preceding nonlinear charged regulator.
In normalized fields $u=(\sqrt I\theta,\sqrt J y)$, let
\[
 A_{\rm m}=\begin{pmatrix}k/I&-k/\sqrt{IJ}\\
              -k/\sqrt{IJ}&(k+K)/J\end{pmatrix},
 \qquad \alpha=\lambda_{\min}(A_{\rm m})>0.
\]
Choose $0<s<1$ and $\mu>0$ with $\alpha>s\mu^2$.
The material operator
\[
 L_{\rm m}=-\partial_x^2+A_{\rm m}+s^2\mu^2I
              -s(s+1)\mu^2\operatorname{sech}^2(\mu x)I
\]
has a strictly positive pointwise mass matrix.
For a supplied kink $Z_0=z_0\tanh(\mu x)$, the well is
$1-Z_0^2/z_0^2$. Using that polynomial in the material potential
keeps it nonnegative for all real $Z$, since $1-Z^2/z_0^2\le1$.
The well can therefore belong to the same positive local action;
the kink's preparation and position fluctuations remain physical inputs.

\begin{proposition}[Finite modes and their reaction bound]
\label{prop:finite-material-mode}
The normalized $b(x)=N\operatorname{sech}^s(\mu x)$ gives two
material modes with squared-frequency matrix $A_{\rm m}$.
The orthogonal material space has gap
$m_{\rm gap}^2=\alpha+s^2\mu^2$.
For compact real $a(x)$, the complete Gaussian form is
\[
 q[\phi,u]=\|\phi'+a u_\Theta\|^2+\langle u,L_{\rm m}u\rangle.
\]
Retaining the carrier and $u=bq_{\rm m}$ gives
\[
 q_P[\phi,q_{\rm m}]=\|\phi'+ab q_\Theta\|^2
                              +q_{\rm m}^T A_{\rm m}q_{\rm m}.
\]
For $z<m_{\rm gap}^2$, elimination of the orthogonal material
space has self-energy $\Sigma(z)\ge0$, with
\[
 \langle\chi,\Sigma(z)\chi\rangle
 \le\frac{\|Q(a\phi'+a^2bq_\Theta,0)\|^2}{m_{\rm gap}^2-z}.
\]
Here $Q$ projects off $b$ in both material components, and the
exact retained operator is $K_{PP}-\Sigma(z)$.
\end{proposition}
\begin{proof}
For $D_s=\partial_x+s\mu\tanh(\mu x)$,
$L_{\rm m}=D_s^*D_s+A_{\rm m}$ and $D_sb=0$.
Its partner is
\[
 D_sD_s^*=-\partial_x^2+s^2\mu^2+
                       s(1-s)\mu^2\operatorname{sech}^2(\mu x)
                       \ge s^2\mu^2.
\]
The nonzero spectra of the two products agree; thus $b$ is the only
spatial bound profile, and its orthogonal space has the stated gap.
The two retained frequencies need not both lie below that gap.
Normalization gives
$N^2=\mu\Gamma(s+1/2)/(\sqrt\pi\Gamma(s))$ and
$\|b'\|^2=s^2\mu^2/(2s+1)$, so its gradient cost is finite and
retained in $A_{\rm m}$. Projection gives $q_P$ exactly.
The omitted block is
$K_{QQ}=Q[L_{\rm m}+\operatorname{diag}(a^2,0)]Q
\ge m_{\rm gap}^2Q$, and
$K_{QP}\chi=Q(a\phi'+a^2bq_\Theta,0)$.
Solving its equation gives
$\Sigma(z)=K_{PQ}(K_{QQ}-z)^{-1}K_{QP}$ and the bound.
\end{proof}

The projected contact has profile $ab$ and its complete square.
The original $k,K,I,J$ are recovered, rather than fitted after the
projection. The remaining material field has not been reset or removed.
For a response bound put
\[
 d_{\rm eff}=\|ab\|^2,\qquad
 \beta_a=4\|a\|_\infty^2(1+d_{\rm eff}/\alpha)
                       +2\|a^2b\|^2/\alpha.
\]
The inequality
$\|\phi'\|^2\le2(1+d_{\rm eff}/\alpha)q_P$ yields
$\Sigma(z)\le\beta_a K_{PP}/(m_{\rm gap}^2-z)$ as forms.
For $\sigma>0$, let
$\epsilon=\beta_a/(m_{\rm gap}^2+\sigma^2)<1$,
$R_{\rm ex}=P(K+\sigma^2)^{-1}P$ on the retained space and
$R_P=(K_{PP}+\sigma^2)^{-1}$. Normalizing the Schur complement by
$K_{PP}+\sigma^2$ gives $I-D$, $0\le D\le\epsilon I$, hence
\[
 \|R_{\rm ex}-R_P\|
             \le\frac{\epsilon}{(1-\epsilon)\sigma^2}.
\]
This certifies the exponentially weighted causal response:
$\int_0^\infty e^{-\sigma t}\sin(t\sqrt K)/\sqrt K\,dt
=(K+\sigma^2)^{-1}$.

For finite-time linear-response forcing $F(t)=K_{QP}\chi(t)$ and zero
omitted response displacement and velocity, its exact reaction is
\[
 u_Q(t)=-\int_0^t K_{QQ}^{-1/2}
       \sin((t-r)\sqrt{K_{QQ}})F(r)\,dr,
 \qquad \sqrt{2E_Q(t)}\le\int_0^t\|F(r)\|dr.
\]
Here $E_Q$ is the energy of that response, not its vacuum fluctuations.
The force uses the actual retained history, including backreaction.
An initial omitted excitation adds its homogeneous evolution.
Near the material threshold, or near a retained resonance, a small
projection alone does not bound unfiltered scattering or rapid gates.
The full local Gaussian propagator retains its light cone; a single
collective coordinate generally does not preserve that sharp front.
The estimates concern linear response and omitted energy, not a
complete quantum-state or interacting thermal equivalence. They give
a finite-body reduction with an explicit remainder while the physical
clock, nonlinear recoil and thermal renewal are still to be matched.

\subsection{A local record drive on the same quantum material}
\label{sec:local-material-actuation}

A finite record can drive the material without replacing its vacuum by
that of the projected oscillator. Keep the complete Gaussian operator
$K$ of Section~\ref{sec:finite-material-mode}, let $\omega=\sqrt K$
and $H=d\Gamma(\omega)$, and use one common material vacuum calibration.
For a real material smearing $h=(0,h_\Theta,h_Y)$ define
\[
 Q(h)=\frac{a(\omega^{-1/2}h)+a^*(\omega^{-1/2}h)}{\sqrt2}.
\]
The bound $q\ge\alpha\|u\|^2$ gives
$\langle h,K^{-1}h\rangle\le\|h\|^2/\alpha$ by the inverse-form
variational formula. Spectral Cauchy--Schwarz then gives
$\|K^{-1/4}h\|^2\le\|h\|^2/\sqrt\alpha$.
Thus these material smearings are defined even when the carrier has no
mass gap. No static infrared dressing unitary is assumed.

Let $f\in C_c^\infty(0,T)$ be real and $F_1=\int|f|$.
An acquired classical bit $r\in\{0,1\}$ selects the drive
$-rf(t)Q(h)$. In the interaction picture its exact unitary is
\begin{align}
 U_{I,r}&=e^{ir^2\Phi_h}W(rz_h),&
 z_h&=\frac{i}{\sqrt2}\omega^{-1/2}
                 \int_0^T f(t)e^{it\omega}h\,dt,\nonumber\\
 \Phi_h&=\frac12\int_{s>t}f(s)f(t)
   \langle h,\omega^{-1}\sin((s-t)\omega)h\rangle\,dt\,ds.
 \label{eq:material-drive}
\end{align}
The commutator is a scalar, so the Magnus expansion stops here.
The Weyl convention and coherent-state overlap are standard~\cite{Weedbrook}.
In the Schr\"odinger picture the mean displacement is the full retarded
response
\[
 \xi_h(t)=r\int_0^t\frac{\sin((t-s)\omega)}{\omega}h f(s)\,ds.
\]
For centred Gaussian input, its added energy is
\[
 E_h=\|\omega^{1/2}rz_h\|^2
   =\frac{r^2}{2}\left\|\int f(t)e^{it\omega}h\,dt\right\|^2
   \le\frac{r^2F_1^2}{2}\|h\|^2.
\]
It equals the actual force work $\int rf\langle h,\dot\xi_h\rangle dt$.
The instantaneous forced form is bounded below by
$-r^2f(t)^2\langle h,K^{-1}h\rangle/2$.
Adding the opposite scalar makes it nonnegative; that offset is not
net work, heat or a construction of the driver.

To address the angle mode, choose smooth $0\le\chi_R\le1$, equal to one
on $|x|\le R$ and zero outside a finite larger interval. Set
\[
 A_R=\int b^2\chi_R,\quad h_R=(0,b\chi_R/A_R,0),\quad
 h_b=(0,b,0),\quad d_R=\|h_R-h_b\|.
\]
Both have retained force $\langle b,h_\Theta\rangle=1$.
Only $h_R$ is the permitted compact drive; $h_b$ is a comparison
smearing in the same full model. If $t_R=\int_{|x|>R}b^2$, then
$d_R^2\le t_R/(1-t_R)$ and $\|h_R\|^2\le(1-t_R)^{-1}$.
For $s=1/2$, $t_R=1-(2/\pi)\arctan(\sinh(\mu R))$.
Reducing this tail grows the actual drive support.

\begin{proposition}[Local actuation with the initial state retained]
\label{prop:local-material-actuation}
Let $\Delta z=z_{h_R}-z_{h_b}$, with both drives evolved under full $K$.
Then
\[
 \|\Delta z\|^2\le\frac{F_1^2d_R^2}{2\sqrt\alpha},\qquad
 \Delta E:=\|\omega^{1/2}\Delta z\|^2
                         \le\frac{F_1^2d_R^2}{2}.
\]
For an actual initial field marginal $\sigma_r$ with its reference
retained, the half trace distance of the two complete outputs obeys
\[
 D_r\le\sqrt{1-|\operatorname{Tr}\sigma_r W(r\Delta z)|^2}.
\]
A classical distribution of records gives joint distance at most
$\sum_r p_rD_r$. On the same Gaussian ground state this becomes
$\sqrt{1-e^{-r^2\|\Delta z\|^2}}$.
On a finite Dirichlet interval, for the actual full-$K$ Gibbs state
$\sigma_\beta=e^{-\beta H}/Z$, the two displaced field states satisfy
\[
 S(\sigma_{h_R}\|\sigma_{h_b})=\beta r^2\Delta E,
 \qquad D_{{\rm field},r}\le\sqrt{\beta r^2\Delta E/2}.
\]
\end{proposition}
\begin{proof}
The bound on $K^{-1/4}$ and the spectral integral in
Eq.~\eqref{eq:material-drive} prove the first two inequalities.
Writing $B_R=\int b^2\chi_R^2\le A_R$ gives
$d_R^2=B_R/A_R^2-1\le A_R^{-1}-1\le t_R/(1-t_R)$.
Purify the actual field and reference state. The overlap of the two
outputs is its Weyl characteristic function, up to a scalar phase.
The pure-state trace-distance formula and partial-trace contraction
give the retained-reference bound. For the Gibbs identity, conjugate
both outputs by $W(-rz_{h_b})$. Displacement preserves entropy,
$\log\sigma_\beta=-\beta H-\log Z$, and the initial mean is zero.
Its energy increment is $r^2\Delta E$; Pinsker gives the field-only bound.
\end{proof}

The Gibbs energy bound concerns the reduced field states. A retained
thermal purification still uses the characteristic-function bound;
its covariance can make that distance larger.
Here $\Delta E$ is the energy of the difference response, not
$E_{h_R}-E_{h_b}$, which has a cross term. No independent ground state
of $K_{PP}$ is substituted. Its dynamic reduction still has the
reaction error of Proposition~\ref{prop:finite-material-mode}.
An actual old field state must use its own characteristic function;
initial correlation is not replaced by a fresh thermal packet.
Coherent record controls also retain the Magnus and Weyl phases and
are not covered by the classical-record distance bound.

For a centred Gaussian input the covariance is unchanged. A quadratic
source, renormalized with the same fixed vacuum subtraction in both
experiments, therefore changes by the classical value $T[\xi_h]$.
Nonzero initial mean adds its bilinear cross term.
The local energy balance has force input
$rf(t)h\cdot\dot\xi_u$; the driver must supply that work.
It uses the complete response, including omitted material modes,
not merely the isolated mode's energy. This is the fixed-texture
Gaussian source, not the full moving quantum apparatus's stress.
The emitted carrier is determined by that material motion. Its equation
is $(\partial_t^2-\partial_x^2)\xi_\phi=\partial_x(a\xi_\Theta)$.
For zero past displacement, outside the compact contact the right and
left outgoing profiles are
\[
 \xi_R(u)=-\tfrac12\int a(y)\xi_\Theta(y,u+y)dy,
 \qquad
 \xi_L(v)=\tfrac12\int a(y)\xi_\Theta(y,v-y)dy,
 \quad u=t-x,\ v=t+x.
\]
Integrating the one-dimensional retarded Green function against
$\partial_y(a\xi_\Theta)$ gives these endpoint integrals directly.
The extra stresses of centred Gaussian input are
$(\xi_R')^2$ and $(\xi_L')^2$. They use the complete material response,
not an independent emitter. With no incoming drive-induced radiation,
the total emitted carrier energy is at most the force work $E_h$;
the remaining material energy is retained. A zero contact $a=0$ emits
no carrier source despite possible stored material energy. Opposite
drives also have equal quadratic stress on centred input, so a command
label alone does not identify the emitted source.

A finite menu of these record-selected unitaries admits the bounded
prepared-bit dilation of Appendix~\ref{sec:local-thermal-preparation}.
For compact smooth $h_R$ on the Gaussian ground state, its Weyl
amplitudes have all positive frequency moments, so the required twice
differentiable $H$ graph orbits exist. Smoothing and autonomous return
then use this same $H$. Other initial states need their own graph-moment
certificate. The extra bit, clock, battery, program and retained
outputs remain resources. The record must reach the entire drive
support by physical transport before the operation; the dilation does
not prove that transport or the clock's covariant stress. Approximating
a smooth local source during autonomous execution also needs its
relative source bound, not trace distance alone. The following estimate
provides it for finite-interval output readings of this material.

\subsection{A source certificate for the quantum material}
\label{sec:Gaussian-source-certificate}

Fix a Dirichlet interval of length $L$ and keep all modes of the
complete Gaussian $K$. Write its pointwise mass matrix as
$mI\le M(x)\le M_{\max}I$, $m>0$. For
$X=(\phi,\Theta,Y)$ its spatial energy form is
\[
 q[X]=\|\phi'+a\Theta\|^2+\|u'\|^2+\langle u,Mu\rangle,
 \qquad u=(\Theta,Y).
\]
With $C_g=2+2\|a\|_\infty^2/m$ and
$V=2\|a\|_\infty^2+M_{\max}$, Cauchy--Schwarz gives
$\|X'\|^2\le C_gq[X]$ and
$q[X]\le2\|X'\|^2+V\|X\|^2$.
Min--max on three Dirichlet components bounds the ordered frequencies:
\[
 c_-j\le\omega_j\le c_+j,\qquad
 c_-=\frac{\pi}{3L\sqrt{C_g}},\quad
 c_+=\sqrt{2\pi^2/L^2+V}.
\]
The full tower, not a finite Fock cutoff, enters the next statement.

Let $\mathcal Y=(\Pi,\phi'+a\Theta,u',\sqrt M u)$ be the energy-factor
fields, so their squared norm divided by two is the energy density.
For a real symmetric bounded spatial matrix $B(x)$,
$\|B(x)\|\le G$, and real $h\in C_c^\infty$, set
\[
 O=\frac12\int h(t)\int
       :\mathcal Y(t,x)^TB(x)\mathcal Y(t,x):\,dx\,dt.
\]
Normal ordering uses the same $K$ vacuum in both compared states.
$B=gI$ tests the local energy. In a free exterior carrier region a
rank-one choice tests $g(\Pi_\phi\mp\phi')^2/4$, the corresponding
null stress. These tests may have negative quantum expectations.

\begin{proposition}[A same-Hamiltonian quadratic-source bound]
\label{prop:Gaussian-source-certificate}
Put $H=d\Gamma(\omega)$, $H_N=\|h^{(N)}\|_1$, $N\ge2$.
The source is a Hermitian form on $\operatorname{Dom}H^{1/2}$,
\[
 O=d\Gamma(A)+P+P^*,\qquad P=\sum_{ij}b_{ij}a_i a_j.
\]
It obeys $|d\Gamma(A)|\le a_0H$, $a_0=G\|h\|_1$.
Define
\[
 \kappa^2=\sum_{ij}\frac{|b_{ij}|^2}{\omega_i\omega_j},\qquad
 \nu^2=\sum_{ij}\frac{|b_{ij}|^2}{\omega_i}.
\]
Both converge, with the explicit sufficient bounds
\begin{align}
 \kappa&\le \frac{GH_N}{2c_-^N}
       \sqrt{\zeta(2N-1)-\zeta(2N)},\nonumber\\
 \nu^2&\le \frac{G^2H_N^2c_+}{8c_-^{2N}}
       [\zeta(2N-2)-\zeta(2N-1)].
 \label{eq:Gaussian-pair-bounds}
\end{align}
For any reference energy $\Lambda>0$, let
$C_O=a_0+2\kappa+\nu/\sqrt\Lambda$.
Then $(H+\Lambda)^{-1/2}O(H+\Lambda)^{-1/2}$ is a bounded
Hermitian form of norm at most $C_O$.
\end{proposition}
\begin{proof}
Choose real eigenmodes $e_j$ of $K$. If $L_s^*L_s=K$ is its
spatial energy factor, the annihilation coefficients of $\mathcal Y$
are
\[
 v_j=(-i\sqrt{\omega_j/2}\,e_j,
                       L_se_j/\sqrt{2\omega_j}),\qquad
 \langle v_i,v_j\rangle=\omega_j\delta_{ij}.
\]
Spatial multiplication by $B$ is bounded by $G$ on this factor
space. Time evolution preserves the energy norm; integration against
$h$ proves the number-form estimate.
The pair coefficient is
$b_{ij}=\tfrac12\int v_i^TBv_jdx\,\widehat h(-\omega_i-\omega_j)$,
where $\widehat h(\xi)=\int h(t)e^{i\xi t}dt$.
Integrating $N$ times by parts gives
\[
 |b_{ij}|\le \frac{G\sqrt{\omega_i\omega_j}H_N}
                         {2(\omega_i+\omega_j)^N}.
\]
Summing over $i+j=k$ uses $k-1$ pairs and
$\sum_{j=1}^{k-1}j=k(k-1)/2$; the frequency bounds give
Eq.~\eqref{eq:Gaussian-pair-bounds}.

For a unit finite-particle spectral vector of energy $E$,
Cauchy--Schwarz applied to $\sqrt{\omega_i}a_i\psi$ and
$\sum_j\overline{b_{ij}}a_j^*\psi/\sqrt{\omega_i}$ gives
\[
 |\langle P\rangle|\le\sqrt{E(\kappa^2E+\nu^2)}
                         \le\kappa E+\nu\sqrt E.
\]
The creation-vector square is a vacuum term $\nu^2$ plus a number
form bounded by $\kappa^2H$: the matrix
$\omega^{-1/2}b\omega^{-1/2}$ has Hilbert--Schmidt norm $\kappa$.
Thus $|\langle O\rangle|\le(a_0+2\kappa)E+2\nu\sqrt E
\le C_O(E+\Lambda)$.
Pair tails converge in these two sums; the uniformly bounded number
forms converge on the finite-energy core. Their normalized forms
therefore extend to the stated bounded Hermitian form.
\end{proof}

This estimate concerns a spacetime-smeared source form, not an
unsmeared density. Its constants depend on $L$; they do not prove a
uniform real-line limit. Compact $W^{2,1}$ time tests also satisfy
the $N=2$ proof. No free chiral energy inequality has been assigned
to the coupled material.

\begin{corollary}[Output source error with references retained]
\label{cor:Gaussian-source-error}
For normalized actual and target vectors $\psi,\varphi\in
\operatorname{Dom}H$, including their actual references, let
$e=\|\psi-\varphi\|$ and $g_H=\|H(\psi-\varphi)\|$.
Then
\begin{align}
 |\langle O\rangle_\psi-\langle O\rangle_\varphi|
 &\le C_O\sqrt{\Lambda e^2+eg_H}\nonumber\\[-2pt]
 &\quad\times\left[\sqrt{\Lambda+\|H\psi\|}
                    +\sqrt{\Lambda+\|H\varphi\|}\right].
 \label{eq:Gaussian-source-error}
\end{align}
\end{corollary}
\begin{proof}
Tensor the bounded normalized form with the reference identity and
polarize. Its difference is bounded by $C_O$ times
$\|(H+\Lambda)^{1/2}(\psi-\varphi)\|$ times the sum of the two
weighted norms. The first squared norm is at most
$\Lambda e^2+eg_H$; each remaining squared norm is at most
$\Lambda+\|H\psi\|$ or $\Lambda+\|H\varphi\|$.
\end{proof}

The energy-regular return theorem supplies $e,g_H$ for a specified
same-$H$ graph-domain operation. It consequently certifies these
output source readings under common subsequent $H$ evolution, not
the missing driver or clock stress during its interaction.
In fact its uniform physical rms bounds $M_a,M_i$ suffice:
$g_H\le M_a+M_i$ in Eq.~\eqref{eq:Gaussian-source-error}.
The resulting source error tends to zero with the autonomous state
error even as the history clock grows; its physical rms bounds do
not grow with that clock. This does not add a permanent halt or
decouple an always-running controller for free.
A mixed input uses its retained purification. On this finite
interval a full-$K$ Gibbs purification has finite energy moments:
with $q_\beta=e^{-\beta c_-}$,
\[
 \langle H\rangle\le \frac{c_+q_\beta}{(1-q_\beta)^3},\qquad
 \operatorname{Var}H\le \frac{c_+^2q_\beta(1+q_\beta)}
                                    {(1-q_\beta)^5}.
\]
The bosonic product partition function converges. Bounding occupations
by $q_\beta^j/(1-q_\beta)$ and summing $j q_\beta^j$ and
$j^2q_\beta^j$ proves the two estimates. All higher polynomial moments
are finite as well. This does not prepare the Gibbs state, prove a
particular gate's graph orbit, or renew a correlated reservoir.

For the vacuum-driven outputs $|z\rangle,|w\rangle$ of two actual
forces on this same material, put $c=e^{-\|z-w\|^2/2}$,
$E_z=\|\omega^{1/2}z\|^2$ and
$\Delta E=\|\omega^{1/2}(z-w)\|^2$.
Choosing only a comparison phase to make their overlap $c$ gives
the exact identity
\[
 \|(H+\Lambda)^{1/2}(|z\rangle-|w\rangle)\|^2
       =c\Delta E+(1-c)(E_z+E_w+2\Lambda).
\]
Indeed $\langle z|H|w\rangle=\langle z|w\rangle
\langle z,\omega w\rangle$.
Their source difference is at most $C_O$ times the square root of
this expression times
$\sqrt{E_z+\Lambda}+\sqrt{E_w+\Lambda}$.
Proposition~\ref{prop:local-material-actuation} bounds its two
difference terms for compact forces. This comparison keeps full $K$;
changing the box or using the projected oscillator changes the
problem. For coherently controlled records the actual relative
phases remain, rather than being physically aligned branch by branch.

Trace distance alone cannot replace these estimates: a trace-small
admixture can have arbitrarily large energy. Nor does matching mean
energy suffice. A small vacuum--two-particle amplitude changes a
pair-source reading to first order while its energy changes to second
order. The source certificate retains both the state and its energy
norm. Matching this finite-interval Gaussian readout to a charged
nonlinear continuum and the complete moving apparatus remains a
separate physical construction.

\section{Finite return through one thermal field}
\label{sec:thermal-return}

The oscillator loop closes at one frequency. A propagating field has an
interval of frequencies, and a sequence timed for its carrier does not
close at every detuning. We now compute what remains, and how a single
thermal line can be used over a finite adaptive protocol without resetting
it between operations. The field theory, phase controls and reading
instruments are still inputs.

\subsection{Finite-width forces and the outgoing record}

Let $\omega=\Omega+\delta$, $|\delta|\le B<\Omega$, and let $g$ be real,
even, square-integrable, and nonzero on an interval. In units $\hbar=1$,
the field Hamiltonian is
$H_F=\int(\Omega+\delta)a_\delta^\dagger a_\delta\,d\delta$.
For a system projector $P$ with both eigenvalues $0,1$, apply
\begin{equation}\label{eq:line-source}
 H_s^I(t)=iP\int g(\delta)
 \left[u(t)e^{i\delta t}a_\delta^\dagger
       -\overline{u(t)}e^{-i\delta t}a_\delta\right]d\delta .
\end{equation}
This is the interaction picture of a prescribed carrier-modulated
force. It supplies two field quadratures and a clock; it is band-smeared,
not compactly localized in space. At finite $u(t)$ the laboratory
Hamiltonian is bounded below by
$-|u(t)|^2\int g^2/(\Omega+\delta)\,d\delta$.
A scalar shift makes it nonnegative without altering the gate.

Use four disjoint rectangular pulses of width $\sigma<\tau$, centred
at $j\tau$, with complex labels $z_j=(1,i,-1,-i)$ and heights
$r z_j/(\sigma\sqrt Z)$, where
\[
 Z=\int g(\delta)^2\operatorname{sinc}^2(\delta\sigma/2)d\delta,\qquad
 f(\delta)=g(\delta)\operatorname{sinc}(\delta\sigma/2)/\sqrt Z,
 \quad \operatorname{sinc}x=\sin x/x .
\]
Thus $\|f\|_2=1$. The individual field displacements are
$\alpha_j=r z_j f e^{ij\delta\tau}$.
Put
\[
 M(t)=\int|f|^2\cos(\delta t)d\delta,\qquad
 C=3M(\tau)-M(3\tau),\qquad
 F=r f(1-e^{2i\delta\tau})(1+ie^{i\delta\tau}).
\]

\begin{proposition}[Continuum loop and its thermal defect]
\label{prop:thermal-loop}
The complete four-pulse unitary is
\begin{equation}\label{eq:thermal-loop}
 U_{\rm loop}=e^{ir^2CP}D(PF),\qquad
 D(PF)=\exp\!\int P(Fa^\dagger-\overline F a)d\delta .
\end{equation}
If $C>0$, choose $r^2=\theta/C$ for a target phase
$e^{i\theta P}$, $\theta>0$. For an initially independent thermal
field with inverse temperature $\beta$, define
\begin{equation}\label{eq:thermal-D}
 D_\beta=\int |F(\delta)|^2
       \coth\!\left[\beta(\Omega+\delta)/2\right]d\delta .
\end{equation}
The thermal state is specified by its Weyl characteristic; a continuum
KMS state need not be a trace-class density matrix on vacuum Fock space.
After removal of the target phase the system channel is dephasing
between the two $P$ blocks, with coherence factor $e^{-D_\beta/2}$.
Its half diamond distance from the target unitary is
$(1-e^{-D_\beta/2})/2$, and its representation defect norm on a
unit-norm off-block Pauli control is $1-e^{-D_\beta}$.

No nonzero compact-time force through this single profile can close the
continuum exactly. This last statement concerns the source
\eqref{eq:line-source}, not all multimode control architectures.
\end{proposition}
\begin{proof}
The source commutator is a scalar field commutator times $P^2=P$,
so the time-ordered exponential is a conditional displacement and a
commuting $P$ phase. Within each constant-phase pulse the self-phase
vanishes after integration over the even band: its sine kernel is odd
in $\delta$. The six inter-pulse Weyl phases add to
$r^2[3M(\tau)-M(3\tau)]$. Summing the four displacement vectors gives
$F$. This derivation includes finite pulse width; the pulses are not
delta kicks.

The thermal Weyl characteristic is
$\langle D(F)\rangle_\beta=e^{-D_\beta/2}$.
Tracing the field gives the stated channel, a mixture of identity and
conjugation by $2P-I$ after the target phase. Orthogonal off-block
superpositions attain the half diamond distance; multiplying its
Heisenberg off-block factor gives the defect.

For a general integrable compact-time $u$ the displacement is
$g(\delta)\int u(t)e^{i\delta t}dt$. Exact unitary control after discarding
the thermal field requires this displacement to vanish. The Fourier
transform of compact-time $u$ is entire. Vanishing on an interval where
$g\ne0$ forces it to vanish identically, and Fourier uniqueness gives
$u=0$. At one oscillator frequency only a single zero is required,
which is why Proposition~\ref{prop:loop} has no such obstruction.
\end{proof}

For $\theta=\pi$ and $P=(I+X)/2$, the target is a corrective flip up
to phase. With $\mu_2=\int\delta^2|f|^2d\delta$,
\begin{align}
 C&\ge2-\tfrac32\tau^2\mu_2, \label{eq:line-Cbound}\\
 \|F\|_2^2&\le8r^2\tau^2(1+B\tau)\mu_2, \label{eq:line-Fbound}\\
 D_\beta&\le
 \coth[\beta(\Omega-B)/2]\|F\|_2^2. \label{eq:line-Dbound}
\end{align}
The first uses $M(\tau)\ge1-\tau^2\mu_2/2$ and
$M(3\tau)\le1$. For the second, $|1-e^{2i\delta\tau}|\le2|\delta|\tau$
and $|1+ie^{i\delta\tau}|^2=2-2\sin(\delta\tau)\le2(1+B\tau)$.
Temperature multiplies the residual record; it does not remove it.

\subsection{One line throughout an adaptive sequence}

Multiplying the reduced one-loop channels would assume a fresh field.
That is not the evolution we consider. The same field and its correlations
remain present, and subsequent projectors need not commute.

\begin{proposition}[Finite reuse without independence]
\label{prop:thermal-reuse}
Start with one thermal line independent of system, controller and retained
reference. A protocol uses at most $L$ loops, interspersed with arbitrary
system--controller instruments whose dilations do not act on the line.
Projectors may be selected from acquired records and may differ between
loops. Suppose every permitted branch has $D_\beta\le D_{\max}$
and its target phase is implemented as in \eqref{eq:thermal-loop}.
The complete actual process, including final records, system and field,
is at operational state distance at most
\[
 \min\{1,L\sqrt{D_{\max}}\}
\]
from the process with ideal target gates and the same initial line.
Every terminal loss in $[0,1]$ has an expectation difference within
this bound, as does the total variation of the complete record laws.
No reset or independent-error hypothesis is used.
Here operational distance is half the norm distance of the states on
their joint observable algebra; it is half trace distance whenever a
density-matrix description applies.
\end{proposition}
\begin{proof}
Purify the initial thermal state, or use its thermal GNS cyclic vector.
For a loop acting on the ideal-prefix state,
\[
 \|(D(PF)-I)|\Psi\rangle\|^2
 =2\langle P\rangle(1-e^{-D_\beta/2})\le D_\beta .
\]
This is uniform over system inputs, retained entanglement and the
record-controlled choice of projector. In the ideal prefix the line
remains thermal and independent; the actual prefix need not.
Telescope the actual unitary dilations against the ideal dilations,
placing the ideal prefix on each difference and leaving the actual
suffix as a norm-preserving factor. The vector difference is at most
$L\sqrt{D_{\max}}$. The norm distance of vector states, followed by
restriction to physical observables or processing records, gives the
stated operational bound.
The proof does not replace any actual correlated prefix by a reset state.
\end{proof}

For constant per-pulse relative amplitude error at most $e_a$, centre
error $d_t$ and carrier-phase error $e_\varphi$, with unchanged pulse
widths and still disjoint pulses, each displacement changes in vacuum
norm by at most
\[
 e=r[e_a+(1+e_a)(Bd_t+e_\varphi)].
\]
The six Weyl phases change by at most $6(2re+e^2)$.
The same proof therefore gives the finite-setting bound
\begin{equation}\label{eq:line-setting}
 L\left[\sqrt{D_{\max}}
 +4\sqrt{\coth[\beta(\Omega-B)/2]}\,e
 +6(2re+e^2)\right].
\end{equation}
It concerns these specified setting errors, not arbitrary time-dependent
source noise or heating added during the force.

As a concrete family take
$g^2=15[1-(\delta/B)^2]^2/(16B)$, $\Omega=10$, $B=1$,
$\beta=.2$ and $\sigma=\tau/2$. The unfiltered second moment is
$B^2/7$; since $\operatorname{sinc}^2(\delta\sigma/2)\le1$,
$\mu_2\le B^2/(7Z)$ and
$Z\ge\operatorname{sinc}^2(B\sigma/2)$ in the stated small-width range.
At $\tau=5\times10^{-5}$, $e_a=e_\varphi=10^{-8}$ and
$d_t=10^{-9}$, conservative rational bounds in
\eqref{eq:line-setting} put an entire 64-loop adaptive protocol below
$.00512$ operational distance. Tight numerical integration gives
one-loop error $1.48\times10^{-9}$. The sequence statement comes
from the joint bound, not from treating those tiny errors as independent.
The clock's frequency units are supplied by $B$; these numbers are not
unlabelled laboratory seconds.

\subsection{Source energy and control resources}

For one loop and an independent zero-mean thermal input, the excess
free-field energy is
\begin{equation}\label{eq:line-energy}
 \Delta E_F=\hbar\langle P\rangle
       \int(\Omega+\delta)|F(\delta)|^2d\delta .
\end{equation}
For the chiral current normalization of Paper~B, its corresponding
free outgoing excess stress is $4\langle P\rangle[\operatorname{Re}
\psi_F(u)]^2$, where
$\psi_F(u)=\int-i\sqrt{\omega/(4\pi)}F(\omega)e^{-i\omega u}d\omega$.
The chosen $g$ vanishes at band edges and gives finite first two time
moments. This is the free field source, not the stress of its support,
phase reference or switching apparatus.

A reused line has extra cross terms. Let
$E_1=\int\omega|F|^2d\omega$ bound the unit-projector displacement
energy of each loop, and let
$n_{\max}=[e^{\beta(\Omega-B)}-1]^{-1}$.
For loops sharing the stated profile, with arbitrary projector choices,
\begin{equation}\label{eq:line-multienergy}
 |\Delta E_F^{(L)}|
 \le\hbar\left[2L\sqrt{(\Omega+B)n_{\max}E_1}+L^2E_1\right].
\end{equation}
Indeed the final annihilation field is its initial value plus
$\sum_jF_jP_j^{\rm past}$, with $\|P_j^{\rm past}\|\le1$.
The quadratic term is at most $L^2E_1$ by the weighted $L^2$ triangle
inequality. Each initial thermal cross term is at most
$\sqrt{(\Omega+B)n_{\max}E_1}$ by Cauchy--Schwarz. Thus the single-loop
energy cannot simply be multiplied by $L$. This bound prices an
excess over the supplied thermal background, not that background itself.

Small final radiation does not mean small driving demands. With
$\|g\|_2=1$, the four rectangular sources have
\[
 \int|u(t)|^2dt=\frac{4r^2}{\sigma Z},\qquad
 |H_{\min}|\le\frac{\hbar r^2}{\sigma^2 Z(\Omega-B)}
\]
during a pulse. The first is a declared quadratic envelope charge, not
heat; the second is the magnitude of a lower energy bound, not consumed
work. Narrowing the loop reduces its escaped record but increases these
source and clock requirements. At the 64-loop settings the single-loop
envelope charge is about $2.51\times10^5$, although its unit-projector
radiated energy is about $4.49\times10^{-8}\hbar$.
Finite pulse heights are allowed; ideal switching of the external
apparatus is still a supplied control.

For fixed band, temperature and target phase, $D_\beta=O(\tau^2)$ as
$\tau\to0$. Taking $\sigma=\tau/2$ and choosing $\tau$ proportional
to $\varepsilon/L$ therefore suffices for joint error $\varepsilon$,
with setting errors also reduced to order $\varepsilon/L$.
Within this family the total quadratic envelope charge scales as
$L^2/\varepsilon$. This is a constructive resource scaling, not a
lower bound for every field-control architecture.

Finite reuse removes the independent-port premise; the initially
thermal line is never reset. Stored records, controller renewal,
exhaust and total source/geometric response still require a compatible
apparatus, not merely a small reduced-channel error.
\section{Repeated renewal of a complete ring body}
\label{sec:ring-renewal}

Use a metric-graph carrier with four complex chiral body fields on
rings of length $w$, each
coupled to its own thermal line over an arc $0<\ell<w$. All spatial
modes remain. Their charges are $(+1,-1,+1,-1)$; the first pair is the
receiver, the second its current purifier. Two colocated, copropagating
rails on each arc obey
\begin{equation}
 i(\partial_t+\partial_x)\Psi=h(x)W(t-x)\Psi,
 \quad h\in C_c^\infty(0,\ell),\quad\int h=1.
 \label{eq:ring-local-transfer}
\end{equation}
The exact ray transfer is $e^{-iW(u)}$, $u=t-x$. Take
$W(u)=\theta(u)\sigma_y$ and include a physical ring connection with
round-trip phase $e^{-i\Delta w}$. The overlap transit is included in
$w$; the voltages and their switching work are supplied controls.

Choose $\tau=Kw$, $K\in\mathbb N$, a real
$\lambda\in C_c^\infty(0,\tau)$, and
$\theta=\sqrt w\lambda$, $|\theta|<\pi/4$. Iteration gives
\begin{align}
 A_w(u)&=\prod_{j=0}^{K-1}\cos[\sqrt w\lambda(u+jw)],
                         &&0\le u<w,\\
 F_w(s)&=\frac{\sin[\sqrt w\lambda(s)]}{\sqrt w}e^{i\Delta s}
          \prod_{j\ge1:\,s+jw<\tau}\cos[\sqrt w\lambda(s+jw)].
 \label{eq:ring-envelopes}
\end{align}
The old-ring multiplier is $A_w$; $F_w$ is the captured zero-mode
envelope. Put $\eta^2=w^{-1}\int_0^w|A_w|^2$. Retaining the full
complementary outgoing field gives $\|F_w\|_2^2=1-\eta^2$ by CAR.
Endpoint flatness makes the last product smooth. With $\Omega=2\pi/w$,
the captured $n$th mode has carrier $\Delta+n\Omega$.

The reference sea fills $n<0$ and leaves $n\ge0$ empty. Sum over all
four species in
\begin{align}
 N_{\rm ex}&=\sum_{n\ge0}d_n^\dagger d_n+\sum_{n<0}d_nd_n^\dagger,\\
 P&=\Omega\left(\sum_{n\ge0}n d_n^\dagger d_n+
                      \sum_{n<0}|n|d_nd_n^\dagger\right),
 \qquad Q=N_{\rm ex}+P/\Omega .
\end{align}
Let $r(E)$ be the occupation of the normalized captured envelope with
carrier $E$ in the actual temperature-$T$ KMS line. Its real envelope
gives $r(-E)=1-r(E)$. For $A_w(w\xi)=\sum_k a_ke^{2\pi ik\xi}$ set
\begin{align}
 \alpha&=\sum_k\sqrt{1+|k|}\,|a_k|,\\
 V_A&=\tfrac12\sum_{k>0}k(k+1)|a_k|^2
       +\tfrac12\sum_{k<0}|k|(|k|+3)|a_k|^2,\\
 S_Q&=(1-\eta^2)\left[r(\Delta)+\sum_{n\ge1}(1+n)
                \{r(\Delta+n\Omega)+r(n\Omega-\Delta)\}\right],\\
 b&=2(\sqrt{V_A}+\sqrt{S_Q}).
 \label{eq:ring-constants}
\end{align}

\begin{proposition}[Uniform renewal with retained history]
\label{prop:ring-renewal}
Suppose $\langle Q\rangle<\infty$, $0<\Delta<\Omega$, and the next
incoming supports have their unconditional KMS marginals. Correlations
with the body and all earlier records are allowed. The actual reset
satisfies
\begin{equation}
 \sqrt{\langle Q\rangle_{\rm reset}}
       \le\alpha\sqrt{\langle Q\rangle_{\rm old}}+b.
 \label{eq:ring-contraction}
\end{equation}
Between resets allow the translation-invariant pumps below and supplied
logical noise on the four zero modes. For post-reset
$x_j=\sqrt{\langle Q\rangle_j}$, if $\sqrt2\alpha<1$, then
\begin{equation}
 x_{j+1}\le\sqrt2\alpha x_j+2\alpha+b,
 \quad\sup_jx_j\le\max\left\{x_0,
                   \frac{2\alpha+b}{1-\sqrt2\alpha}\right\}.
 \label{eq:ring-uniform}
\end{equation}
Against the actual outside marginal, including all old references,
\begin{equation}
 \tfrac12\|\rho_{{\rm body},R}-|{\rm sea}\rangle\langle{\rm sea}|
                             \otimes\rho_R\|_1\le2x_j.
 \label{eq:ring-decouple}
\end{equation}
At fixed cycle duration and $T$, one finite choice of pulse strength,
gap and ring length makes this error arbitrarily small uniformly in
the number of cycles, with four active outlet species.
\end{proposition}
\begin{proof}
Apply all weighted wrong-occupation annihilators to a purification of
the actual joint state. Weighted convolution, using
$1+|n|\le(1+|m|)(1+|n-m|)$, bounds their old excitation-annihilation
part by $\alpha\sqrt{\langle Q\rangle}$. CAR bounds the creation
covariance by the identity; its squared norm is at most $V_A$ per
species. The incoming vector has squared norm $S_Q$. Minkowski's
inequality on the four-species direct sum proves
\eqref{eq:ring-contraction}, without factoring the old state from the
line. Pumps and logical noise preserve $P$, and
$N_{\rm ex}\le P/\Omega+4$, so
$Q_{\rm after\ gate}\le2Q_{\rm before\ gate}+4$ in expectation.
This gives \eqref{eq:ring-uniform}. The nonsea projection is bounded
by $N_{\rm ex}\le Q$. Projection of a joint purification onto the
pure sea and trace-distance contraction give \eqref{eq:ring-decouple}.

The full sums are finite before imposing any numerical cutoff.
Set $\Gamma=\int\lambda^2$ and $g_w=\log\cos(\sqrt w\lambda)$,
extended by zero. For each $m\ge3$,
$G_m=\sup_{w\le w_0}w^{-1}\|g_w^{(m)}\|_1<\infty$ if
$\sqrt{w_0}\|\lambda\|_\infty<\pi/4$. Poisson summation bounds the
nonconstant Fourier coefficients of $\log A_w$ by
$G_mw^m/(2\pi|k|)^m$. Weighted Wiener exponentiation therefore gives
\begin{equation}
 \alpha\le e^{-\Gamma/2+b_m},\quad
 V_A\le2e^{-\Gamma}(e^{b_m}-1)^2,\quad
 b_m=\frac{2G_mw^m}{(2\pi)^m}[\zeta(m)+\zeta(m-1)].
\end{equation}
Write $F_w=e^{i\Delta s}f_w$. Every derivative of its logarithmic
product is a sum of $O(1/w)$ terms of size $O(w)$, with flat endpoints.
Thus $D_m=\sup_{w\le w_0}\|f_w^{(m)}\|_2^2<\infty$. Splitting the
Fourier integral at $E/2$ bounds the exact Fermi occupation by
\begin{equation}
 r(E)\le e^{-E/(2T)}+\frac{(2/E)^{2m}D_m}{1-\eta^2},\qquad E>0.
\end{equation}
For $\Omega\ge2\Delta$ and $z=e^{-\Omega/(4T)}$,
\begin{align}
 S_Q\le{}&(1-\eta^2)r(\Delta)
     +2(1-\eta^2)\frac{z(2-z)}{(1-z)^2}\notag\\
     &+2D_m(4/\Omega)^{2m}[\zeta(2m)+\zeta(2m-1)].
\end{align}
Choose $\Gamma$ large, then $\Delta$ large, then $w=\tau/K$ small.
All errors become small simultaneously. A finite initial reset puts
$x_0$ at the same tolerance. Earlier nonselective even operations on
disjoint line characteristic supports leave the next unconditional
KMS marginal unchanged; guards longer than the overlap transit provide
that separation.
\end{proof}

The mixed target is restored by a local neutral quartet pump
\begin{equation}
 V_4=\int_0^w(D_1^\dagger D_2^\dagger D_3^\dagger D_4^\dagger
                          +\mathrm{h.c.})\,dx .
\end{equation}
It commutes with $P$. On the sea, creation requires four nonnegative
mode indices summing to zero: only the four zero modes occur. The
reverse transition obeys the same argument. Their two-dimensional
subspace is exactly invariant, with matrix element $1/w$.

For the full operator, bosonize the periodic fields with
$q_i\in\mathbb Z+1/2$. The charge coordinate
$\phi_c=\sum_i\phi_i/2$ gives the weight-two vertex $e^{2i\phi_c}$.
At fixed transverse charges, $q_c=q_0+2\mathbb Z$ with
$q_0\in\tfrac12\mathbb Z$, a fixed parity sector of a charge fermion;
the other three bosons and the original Klein phases remain. The
weight-two vertex splits this stress into two positive $c=1/2$
Virasoro stresses with velocities $1\pm\kappa/\pi$~\cite{Naud1999}.
The four Ramond ground weights sum to $1/2$, hence
\begin{equation}
 \|V_4\psi\|\le\pi^{-1}\|(P+\Omega/2)\psi\|.
\end{equation}
Finite-dimensional $P$ eigenspaces are invariant; their direct sum
defines the self-adjoint pump. The pair pump
$B=\int(D_1^\dagger D_2^\dagger+\mathrm{h.c.})dx$ obeys
$|B|\le P/\Omega+1$: every active nonzero pair costs at least
$\Omega$ in $P$. For sea-normal-ordered number $N\ge-P/\Omega$,
\begin{equation}
 P+\Delta N+\kappa V_4+b_2B\ge
 \left(1-\frac\Delta\Omega-\frac{|\kappa|}\pi
                        -\frac{|b_2|}\Omega\right)P
                           -\frac{|\kappa|}{w}-|b_2|.
 \label{eq:ring-positive}
\end{equation}
Small enough pump peaks give a common finite physical ground offset
and positive Hamiltonian. This offset is part of the apparatus source.

A quartet pulse with phase $e^{-4i\Delta t}$ and area
$w^{-1}\int\kappa(t)dt=\arcsin\sqrt\epsilon$ prepares
\begin{equation}
 |\Phi_\epsilon\rangle=\sqrt{1-\epsilon}|00\rangle_A|00\rangle_R
                       +\sqrt\epsilon|11\rangle_A|11\rangle_R.
\end{equation}
The pair pump selects logical $Z$ or $X$, retaining spectrum
$(1-\epsilon,\epsilon)$ and the actual current purifier. Unitary
contraction preserves \eqref{eq:ring-decouple}; the next reset exports
that purifier and its correlations through the same lines. Fixed gate
duration requires $\kappa_{\rm peak}=O(w)$, compatible with
\eqref{eq:ring-positive} and the ordered reset parameters.

The body has finite energy-constrained storage capacity rather than
finite Hilbert dimension. Its positive-energy Gibbs partition function
is finite for every inverse temperature $\beta>0$, giving
$S\le\beta\langle H\rangle+\log Z(\beta)$. Four complex outlets,
any anomaly-cancelling partners, gap energy, drive work and stored
outputs are charged. Fixed per-cycle error does not imply vanishing
error for the whole growing history; that stronger task requires
decreasing error per cycle and increasing resources. Autonomous
covariant driving is treated below. Direct renewal exports the
current purifier's two correlated occupation records into different
species. Against a thermal reference that factorizes across species,
their redundancy gives a positive complete-output relative-entropy
rate whenever the reset and full-joint clock errors are small enough
\cite[Proposition ``Output deficit and direct renewal'']{Liu2026HAFFD}.
Improving this reset alone therefore does not approach the sharp
heat-current limit. That task requires a physical output converter
with its own work, memory and charge account.

\begin{proposition}[Neutral block compression by physical gates]
\label{prop:neutral-buffer-code}
Buffer $L$ original noisy A/R quartets in $L$ equal-circumference
slots of each physical charge species $(+,-,+,-)$. Permit local
current phases, the pair and quartet pumps above, and same-charge
quadratic mixers between connected slots. These prescribed gates
can approximate a reversible code which confines their $4^L$-dimensional
support to $m_L$ active zero modes and fixes the remaining modes,
retaining arbitrary old references and intercycle correlations, where
\begin{equation}
 m_L=\min\left\{r:\binom r{\lfloor r/2\rfloor}\ge4^L,
 \quad0\le2L-\lfloor r/2\rfloor\le4L-r\right\}
 =2L+\tfrac12\log_2L+O(1).
 \label{eq:neutral-buffer-size}
\end{equation}
No additional pure working register or outgoing species is used.
The buffer, initialization, control, offsets and eventual discharge
are physical resources.
\end{proposition}
\begin{proof}
Relabel only the existing negative-charge zero modes by
$f_i=d_i^\dagger$ and the positive-charge ones by $f_i=d_i$.
Neutrality is $\sum f_i^\dagger f_i=2L$ on $4L$ modes; this
relabeling leaves the full sea and original charge representation
unchanged. Pair pumps and same-charge mixers become one-body
rotations. Connected rotations and currents generate $U(4L)$.
The quartet becomes, with its calibrated Klein phase,
$Q_{ab;cd}=f_a^\dagger f_b^\dagger f_cf_d+\mathrm{h.c.}$.
Permutations move it to any four distinct labels.

Set $n_i=f_i^\dagger f_i$. Balanced rotations in $(a,c),(b,d)$,
followed by averaging over their sixteen current sign phases, give
$-\tfrac12(n_a-n_c)(n_b-n_d)$.
This average is implemented by a deterministic Trotter word,
without a discarded random record. Quadrilateral differences span
all two-body diagonal terms modulo currents on the fixed-number
sector: their missing coefficient arrays are
$\beta_{ij}=u_i+u_j$, for which
$\sum_{i<j}\beta_{ij}n_in_j=(2L-1)\sum_i u_in_i$.
Choose a diagonal $H_D=\sum\alpha_i n_i+\sum\beta_{ij}n_in_j$
whose finitely many hopping transition frequencies are nonzero
and distinct up to sign. Such coefficients avoid finitely many
proper hyperplanes. Vandermonde interpolation in iterated
commutators with $H_D$ isolates the real and imaginary generators
of each individual transition between occupied $2L$-subsets.
These subsets form a connected Johnson graph. Its edge generators
give the full special unitary group on the neutral sector; finite
group-commutator and Trotter words approximate each required
Givens rotation to the stated tolerance.

Enumerate the original $4^L$ pair-code basis states and inject them
into the first $4^L$ occupation strings on $m_L$ active modes with
$\lfloor m_L/2\rfloor$ particles. Fix the inactive pattern to
$2L-\lfloor m_L/2\rfloor$ particles and extend the injection to
a neutral-sector unitary. Its action is $U_L\otimes I$ on every
old reference. Stirling's formula gives \eqref{eq:neutral-buffer-size}.
For the physical embedding, the translation-invariant equal-ring
integrals commute with $P$. Its complete $P=0$ sector is precisely
the zero-mode Fock space tensored with the same spectator sea;
the control word never projects out the full field tower. The
previous relative-energy margins, with finite offsets and small
primitive amplitudes, define the full-field gates. Unitary
contraction carries a certified input-reset error through the word.
\end{proof}

For example, two quartets occupy eight zero modes but need only six
active modes after this encoding. The fixed occupations result
from reversible redistribution, not from a recurring purity supply.
Two paid banks can hold old batches while new states arrive; both
must subsequently be discharged and reused through the same four
lines. The compiler may be very long and expensive. It proves
finite prescribed controllability, not a sustainable autonomous
clock or an efficient fixed-rate layout. It also preserves the
data spectrum: concentrating it does not prepare thermal
occupations of the outlet translation generator. The next two
propositions supply finite physical loading, sorting and return;
their clock, source and repeated-inlet accounts remain distinct.

\begin{proposition}[A finite-width packet interface]
\label{prop:packet-interface}
Fix a physical charge species, $m$ paid equal-ring memories and one
of its original outlet lines. Let $v(t)$ be a smooth row of $m$
overlapping envelopes, flat outside an interior guard of $(0,\tau)$,
with $\int v^\dagger v\,dt=(1-\eta^2)I_m$, $0<\eta<1$.
A local Hermitian mixer on a finite common arc emits these memory
columns with $O(w)$ packet error and terminal amplitude
$\eta I_m+O(w)$ as the ring round-trip time $w$ decreases.
An analogous mixer captures any such incoming overlapping packet
span. All original body modes and the complementary outgoing field
remain in the canonical map.

For a finite physical carrier of length $V$ with no return during
the window, let the joint input have physical excitation energy at
most $E$, retaining arbitrary old references. With
$\Omega=2\pi/w$, $M_\Lambda=O(V\Lambda)+m$ low modes and any
fixed sufficiently large smoothness order $r$, the actual and
comparison full outputs have half trace distance at most
\begin{equation}
 C\left(\sqrt{E/\Lambda}+\sqrt{E/\Omega}\right)
 +C_{m,V,\tau,\eta,r}\left[
 (M_\Lambda+1)w(1+\Lambda)^2
 +(1+\Lambda)^{1-r}+w^{r-1}\right].
 \label{eq:packet-joint-error}
\end{equation}
The constants include the supplied envelope norms. First choose
$\Lambda$, then $w$. This is a finite-carrier, prescribed-control
result; it is not a global trace theorem for an infinite KMS rail.
\end{proposition}
\begin{proof}
Set $K(t)=I-\int_0^t v^\dagger v\,ds$, $\Phi=K^{1/2}$ and
\[
 \lambda^\dagger=v\Phi^{-1},\qquad
 A=\dot\Phi\Phi^{-1},\qquad
 H=i\left(A+\tfrac12\lambda\lambda^\dagger\right).
\]
Since $K\ge\eta^2I$ and
$A+A^\dagger=\Phi^{-1}\dot K\Phi^{-1}
=-\lambda\lambda^\dagger$, $H$ is smooth and Hermitian.
The comparison equations
\[
 \dot d=Ad-\lambda\psi_{\rm in},\qquad
 \psi_{\rm out}=\psi_{\rm in}+\lambda^\dagger d
\]
have homogeneous propagator $\Phi(t)$ and emitted columns
$\lambda^\dagger\Phi=v$. For capture, choose a column $F$ with
$\int FF^\dagger=(1-\eta^2)I$, and use
\[
 R(t)=\left[I-\int_t^\tau FF^\dagger\,ds\right]^{1/2},\quad
 \lambda=R^{-1}F,\quad A=-R^{-1}\dot R.
\]
The same identity gives Hermitian $H$; the future propagator is
$G(\tau,t)=R(t)$, so the capture column is $-F(t)$.
Its sign is an ordinary current phase, not reversed chiral transport.

In units of propagation speed one, choose a smooth arc profile
$h_w(x)=w^{-1}h_0(x/w)$ of integral one and use the actual local
potential
\begin{equation}
 \mathcal F_w(t,x)=h_w(x)W_w(t-x),\qquad
 W_w=\begin{pmatrix}
 wH&-i\sqrt w\lambda\\ i\sqrt w\lambda^\dagger&0
 \end{pmatrix}.
 \label{eq:packet-star}
\end{equation}
The internal gap is lowered for this window, with its work charged.
Along a characteristic the exact arc transfer is
$S_w=e^{-iW_w}=\left(\begin{smallmatrix}B_w&C_w\\D_w&E_w\end{smallmatrix}\right)$.
Thus the full fields obey
\[
 b(t)=B_w(t-w)b(t-w)+C_w(t-w)\psi_{\rm in}(t-w),\qquad
 \psi_{\rm out}(t)=D_w(t)b(t)+E_w(t)\psi_{\rm in}(t).
\]
In each fixed derivative norm,
$B_w=I+wA+O(w^2)$,
$C_w=-\sqrt w\lambda+O(w^{3/2})$ and
$D_w=\sqrt w\lambda^\dagger+O(w^{3/2})$.
Telescoping the noncommuting ordered products over $\tau/w$
returns gives the $O(w)$ comparison. The complete star is unitary,
so its emitted Gram matrix plus its residual-body Gram matrix is
exactly $I_m$.

The full tower also contributes. Each positive derivative of $B_w$
is $O(w)$; differentiating an ordered product sums $O((\tau/w)^j)$
terms with $j$ differentiated factors, each supplying $O(w^j)$.
Its derivatives are therefore bounded uniformly in small $w$.
Flat guards make the residual multiplier periodic on the ring.
Its nonconstant Fourier coefficients satisfy
$\|a_k\|_{\rm HS}\le C_r w^r/|k|^r$.
The emitted and captured envelopes have the corresponding uniform
smooth bounds; their $n$th body aliases retain the exact factors
$e^{\pm in\Omega t}$.

The complementary line kernel is the direct term $E_w(t)\delta(t-s)$
plus its causal delayed returns
\[
 \sum_{j\ge1}D_w(t)B_w(t-w)\cdots B_w(t-(j-1)w)
 C_w(t-jw)\delta(s-t+jw).
\]
Each coefficient and each fixed center derivative is $O(w)$,
with $O(\tau/w)$ terms. The center derivative
$\partial_t+\partial_s$ leaves both these delayed diagonals and
the comparison Volterra factor $\Theta(t-s)$ unchanged.
Fourier integration by parts therefore bounds its nonidentity
kernel by $C_r(1+|\omega-\omega'|)^{-r}$, uniformly in $w$.
Across opposite sea polarizations,
$|\omega-\omega'|=|\omega|+|\omega'|$; the squared tail above
$\Lambda$ is $O((1+\Lambda)^{2-2r})$.
On a fixed low-frequency subspace, the causal sum is a Riemann
approximation with one-particle error $Cw(1+\Lambda)^2$.
Together with the body alias bounds this controls the relative
transformation's opposite-polarization Hilbert--Schmidt block.

For the state estimate, a proof-only energy cutoff excludes field
excitations above $\Lambda$ and nonzero body excitations with
probability at most $E/\Lambda+E/\Omega$.
Gentle projection works with the old reference retained.
The remaining finite excitation space has $M_\Lambda$ modes, so
its Fock lift costs at most $M_\Lambda$ times the one-particle error.
The canonical pair-rotation decomposition bounds the vacuum-lift
error by the opposite-polarization Hilbert--Schmidt norm
\cite[Section~7.4, Theorem~31]{Derezinski2005}.
The actual normal-ordered lifts may differ by an overall phase,
irrelevant for this prescribed channel. Combining the cutoff,
low-band and full-tower estimates gives
\eqref{eq:packet-joint-error}.
\end{proof}

Capture retains the actual covariance
$C_{\rm cap}=\int F_w(t)C_{\rm in}(t,s)F_w(s)^\dagger\,dt\,ds$,
including off-diagonal packet correlations and old body--line terms.
Orthogonal wavefunctions need not have independent thermal occupations.
A quantum program coherently selecting different loaders also requires
their actual relative lift phases to be calibrated and charged.
For the same finite graph, the source comes from its full covariance
and local point splitting, and the work is
$\int\langle\partial_tH_{\rm graph}\rangle dt$,
including gap changes, contact energy and physical ground offsets.
With $C_h=\int h_0^2$, the local quadratic scale satisfies
\[
 \int\operatorname{tr}\mathcal F_w^2dx
 =C_h\left[w\operatorname{tr}H^2+2\|\lambda\|^2\right].
\]
This does not bound all higher spatial derivatives or clock resources
uniformly. Nor does \eqref{eq:packet-joint-error} alone control heat:
an old spectator occupied with probability $E/\Omega$ can export
energy $E$ while disappearing in trace norm. Heat or source convergence
requires propagated uniform energy integrability, such as a controlled
higher energy moment or explicit weighted output tails. The packet
map supplies the loading step for the following thermal converter.
Its energy estimate uses the actual carrier rather than packet
occupation numbers assigned independently of it.

\begin{proposition}[Charge-sector thermal sorting on a retained carrier]
\label{prop:thermal-buffer-sort}
Use the original four charge species and neutral gates, and retain
four physical carrier loops of round-trip time $V=L/\nu$ in their
complete finite-volume Gibbs state $\tau_{T,V}$, $T>0$.
Include its purification and all old references. Supply $L$
independently realized noisy original quartets with marginal
$\rho_S=\rho_0^{\otimes L}$, where
\[
 \rho_0=\begin{pmatrix}
 (1-q)(1-\epsilon)&0&0&(1-2q)\sqrt{\epsilon(1-\epsilon)}\\
 0&q\epsilon&0&0\\0&0&q(1-\epsilon)&0\\
 (1-2q)\sqrt{\epsilon(1-\epsilon)}&0&0&(1-q)\epsilon
 \end{pmatrix},\qquad h_\lambda=S(\rho_0).
\]
The source and complete carrier initially have a product marginal;
their external purifications are retained. There are prescribed
finite words of the original smooth gates that return the source
near its original blank and achieve, in natural units,
\begin{equation}
 \lim_{\Lambda\to\infty}\limsup_{L\to\infty}
 \frac{D(\rho_{B,\rm out}\Vert\tau_{T_*,V})}{V}=0,
 \qquad T_*=T+\frac{6\nu h_\lambda}{\pi c_{\rm ex}},
 \quad c_{\rm ex}=4.
 \label{eq:finite-sort-deficit}
\end{equation}
Implementation error can be made arbitrarily small after choosing
each finite block and cutoff $\Lambda$, with physical energy tails
controlled. The entire output $B$, its uncaptured complement and
every atypical label are present in this relative entropy.
Finite paid internal pipeline stages permit the specified throughput
without adding external outlets. This is a finite-carrier converter;
an actual stationary inlet, repeated preparation and an autonomous
reused controller require further physical constructions.
\end{proposition}
\begin{proof}
On each loop let $D=-i\partial_x$, $P_\Lambda$ its Fourier subspace
with $|\omega|\le\Lambda$, and $\chi$ a smooth cutoff flat zero in
endpoint guards of length $g$. If
$P_\Lambda\chi^2P_\Lambda u_j=g_j u_j$, retain $g_j\ge1-e$ and
set $f_j=\chi u_j/\sqrt{g_j}$. These are orthonormal smooth
envelopes for Proposition~\ref{prop:packet-interface}.
For their isometry $F$, the actual captured matrices are
\[
 C_T=F^*(1+e^{D/T})^{-1}F,\qquad h_F=F^*DF.
\]
They need not commute. A physical one-body rotation diagonalizes
$C_T$ for its spectral probabilities, whereas the final occupation
states use the eigenbasis of $h_F$. Relabel existing negative-charge
modes by $f=d^\dagger$ as in the neutral code, retaining their
energy $-h_F^T$ and the physical sea constant.

Here is an exact finite count. In total charge sector $Q$, let
$E_{\beta,Q},S_{\beta,Q},\sigma_{\beta,Q}^2$ be the energy,
entropy and energy variance of the actual carrier-sector Gibbs
matrix. For $b_Q>\sigma_{\beta,Q}$ there are at least
\begin{equation}
 \left(1-\frac{\sigma_{\beta,Q}^2}{b_Q^2}\right)
           e^{S_{\beta,Q}-\beta b_Q}
 \label{eq:thermal-sector-count}
\end{equation}
occupation states below $E_{\beta,Q}+b_Q$. Chebyshev gives this
shell's mass; each state in it has probability at most
$e^{-S_{\beta,Q}+\beta b_Q}$. Map the charge-resolved good input
eigenlabels injectively into these low-energy states with the source
blank, then complete the injection to a sector bijection.
Every bad label occupies a remaining physical state. If their
probability is $p_{\rm bad}$,
\[
 \Pr(S_{\rm out}\ne b)\le p_{\rm bad},\qquad
 S(S_{\rm out})\le h(p_{\rm bad})+4L p_{\rm bad}\log2.
\]
For the depolarizing source specified in the proposition, exact blank
restoration would violate the full-rank finite input's rank count.
No label is postselected away.

At fixed sufficiently large $\Lambda$, the source eigenvalue
strings and the Gaussian carrier's actual spectral probabilities
have information variance $O(L)$. An information margin $L^{9/10}$
and charge window $|Q|\le L^{5/8}$ lose $o(1)$ probability.
Choose the finite-bank target entropy larger than the input entropy
by $3L^{19/20}$ and put $b_Q=L^{17/20}$.
For these charges exponential tilting of the bounded occupation
Bernoullis gives $\log\Pr_\beta(Q)=-O(Q^2/L+\log L)$.
Fourier inversion bounds the tilted probability at its mean below
by $cL^{-1/2}$: its characteristic function is bounded by
$e^{-cL\sin^2(\theta/2)}$, and its scaled central part converges
to a nonzero Gaussian integral. Conditioning its $O(L)$ energy
variance therefore gives $\sigma_{\beta,Q}^2=O(L^{3/2})$ and
$|E_{\beta,Q}-E_\beta|=O(|Q|+L^{3/4})$.
The count loss in \eqref{eq:thermal-sector-count} is smaller than
the entropy margin. Hence $p_{\rm bad}\to0$, source entropy is
$o(L)$, and the carrier energy excess is $o(L)$.

The bijection uses the original gates. For a hop $a\to b$, a
quadrilateral interaction $(n_a-n_c)(n_b-n_d)$ has difference
$n_d-n_c$. Two distinct spectator subsets have the same differences
for every $c,d$ only when one is empty and the other full.
Generic current and quadrilateral coefficients therefore separate
all other oriented edges; finite Vandermonde filtering and the
connected occupation graphs give independent special-unitary
actions in sectors $2\le k\le m-2$. The linked $k=1,m-1$
sectors are unnecessary: the original source has $k_S=2L$ on
$4L$ modes, so every occupied source--carrier sector is interior.
Spectral diagonalization and the chosen bijection are finite
commutator/Givens words. Their phases are retained.

Capture and return act on the same complete carrier as
$U_{\rm load}^*U_{\rm sort}U_{\rm load}$.
For commensurate loops $U_0(V)=1$ up to its declared sea phase.
Thus the inverse of a free interval $\delta$ is the positive wait
$V-\delta$. Smooth signed current kicks invert the local mixing
steps: evolution with $H_0+(a/\varepsilon)J(F)$ for time
$\varepsilon$ approaches $e^{-iaJ(F)}$ on the smooth energy
core, with error at most
$\varepsilon\int_0^1\|H_0e^{-isaJ(F)}\psi\|ds$.
Higher weighted versions follow from smooth multiplication and
the summable opposite-sea tails. Reverse a finite Trotter word
using these waits and kicks. Their work, peak ground shifts,
relative phases and contact sources are paid. Its $G$ whole-loop
waits require $G$ internal storage stages; narrowing a contact
does not remove that latency. All stages have the original species
and only the original four external outlets.

To pass from the bank to the complete physical carrier, let $d$
be the number of discarded Fourier vectors and $M$ the bank size.
Then $d\le C(\Lambda g+1)/e$. The retained Fourier span need
not be invariant under $D$; its discarded cross block contributes
$C\Lambda d$, in addition to the cutoff error
$CM(\Lambda\sqrt e+1/g)$.
Taking $g=V^{1/2}$ and $e=V^{-1/4}$ makes these and the entropy
compression error $o(V)$ at fixed $\Lambda$.
The unselected thermal tails have energy and entropy per $V$
at most $C(1+\Lambda)^2e^{-\Lambda/T_*}$.
The ideal conjugation preserves the entropy of source plus full
carrier, so Araki--Lieb and subadditivity give
\[
 \left|S(B_{\rm out})-S(\tau_{T,V})-Lh_\lambda\right|
       \le S(S_{\rm out}).
\]
Together with the physical energy estimate and Gibbs identity,
this bounds the deficit in \eqref{eq:finite-sort-deficit} by
$o_\Lambda(1)+C(1+\Lambda)^2e^{-\Lambda/T_*}$.
Full finite-carrier Gibbs moments and the smooth gate graphs
control implementation tails; Gibbs energy-constrained entropy
continuity then gives the stipulated finite accuracy. Trace error
alone would not give this energy or entropy conclusion.
\end{proof}

\begin{remark}[Other supplied neutral source spectra]
The proof of Proposition~\ref{prop:thermal-buffer-sort} also applies
to a product source $\bigotimes_{j=1}^L\sigma_j$ from a fixed finite
alphabet of known states in the same neutral quartet code, with the
same joint source--carrier product marginal, paid gates and controlled
physical energy tails. Joint operations include the receiver and its
current, resettable purifier; external old references remain untouched.
Put $s_L=\sum_jS(\sigma_j)$ and suppose $s_L/L\to s$ and the
source information variance is $O(L)$. Replace $h_\lambda$ by $s$
in $T_*$ and $Lh_\lambda$ by $s_L$ in the entropy identity.
The finite-bank target uses the actual $s_L$ plus the same subextensive
entropy margin; its target temperature tends to $T_*$. No rate of
convergence of $s_L/L$ is needed.
Eigenlabel concentration uses this variance and entropy. The source
still has $k_S=2L$, so the neutral-sector compiler remains in its
interior sectors. The final full-carrier entropy bound uses $s_L$;
thermal counts, energy tails and paid loading/unloading are unchanged.
HAFF-D's pointer-aligned dephasing supplies such rank-two sources at
fixed noise parameters \cite{Liu2026HAFFD}. This is supplied-state
renewal, not recovery of arbitrary unknown content against a forbidden
purifier or a proof of the product marginal for an adaptive source.
\end{remark}

The source debit is the complete quartet entropy $h_\lambda$.
For $0<q,\epsilon<1/2$, its difference from the bare receiver change is the correlation
lost in the logical depolarizing step:
\[
 h_\lambda-[h(\epsilon_q)-h(\epsilon)]
 =I(A:R)_{\rm before}-I(A:R)_{\rm after}
 \ge4q\,h(\epsilon)>0.
\]
The equality uses the pure corrected logical quartet and its noisy
marginals; joint convexity for the depolarizing mixture gives the bound.
With a conserved classical axis $P$, the source entropy conditional
on $P$ is $Lh_\lambda$. Its near-blank return removes
$s^{PS}=Lh_\lambda-o(L)$ from the complete source--program system.
Entropy conservation with the independent thermal carrier gives
\[
 S(B_{\rm out}|P)-S(\tau_{T,V})
 =s^{PS}+I(S_{\rm out}:B_{\rm out}|P).
\]
No noise purifier is granted as usable side information. An autonomous
controller's sensing and renewal must also enter its complete ledger.
HAFF-D gives the corresponding conditional deficit account~\cite{Liu2026HAFFD}.

The fixed $Z$ and $X$ noisy quartets are isospectral. Their
entropy-optimal outlet heat agrees to leading order, although their
retained correlations and recovered work can differ. The direct source
contrast below belongs to its specified unsorted reset. Its causal
intervention remains valid, but does not certify persistent heat gain
for the sharp converter.

\begin{corollary}[The selected axis changes the maintenance source]
\label{cor:ring-source}
Fix $0<\epsilon<1/2$ and $0\le q<1/2$. On the same complete
post-reset state, compare the two physical axis preparations above,
followed by the logical depolarizing channel with
$\epsilon_q=\epsilon+q(1-2\epsilon)$ and the same reset. Do not undo
the $X$ rotation. There are fixed finite reset parameters for which,
uniformly after the initial preparation cycles,
\begin{equation}
 \int\zeta(t)(\langle T_{\rm out}\rangle_X-
                    \langle T_{\rm out}\rangle_Z)dt
 \ge\frac{\Delta(1-\eta^2)}2(1-2q)(1-2\epsilon)>0,
 \label{eq:ring-source-contrast}
\end{equation}
where $\zeta\ge0$ is smooth and equals one on the emitted pulse.
This is a paired causal intervention on one pre-preparation state,
not a comparison of histories selected by conditioning on different
acquired outcomes.
\end{corollary}
\begin{proof}
The ideal total zero occupations are $2(\epsilon_q+\epsilon)$
and $1+2\epsilon$. Their difference is
$D=(1-2q)(1-2\epsilon)$. Let $v_w$ be the real emitted zero-mode
envelope. The emitted $n$th mode is
$v_n=e^{-i(\Delta+n\Omega)t}v_w$, and
$J_k=\int\zeta v_w^2e^{-ik\Omega t}$, $J_0=1-\eta^2$.
The exact old-body matrix of the smeared stress is
\[
 B_{mn}=[\Delta+\Omega(m+n)/2]J_{n-m}.
\]
Its diagonal is $J_0(P+\Delta N_s+\Delta N_0)$ after sea normal
ordering. Both interventions preserve $P$; $|N_s|\le P/\Omega$ and
$\|N_0\|=4$. If $d$ bounds the half-trace reset error and $M_E$
bounds $\langle P/\Omega\rangle$, their diagonal contrast is at least
$\Delta J_0(D-16d-2M_E)$.

These errors can be controlled independently of the zero occupation.
Separating the old zero mode in the excitation-vector proof gives
\[
 \sqrt{\langle P/\Omega\rangle_{j+1}}
 \le\alpha\sqrt{\langle P/\Omega\rangle_j}
       +2(\sqrt{Z_A}+\sqrt{V_P}+\sqrt{S_P}),
\]
where $Z_A=\sum_{n\ne0}|n||a_n|^2$, $V_P\le V_A$ is the
nonzero-mode creation leakage, and $S_P$ is $S_Q$ with its zero term
removed and weight $n$ instead of $1+n$. Thus $M_E\to0$ as
$w\to0$ at fixed $\Gamma,\Delta$. Before reset take
$M_Q=4+2M_E$ as a bound on $\langle Q\rangle$.

Set $R_2=\sum_{k\ne0}(1+|k|)^2|J_k|$. Weighted Schur estimates for
equal sea polarizations and a Hilbert--Schmidt estimate across the
sea give a two-branch off-diagonal error at most
$2(\Delta+\Omega)R_2(M_Q+4\sqrt{M_Q})$.
For $m\ge5$, integration by parts gives
\[
 R_2\le2\|(\zeta v_w^2)^{(m)}\|_1\Omega^{-m}
                   [\zeta_R(m)+2\zeta_R(m-1)+\zeta_R(m-2)],
\]
where $\zeta_R$ denotes the Riemann zeta function. The envelope
derivatives are uniformly bounded as $w$ decreases.

Old body--line cross terms also remain. Let $S_f$ be the exact
line-to-line scattering operator and
$g_n=S_f^*(i\zeta\partial_t+i\zeta'/2)v_n$.
Its delay kernel contains shifts $kw$. Since
$e^{-in\Omega kw}=1$, direct substitution gives
$g_n=e^{-iE_nt}(E_na_w+b_w)$, $E_n=\Delta+n\Omega$, with all
derivatives of $a_w,b_w$ uniformly bounded independently of $n,\Delta$.
The wrong KMS occupation of this unnormalized mode is bounded by
\[
 2(E_n^2\|a_w\|_2^2+\|b_w\|_2^2)e^{-|E_n|/(2T)}
 +2(2/|E_n|)^{2m}(E_n^2\|a_w^{(m)}\|_2^2+\|b_w^{(m)}\|_2^2).
\]
Sum this over all integers $n$ and the four species to obtain $G_T$.
The sum converges for $m\ge2$. Its zero term tends to zero with
$\Delta$; its nonzero terms tend to zero with $\Omega\ge2\Delta$.
CAR and Cauchy--Schwarz bound the difference of the actual cross
stresses by $4\sqrt{4+M_E}\sqrt{G_T}$, without a product-state premise.
The incoming-only term cancels because body-only interventions leave
the same incoming marginal unchanged. The full lower bound is therefore
\begin{align}
 \Delta J_0(D-16d-2M_E)
 &-2(\Delta+\Omega)R_2(M_Q+4\sqrt{M_Q})\notag\\
 &-4\sqrt{4+M_E}\sqrt{G_T}.
\end{align}
Choose $\Gamma$ large, then $\Delta$ large, then $w$ small. The
reset bounds and the displayed convergent tails make the error less
than $\Delta J_0D/2$. These choices remain compatible with the pump
positivity condition \eqref{eq:ring-positive}.
\end{proof}

The energy difference is paid by preparing and maintaining different
physical axes at a fixed gap. It is emitted by the reset itself, not
by an additional program-tagged source. Conditional histories can
have different incoming backgrounds; the causal result makes no
claim that those observed conditional means obey the same formula.

\section{A spatial work clock for the retained thermal body}
\label{sec:physical-window-clock}

The waveforms above can be generated by a physical clock field over a
specified finite horizon. Its source depends on the retained quantum
state, including the body--line correlations created by earlier gates.
For this construction supply a static material frame $(U,V)$ and
reflecting endpoints. Let
\begin{equation}
 S_\chi=\frac I2\int\sqrt{-g}
 [(U\cdot\nabla\chi)^2-C(V\cdot\nabla\chi)^2],
 \qquad 0<C<1,
 \label{eq:work-clock-action}
\end{equation}
with Neumann data, initially $\chi=0$, $\chi_t=1$, $\chi_x=0$.
Gate profiles are smooth functions of $\chi$ and the material coordinate.
Use flat guard intervals between a zero-gap, fixed real quartet
quadrature and the quadratic gap, pair and mixer gates. The actual
gap-lowering and restoration work, offsets and finite normal-order
terms are included in the Hamiltonian. Internally mixed components
have a common positive free speed $v_0$.

\begin{proposition}[Retained-state field--clock evolution]
\label{prop:retained-window-clock}
Fix a finite list of such windows and a finite horizon $H$. Prepare
the original mixed four-mode target and its purifier, the empty-zero
body reference and an actual finite-volume Gibbs reservoir; retain
every subsequent state and reference. Assume smooth stress and kernel
remainders, sufficiently many finite energy moments, and the strict
quartet inequalities \eqref{eq:quartet-clock-margins} below.
Choose the reservoir long enough to prevent a return before $H$.
Then sufficiently large finite physical inertia $I_H$ gives a
timelike work clock and the same sequence of gate types throughout
$H$. Its force is that of the actual evolved full-CAR state.
The quadratic windows and their delayed matrix boundary close at the
same derivative order. The quartet window admits a reflected
symmetrizer at that order. In each specified lower source norm the
field and the paid clock-work source have a comparison error
$O_H(I^{-1/2})$; a quadratic window with identical incoming data has
$O_H(I^{-1})$ error. These estimates retain the clock's $O(I)$
baseline stress and its finite work contribution.
\end{proposition}

\begin{proof}
During a quadratic window the Heisenberg multiplet obeys
$(\partial_t+\sigma v_0\partial_x)\Psi=-iF\Psi$, where
$F=F(\chi,a)$ is Hermitian and $\sigma=\pm1$ is chirality.
Pair gates use constrained Majorana components; their doubling adds
no physical outlets. Subtract the fixed symmetric spatial singularity
$-i\sigma\kappa_J/(x-y)$, $\kappa_J=1/(2\pi)$, from
$\Gamma(x,y)=\langle\Psi^\dagger(y)\Psi(x)\rangle$.
Let $J=\Gamma_{\rm reg}(x,x)$ and
$K=-i\sigma(\partial_x-\partial_y)\Gamma_{\rm reg}/2|_{y=x}$.
Point splitting gives the full matrix equations
\begin{align}
 J_t+\sigma v_0J_x&=-i[F,J]-\sigma\kappa_JF_x,\notag\\
 K_t+\sigma v_0K_x&=-i[F,K]-\frac\sigma2\{F_x,J\},
 \qquad F_x=F_\chi\chi_x+F_a a_x.
 \label{eq:matrix-clock-ward}
\end{align}
The force in \eqref{eq:work-clock-action} is
$-\sum_\sigma\langle F_\chi,J_\sigma\rangle-\partial_\chi E_o$.
For complex components the contraction is a trace; for constrained
Majoranas it is half the trace. Finite sea shifts are retained:
the periodic empty-zero convention has
$[L_n^{\rm sea},J_m^{\rm sea}]=-mJ_{n+m}^{\rm sea}
-(n/2)\delta_{n+m,0}$. Passing to canonical currents uses
$J_n=J_n^{\rm sea}-\delta_{n0}/2$ and
$L_n=L_n^{\rm sea}+\delta_{n0}/12$, with the corresponding
fixed Hamiltonian subtraction. It leaves the physical zero
occupation unchanged.

For a mixing arc of length $\ell<w$, set
$d_a=\ell/v_0$, $d_{\rm out}=(w-\ell)/v_0$ and $d=w/v_0$.
The arc Wilson matrix $G(u)$ maps entry time $u$ to exit
$u+d_a$. With its body/line blocks $G_{BB},G_{BF}$ and the actual
outside phase $Q$, the incoming body satisfies
\begin{equation}
 b(u)=A(u)b(u-d)+B(u)f(u-d),\quad
 A(u)=QG_{BB}(u-d),\quad B(u)=QG_{BF}(u-d).
 \label{eq:actual-delayed-clock-boundary}
\end{equation}
Finite iteration expresses its cross kernel as
\[
 \Gamma_{BF}(u,v)=R_M\Gamma_{B_0F}(u-Md,v)
 +\sum_{j=0}^{M-1}R_jB(u-jd)
                     \Gamma_{FF}(u-(j+1)d,v),
\]
where $R_j=A(u)\cdots A(u-(j-1)d)$ and $R_0=1$.
These are kernels of the original joint state. At $v=u$ every
reservoir separation is at least $d$; all clock-dependent incoming
data precede $u$ by at least $d_{\rm out}>0$.

The incoming cross stress requires covariant relative time
differentiation. For constant exterior connections $F_{0,B},F_{0,F}$,
\[
 K_{BF}=\frac{i}{2v_0}(\partial_u-\partial_v)\Gamma_{BF}|_{v=u}
 -\frac{F_{0,B}\Gamma_{BF}+\Gamma_{BF}F_{0,F}}{2v_0}.
\]
For the full right-moving transfer put
$S=(iG'G^\dagger+GF_{0,\rm in}G^\dagger-F_{0,\rm out})/v_0$.
The complete spatial point-split conversion is
\begin{align}
 J_{\rm out}&=GJ_{\rm in}G^\dagger+\kappa_JS,\notag\\
 K_{\rm out}&=GK_{\rm in}G^\dagger
 +\tfrac12\{S,GJ_{\rm in}G^\dagger\}
 +\tfrac{\kappa_J}2S^2.
 \label{eq:covariant-matrix-transfer}
\end{align}
Distinct body/line finite-circle backgrounds are added before and
subtracted after this conversion. Differentiating unitarity twice in
the symmetric split cancels the anti-Hermitian second derivative,
including noncommuting connection terms. A scalar total stress would
lose the $BF$ data in this equation.

Differentiate the arc Wilson matrix by Duhamel's formula. Its
$(m+1)$st entry-time derivative needs only derivatives through $m$
of $\chi_t$, a first-order clock variable. The initial two-time
kernel remainders and their required derivatives are bounded away
from the declared diagonal singularity. Finite history products
therefore give a locally Lipschitz $C^m$ boundary functional in
$(\chi_t,\chi_x,J,K)$, $m\ge3$. On time steps shorter than
$d_{\rm out}$ its incoming data are already fixed. Characteristic
iteration of \eqref{eq:matrix-clock-ward} and the acoustic-clock
wave equation uses these same jets and yields the quadratic-window
estimate. The constants depend on $H/d$ and the initial norms.

For a quartet window write $v_s=v_0+s a(\chi)$,
$g_s=s a_\chi$, $s=\pm$, and $\gamma_s=1/(48\pi)$.
The actual chiral stress satisfies
$T_{\sigma s,t}+\sigma v_sT_{\sigma s,x}
+2\sigma v_{s,x}T_{\sigma s}=\sigma\gamma_s v_{s,xxx}$.
Use $Y_{\sigma s}=(1-\partial_x^2)^{-1}T_{\sigma s}$ and
$Z_{\sigma s}=\partial_xY_{\sigma s}$, retaining their means.
For $r=\chi_t$, $p=\chi_x$, the principal matrix has
$A_{rp}=-C$, $A_{rZ}=-g_s/I$, $A_{pr}=-1$,
$A_{Zp}=\sigma\gamma_sg_s$, $A_{ZZ}=\sigma v_s$.
Its symmetrizer has $H_{rr}=I$, $H_{pp}=IC$,
$H_{pZ}=g_s$, $H_{ZZ}=v_s/\gamma_s$.
Direct Schur elimination gives
\begin{equation}
 IC>2\gamma_s a_\chi^2\sum_s v_s^{-1},\qquad
 I(1-C)>2\gamma_s a_\chi^2\sum_s\frac{v_s}{1-v_s^2}.
 \label{eq:quartet-clock-margins}
\end{equation}
Thus $H$ and $H\pm HA$ are positive. The exact remainder follows
by commuting $1-\partial_x^2$ through the Ward equation; its terms
with $v_{s,xx}Z$ gain one derivative under
$\partial_x(1-\partial_x^2)^{-1}$. They are locally Lipschitz
at the same Sobolev order. At a reflected endpoint $p=0$,
$Y_{Rs}=Y_{Ls}$, $Z_{Rs}=-Z_{Ls}$, so the energy flux vanishes.
Smooth higher classical energies recover the desired lower-norm
actual stress. Increasing finite $I$ preserves the flat guard
intervals and a timelike clock over the fixed finite list.

For the solved smooth profiles, the full-tower propagator carries
the actual state through each seam. Its distributions satisfy these
same linear Ward and kernel equations; uniqueness identifies the
mean force with that retained state. Finitely higher initial graph
and classical norms can be selected for the finite list. In a gate
controlled by a conserved record projector use unnormalized weighted
kernels; each c-number anomaly carries that projector's actual
weight. The controlled unitary also retains off-diagonal record
coherence. Register updates occupy their separately paid slots.

Finally the actual clock work is
\[
 \frac d{dt}\frac I2\int(r^2+Cp^2)
 =-\int r\,[\sum_\sigma\langle F_\chi,J_\sigma\rangle
                         +\partial_\chi E_o],
\]
with the quartet force replacing the bracket during its slot.
Write $\chi=t+\eta_I/I$ and solve
$\eta_{0,tt}-C\eta_{0,xx}=-\mathrm{force}_{\rm ref}$
with the same corrected initial data. The source estimate controls
$\eta_I-\eta_0$; the finite retarded work term and the
$O(I^{-1})$ quadratic clock remainder remain in the comparison.
No growing clock baseline multiplies an uncontrolled state error.
\end{proof}

\begin{corollary}[Full joint accuracy of the finite schedule]
\label{cor:joint-window-accuracy}
Use the original charge-glued field space, its finite reservoir and all
retained quantum registers and references in
Proposition~\ref{prop:retained-window-clock}. In fixed energy units put
$G_{\rm en}=1+\sum_{\rm rails}v_0P_{\rm rail}$, the positive physical
excitation energy, including holes and every spectator. The finite
zero-mode multiplicity is retained. Suppose the joint initial state
has $M_1^2=\operatorname{Tr}(\rho_0G_{\rm en}^2)<\infty$, with the
finitely higher smooth input norms required for the schedule.
Compare its solved clock history with the prescribed history on the
same field space and in the same initial state. If
\[
 a_I(t)=\|(H_I(t)-H_0(t))G_{\rm en}^{-1}\|,\qquad
 b(t)=\|[G_{\rm en},H_0(t)]G_{\rm en}^{-1}\|,
\]
then, for every $t\le H$, the half trace distance of their full joint
quantum outputs is bounded by
\begin{equation}
 e_H=M_1\int_0^H a_I(s)
                  \exp\!\left(\int_0^s b(r)\,dr\right)ds
       =O_H(I^{-1/2}).
 \label{eq:joint-window-error}
\end{equation}
At a reset checkpoint whose prescribed-history whole-body error is
$d_{\rm ring}$ in \eqref{eq:ring-decouple}, the actual output satisfies
\begin{equation}
 \tfrac12\|\rho_{{\rm body},R}^{I}
   -|{\rm sea}\rangle\langle{\rm sea}|\otimes\rho_R^{I}\|_1
       \le d_{\rm ring}+2e_H.
 \label{eq:joint-clock-reset}
\end{equation}
Here $R$ is the actual outside system, including every old reference.
The same joint bound applies after preparing the original mixed
receiver with its current purifier and before its next specified
readout. It is an energy-constrained state estimate, not an
unconstrained diamond-norm assertion.
\end{corollary}
\begin{proof}
The local quartet is a stress smearing in its original sector
representation. The positive-energy bound
\cite[Sections~2.2--2.3]{CarpiTanimotoWeiner2022}
\[
 \|L_n\psi\|\le\sqrt{1+c/12}(1+|n|^{3/2})
                         \|(1+L_0)\psi\|
\]
and Fourier Cauchy--Schwarz control a profile difference by its
$H^m$ norm, $m>2$. The energy commutator needs one further derivative.
For a quadratic current, pair or mixer, separate the same-sea number
part from the pair part. The first is bounded relative to excitation
number and hence to $G_{\rm en}$ on the fixed finite hardware. A
smooth compact-arc profile gives pair matrix entries bounded by
$C_s(1+E_n+E_m)^{-s}$. A sum-frequency shell has $O(q)$ entries:
absolute summation requires $s>2$, and the energy commutator requires
$s>3$. Three and four integrable spatial derivatives, respectively,
are safe integer choices. Different circumferences enter through
their physical energies and finite spacing constants; an unweighted
sum of $L_0$'s would not supply this mixer commutator. The estimates
act across the original CAR sectors, including record coherences and
their scalar phases. They use no Gaussian replacement at a seam.

The smooth propagators preserve the common energy core. For a
purification $\Psi_0$ of the actual initial state, the commutator
estimate gives
$\|G_{\rm en}U_0(s)\Psi_0\|\le M_1\exp(\int_0^s b)$.
Duhamel's formula then bounds
$\|(U_I(t)-U_0(t))\Psi_0\|$ by \eqref{eq:joint-window-error};
every reference carries the identity. The clock comparison at the
chosen lower coefficient norm gives
$\int_0^H a_I\le C_H I^{-1/2}$, after supplying the required higher
norms. Its existence interval already covers the entire schedule.
Triangle inequality with the prescribed reset costs $e_H$; replacing
its outside marginal by the actual one costs another $e_H$.
Re-preparation and the same subsequent instrument contract this error.
\end{proof}

The delay in reset bounds is the round-trip time $d$, whereas $w$
is spatial length. For example, $v_0=.5$, $w=10^{-8}$,
$\Delta=10^7$, $T=\tau=1$ and $\Gamma=32$ admit a smooth
flat reset pulse with $D_0<9.9\times10^8$ and
$r_d<4\times10^{-5}$. The four-mode joint reset bound is
$0.026503$. Finite additional clock accuracy is obtained by
increasing $I_H$; this argument does not provide its numerical value.
Restoring the gap preserves the original $X/Z$ work difference
$\Delta(1-2q)(1-2\epsilon)$. For a constant arc background
$F_0=\operatorname{diag}(\Delta,0)$ the actual local connection
\[
 F=F_0+v_0h(a)e^{-iF_0a/v_0}
                 W(\chi-a/v_0)e^{iF_0a/v_0},\qquad \int h\,da=1,
\]
gives the full circuit phase $e^{-i\Delta d}$, retaining its
spatial force and emitted carrier energy.

The evolving work clock still has a supplied frame/endpoints;
traction and geometric reaction need their own equations. Equal mixer
speeds, smooth inputs and separated quartet/current slots are required.
Horizon-dependent inertia and reservoir length leave fixed-capacity
sharp exhaust open.

\subsection{A local gate on a massive body}
\label{sec:guarded-mass-gate}

A spatial body need not have the massless ring's energy graph.
Pay instead for a finite physical circle with a mass-zero central
interval of length $\ell$ and a positive smooth Dirac mass throughout
the rest of its return arc. Its four thermal outlet fields remain
separate, with the original charges and capacity. Both mass layers,
the exterior polarization, support pressure and initial preparation
belong to the apparatus. This layout replaces a compact barrier
followed by a long gapless body exterior; it does not give that old
layout a gap by assumption.

\begin{proposition}[A guarded local gate with full physical energy control]
\label{prop:guarded-mass-gate}
Let $h_M=-iv\sigma_z\partial_x+M(x)\sigma_x$ on this circle.
There is a nonempty smooth mass-profile class with a simple lowest
positive mode $E_0$, no zero mode and next positive energy $E_1>E_0$.
Set $H_M=d\Gamma(|h_M|)$, with particles and holes above the
actual filled mass sea, and $A=1+H_M$. Let $h\ge0$ be smooth,
supported strictly inside the mass-zero region. On the original
charge-glued space define the local vertex
\[
 Q_+(h)=\int h(x)
 [R_1^\dagger R_2^\dagger R_3^\dagger R_4^\dagger+
  L_1^\dagger L_2^\dagger L_3^\dagger L_4^\dagger]dx,
 \qquad Q=Q_++Q_+^\dagger.
\]
All modes and old references are retained. For each fixed graph
order $s$, sufficiently weak smooth resonant drives
\begin{equation}
 H_{\rm lab}(t)=H_M+\lambda b(\lambda t)
 [e^{-4iE_0t}Q_+(h)+e^{4iE_0t}Q_+(h)^\dagger]
 \label{eq:guarded-local-pump}
\end{equation}
implement a prescribed full logical angle on the four-ground-mode
code $P_0$, in physical time $T=U/\lambda$, with
\begin{equation}
 \|A^s(U_{\rm lab}(T)-U_{\rm target,lab}(T))P_0\|
 \le C_s\lambda.
 \label{eq:guarded-full-angle}
\end{equation}
The target includes its calibrated laboratory phase.
The error is uniform in external old references, not an
unrestricted-energy diamond bound. Smooth physical mass-force
and stress tests satisfy
$\|A^{-4}S_fA^{-4}\|<\infty$, so weighted state errors also
control their actual renormalized expectations at this fixed mass.
\end{proposition}
\begin{proof}
The original orthogonal boson rotation and re-fermionization give,
for each chirality, the operator identities
\[
 Q_\sigma(h)=\pi^{-1}(T_{\sigma,1}(h)-T_{\sigma,2}(h)),\qquad
 T_{\sigma,1}+T_{\sigma,2}+\sum_{a=1}^3T_{\sigma,a}
      =\sum_{i=1}^4T_{\sigma,\mathrm{Dirac},i}.
\]
The individual rotated stresses need not be quadratic in the
massive species. Their total on the right is quadratic.
The commuting sector stresses have positive shifted forms by
the chiral QEI~\cite[Theorem~4.1 and Remarks~2--3]{FewsterHollands2005}.
Their joint spectral calculus therefore bounds $\|Q(h)\psi\|$
by $C\|(K_h+C_h)\psi\|$, where $K_h$ is the actual original
localized kinetic stress. Global charge gluing is retained;
no independently prepared sector vacua are supplied.

In the physical mass polarization its one-particle operator is
\[
 D_h=-iv\sigma_z(h\partial_x+h'/2)
     =h h_M-iv\sigma_z h'/2,
\]
since $hM=0$. For a number block $b$ and $\omega=|h_M|$,
$b^2\le C^2\omega^2$ and hence $|b|\le C\omega$.
On an $n$-particle tensor space, different-particle factors
commute, so
\[
 (\sum_i b_i)^2\le\sum_i b_i^2+
     2\sum_{i<j}|b_i||b_j|
       \le C^2(\sum_i\omega_i)^2.
\]
This gives a uniform Fock bound without a spurious $\sqrt n$.
The number-block closure follows from the commutator criterion
with $1+\omega^2$ on the smooth elliptic core. On the guard all
mass jets vanish. The complete symbols of the spectral projections
there agree with the massless chiral projections; consequently
$P_+D_hP_-$ is smoothing. Its actual mass-eigenbasis coefficients
decay faster than every fixed power of $E_n+E_m$.
Number and pair estimates give $QA^{-1}$ bounded.
For signed derivatives use a wider guarded positive cutoff,
not a bound on $h'/h$.
Local translation gives
$\operatorname{ad}_{H_M}^jQ(h)=(-iv)^jQ_R(h^{(j)})+
(iv)^jQ_L(h^{(j)})$.
Thus every fixed required energy commutator is relatively bounded
and the smooth weak-drive propagator preserves the corresponding
graphs over its entire gate time.

The physical spectral class can be constructed, rather than
postulated. For a step reference with mass $m$ on the exterior
arc of length $d$, the subthreshold parity equations are
\[
 \tan\frac{E\ell}{2v}
 =\sqrt{\frac{m-E}{m+E}}\tanh\frac{d\sqrt{m^2-E^2}}{2v}
 \quad(+),
\]
with the negative reciprocal square-root on the right for parity
$(-)$. If $m>3\pi v/(2\ell)$, the first positive root is simple,
$0<E_0<\pi v/(2\ell)$, while $E_1>\pi v/\ell$.
At zero the transfer is $\exp((\int M/v)\sigma_y)$ and has no
unit eigenvalue. Smooth the two steps just outside the body.
Transfer Duhamel bounds the perturbation on fixed complex contours
by $C\|M_{\rm smooth}-M_{\rm step}\|_1$; choosing it below the
positive minimum of the explicit step characteristic function
preserves these root counts and a strict gap. This gives a real
smooth class, without treating a sampled eigengrid as its proof.
In the body its normalized ground mode has
$R=e^{iE_0x/v}/\sqrt{2Z}$ and $L=e^{-iE_0x/v}/\sqrt{2Z}$.
For $\operatorname{supp}h\subset(-\ell/8,\ell/8)$,
\[
 g_4=\frac1{2Z^2}\int h\cos(4E_0x/v)dx>0,
 \qquad g_2=\frac1Z\int h\cos(2E_0x/v)dx>0.
\]

Let $N$ be signed flavour number, particles minus holes.
It is not excitation number or total electric charge.
Then $K_0=H_M-E_0N\ge0$ has the sixteen-dimensional kernel
$P_0$ and gap $\delta=\min(E_1-E_0,2E_0)>0$; its graph is
equivalent to $A$. Because $[N,Q_+]=4Q_+$,
$e^{-iE_0tN}$ transforms \eqref{eq:guarded-local-pump} exactly
to $K_0+\lambda b(\lambda t)Q$.
This is a proof coordinate, not a physical global-number pulse.
The compression is $P_0QP_0=g_4C_4^\dagger+\overline g_4C_4$;
all off-code terms remain in the laboratory Hamiltonian.

On $|z|=\delta/2$ the weighted resolvent norms
$\|A^jQ(K_0-z)^{-1}A^{-j}\|$ are finite.
Choose $\lambda\|b\|_\infty$ within their Neumann tube and
the physical relative-bound margin. The contour projection and
its Kato transport give an effective code matrix
$aP_0QP_0+O_s(a^2)$, $a=\lambda b$.
Its integrated remainder is $O_s(\lambda)$ because
$\int_0^T a^2dt=\lambda\int_0^U b^2du$.
The transport residual costs
$\int_0^T|\dot a|dt=\lambda\operatorname{TV}(b)$.
Energy propagation costs $\exp(C_s\int|b|du)$, independent
of $\lambda$. Choosing $\int bdu=\theta/g_4$ proves
\eqref{eq:guarded-full-angle}. Pair resonance at $2E_0$ and
guarded currents supply the paired-code rotations and phase.
For $\theta=\arcsin\sqrt\epsilon$ the vacuum preparation is
the original mixed A/R target, at final energy $4E_0\epsilon$
up to the weighted error, supplied by the pump.

Finally a smeared physical mass bilinear has order zero and a
metric stress order one. A safe pair coefficient bound is
$C_f(1+E_n+E_m)^d$, $d\le1$.
Weighting both sides by $A^{-4}$ bounds each creation term
by $(1+E_n+E_m)^{-4}$. One-dimensional energy-shell counting
has $O(q)$ pairs, so the absolute sum converges.
Number terms obey the previous estimate. With one fixed
Hadamard subtraction scheme, for
$d_4=\frac12\|A^4(\rho-\sigma)A^4\|_1$ this yields
$|\operatorname{Tr}S_f(\rho-\sigma)|\le2\|A^{-4}S_fA^{-4}\|d_4$.
Counterterm changes between backgrounds are not cancelled by
this fixed-background comparison.
\end{proof}

This gate acts on the full massive carrier, with the actual pump
phase, work and wall source. It does not use a hardware projector
onto four smeared modes. Its source certificate is for smooth
tests at a fixed mass, not pointwise source jets on moving walls.
The massive carrier's outlet mixer and whole-sea renewal must be
computed on its own spectral generator; the earlier massless
ring's reset coefficients cannot be reused. Nor are its ground
modes already the massless short-ring storage modes in
Proposition~\ref{prop:thermal-buffer-sort}. Physical transfer,
moving support and a reused work clock must join these layouts
before they form one sustained device.

\section{A finite conserved control core}
\label{sec:conserved-core}

A finite prescribed circuit, energy supply and thermal ports can share
a time-independent Hamiltonian. We combine energy-ladder
control~\cite{Aberg2014} with a finite clock~\cite{Woods2017},
using complete energy sectors and an exactly solvable history clock.
Register, gate permissions and initial resources are supplied inputs.

\subsection{A guarded battery for the whole circuit}

Let $H_D=\delta\sum_a q_a|a\rangle\langle a|$, with integers
$0\le q_a\le Q$, and let
$H_B=\delta\sum_{n=0}^K n|n\rangle\langle n|$, $K\ge2Q$.
The register $D$ includes its program, retained system and all thermal ports;
an inert reference can be appended. Define
\begin{equation}\label{eq:finite-lift}
 V(U)|a,E-q_a\rangle
   =\sum_b U_{ba}|b,E-q_b\rangle,\qquad Q\le E\le K,
\end{equation}
for any unitary $U$, and take the identity on incomplete energy sectors.

\begin{proposition}[Finite energy-conserving circuit lift]
\label{prop:finite-lift}
The operator $V(U)$ is unitary, commutes with $H_0=H_D+H_B$, and
$V(U)V(W)=V(UW)$. Put $K=2Q+L-1$ and prepare
\[
 |\beta_L\rangle=\sqrt{\frac2{L+1}}
 \sum_{j=1}^L\sin\frac{\pi j}{L+1}|Q+j-1\rangle,
 \qquad e_B=2\sin\frac{\pi}{2(L+1)}.
\]
For any input of $D$ entangled with a reference and initially independent of
this battery, the joint half trace distance between the lifted circuit
$V(U_m)\cdots V(U_1)$ and $U_m\cdots U_1\otimes I_B$ is at most
$\min\{1,2Qe_B\}$. The comparison retains the battery and all correlations.
\end{proposition}
\begin{proof}
Every complete energy sector carries a copy of the entire data space;
the others carry the identity. This proves unitarity, conservation and composition.
Extend the battery vector by zero to a bilateral ladder with shift $S$.
Its adjacent overlap is $\cos[\pi/(L+1)]$, so
$\|(S^k-I)\beta_L\|\le |k|e_B$.
With $D_q=\sum_a|a\rangle\langle a|\otimes S^{-q_a}$,
$V(U)=D_q(U\otimes I)D_q^*$ on the guarded inputs.
Each dressing difference has vector norm at most $Qe_B$, even with a retained
reference. Their sum bounds the joint state distance after purification.
The guards make this bilateral calculation agree with the finite-sector
action; no shift wraps through an endpoint.
\end{proof}

The bound concerns the whole compiled circuit at a fixed register energy range.
Fresh thermal units in a growing protocol enlarge $Q$. Battery coherence,
mean energy and preparation remain resources; the battery is not an exact catalyst.

\subsection{Autonomous timing and a usable output window}

Let $V_j=V(U_j)$ and add clock states $|0\rangle,\ldots,|m\rangle$.
For $J>0$ and $\Lambda=Jm$, set
\begin{equation}\label{eq:positive-history-clock}
 H=H_0+\Lambda I+
 \sum_{j=1}^m J\sqrt{j(m-j+1)}
 \bigl(|j\rangle\langle j-1|\otimes V_j+\mathrm{h.c.}\bigr).
\end{equation}
The clock-dependent interaction is part of the apparatus energy.

\begin{proposition}[Positive autonomous execution]
\label{prop:positive-history-clock}
The Hamiltonian is nonnegative and commutes with $H_0$. From clock state
$|0\rangle$, at $T=\pi/(2J)$ it gives
\[
 e^{-iHT}(|0\rangle\otimes|\psi\rangle)
 =e^{-i\Lambda T}(-i)^m|m\rangle\otimes
 e^{-iH_0T}V_m\cdots V_1|\psi\rangle .
\]
If only the first $r\le m$ edges perform gates and the rest are identities,
then throughout $[T/2,3T/2]$ the data's half trace distance from the completed
lifted circuit, including its free evolution, is at most
\begin{equation}\label{eq:idle-window}
 P_{m,r}=\Pr\{\operatorname{Bin}(m,1/2)<r\}.
\end{equation}
The joint data--clock distance to the completed data and the normalized
completed-tail clock state is at most $\sqrt{P_{m,r}}$.
\end{proposition}
\begin{proof}
Put $W_j=V_j\cdots V_1$, $W_0=I$, and
$\mathcal W=\sum_j|j\rangle\langle j|\otimes W_j$.
Conjugation by $\mathcal W^*$ reduces the hopping term to the spin-$m/2$
chain $2JJ_x$, commuting with $H_0$. Its spectrum lies in $[-Jm,Jm]$;
the shifted block lies in $[0,2\Lambda]$.
The spin rotation transfers $0$ to $m$ at $T$, as in perfect-transfer
chains~\cite{Christandl2004}. Its position law is
$\operatorname{Bin}(m,\sin^2 Jt)$. Every prefix $j\ge r$ has the
same completed data. Tracing the clock leaves at most the incomplete
probability as error. Projection onto the completed tail has squared
overlap one minus this probability, giving the square-root joint bound.
On the stated window $\sin^2 Jt\ge1/2$.
\end{proof}

The joint target uses the normalized tail, not the unchanged full clock
wavefunction; the latter comparison instead has bound $2\sqrt{P_{m,r}}$.
For $r=10$, $m=64$ and $\Lambda=.4$,
$P_{m,r}\approx1.77111\times10^{-9}$ and $T\approx251.327$ in the declared energy units.
The clock has 65 states. This finite window is not permanent halting:
the closed clock eventually returns and reverses its gates. Its initialization,
idle states and any output receiver remain charged resources.

\subsection{Two cycles through the same thermal contact}

The core has seven qubits: detector $S$, retained body $Q$, contact $F$,
output cells $E_1,E_2$ and dephasing cells $D_1,D_2$. The body and dephasing
cells have zero bare gap. The other four have gap $\delta$, so the register
has maximum charge 4. Prepare $S,F,E_1,E_2$ independently in
$\tau=\operatorname{diag}(1-\epsilon,\epsilon)$,
$\beta\delta=\ln[(1-\epsilon)/\epsilon]$, and each $D_i$ in its
zero-gap Gibbs state $I/2$. The body may be entangled with a retained
reference and is never re-prepared.

Cycle $i$ consists, in order, of
\[
 e^{-i\theta_i X_S P_Q},\quad P_Q=|1\rangle\langle1|_Q,\quad
 \operatorname{CZ}_{S,D_i},\quad
 \operatorname{CX}_{S,Q},\quad
 \operatorname{SWAP}_{S,F},\quad
 \operatorname{SWAP}_{F,E_i}.
\]
The dephased detector directly controls the body operation. Each gate is
lifted and placed on an edge of Eq.~\eqref{eq:positive-history-clock};
no external switching times enter. The first swap transfers the old
detector and all its correlations to $F$, putting an independent Gibbs
state in $S$. The second moves the old state to $E_i$ and restores $F$
from the unused Gibbs cell. Both conserve bare energy. Replacement is exact
against every retained reference, including the battery, and returns a
thermal state rather than a pure blank.

The first recorded bit has probability $\epsilon$ when the body is $|0\rangle$
and $\epsilon+(1-2\epsilon)\sin^2\theta_1$ when it is $|1\rangle$.
This is a genuine body reading. The command acts on the already retained
body and can damage its initial coherence; it is not a coherence-preservation claim.
For $\epsilon=.1$, $\delta=.1$, $\theta_1=\pi/3$ and $\theta_2=\pi/2$,
with the body initially half of a Bell pair, output populations are $.4,.26$.
Net heat $Q_E=.046$ equals the battery's mean energy loss. The joint output obeys
\begin{equation}\label{eq:core-output-heat}
 \beta Q_E=S(\rho_{E_1E_2})-2S(\tau)
       +D(\rho_{E_1E_2}\Vert\tau\otimes\tau).
\end{equation}
This follows from $-\ln\tau=\beta H_E+\ln Z$. Other ports, the body,
reference, clock and battery remain in the full entropy account.

At $L=31$ the battery has 39 levels and mean initial energy $1.9$.
The two sensing rotations give a joint error bound
$(\sin\theta_1+\sin\theta_2)e_B=.183123$; the full purification gives
the smaller upper bound $.070347$. For $\Lambda=.4$, the complete
Hamiltonian has cutoff $5.0$ and initial mean $2.34$. At $L=127$,
the purification upper bound is $.017610$, mean energy $7.14$ and
cutoff $14.6$. These are executions retaining the complete correlated state.

The cells hold a final energy $.066$, including $.02$ initially present;
net heat $.046$ is not automatically emitted radiation. Spatial transport,
release recoil and a stationary thermal outlet are further operations.
HAFF~D gives a material-support branch compatible with the retained core's
mean energy~\cite{Liu2026HAFFD}. It changes the support's dilaton coupling;
the spin cells are not identified with the preceding continuum.

\section{A local outlet with an autonomous clock}
\label{sec:local-outlet}

Heat stored in a cell can leave through a local Hamiltonian. Let the
output cell be site 0 of a semi-infinite fermion chain,
\begin{equation}\label{eq:local-outlet}
 H_{\rm lead}=\omega\sum_{j\ge0}n_j
       +g\sum_{j\ge0}(c_j^\dagger c_{j+1}+c_{j+1}^\dagger c_j),
 \qquad \omega>2g>0.
\end{equation}
Finite open chains have a nonnegative many-body Hamiltonian and spatially
nearest-neighbour coupling. Equivalently they are open spin chains with
XY exchange under the Jordan--Wigner map; the reference-uniform cell
channel below uses this tensor-product realization. The infinite formula denotes its local
dynamics, not a finite total-energy thermal bath. Initially site 0 has
population $p$ and may be correlated with old references. The
$\beta$-KMS bath on $j\ge1$ is initially independent of that entire
cell--reference state. Use as comparison the same bath and the
site population $\epsilon=(1+e^{\beta\omega})^{-1}$.

\begin{proposition}[Spatial export and reduced thermalization]
\label{prop:local-outlet}
Write $w_j(t)$ for the one-particle propagator from site 0.
The excess covariance and on-site energies are
\begin{equation}\label{eq:outlet-energy}
 \Delta\Gamma(t)=(p-\epsilon)|w(t)\rangle\langle w(t)|,
 \qquad \Delta E_j=\omega(p-\epsilon)|w_j(t)|^2.
\end{equation}
For the homogeneous chain the excess hopping energy is zero, and
the energy beyond the cell is
$\omega(p-\epsilon)(1-|w_0|^2)$.
On the semi-infinite chain,
\begin{equation}\label{eq:outlet-survival}
 \alpha(t):=w_0(t)=e^{-i\omega t}\frac{J_1(2gt)}{gt}.
\end{equation}
Let $b(t)$ be the population produced at site 0 from the bath alone and
$n_{\rm base}(t)=|\alpha|^2\epsilon+b(t)$.
Against every old external reference, the reduced cell's half trace
distance from its Gibbs replacement is at most
\begin{equation}\label{eq:outlet-reference-error}
 \min\{1,\ |\alpha|^2+|\alpha|/2
                    +|n_{\rm base}(t)-\epsilon|\}.
\end{equation}
The retained lead and its initial purification are excluded from that
product replacement claim.
\end{proposition}
\begin{proof}
Quadratic fermion evolution sends a covariance to
$\Gamma(t)=U(t)\Gamma(0)U(t)^*$, giving the rank-one difference.
On a bipartite chain successive components of $w$ have relative phase
$-i$ after removing $e^{-i\omega t}$, so
$\operatorname{Re}(w_j^*w_{j+1})=0$.
The surface spectral density is semicircular; its Fourier transform gives
Eq.~\eqref{eq:outlet-survival}. This is standard tight-binding propagation
with non-Markovian returns~\cite[Sec.~III]{RufeilFiori2009}.
The thermal bath is even in fermion parity. The cell channel therefore
sends population $p$ to $|\alpha|^2p+b$ and off-diagonal entry
$\rho_{01}$ to $\alpha^*\rho_{01}$.
For a state with reference blocks $\rho_{ab}$, the diagonal difference
from replacement at population $n_{\rm base}$ is
$|\alpha|^2[(1-\epsilon)\rho_{11}-\epsilon\rho_{00}]$, of trace norm
at most $|\alpha|^2$. Positivity gives $\|\rho_{01}\|_1\le1/2$.
Adding the change from $n_{\rm base}$ to $\epsilon$ proves the bound.
\end{proof}

The lead need not be exactly Gibbs for the bare cell gap. The bath band
has width $4g$, and the derivative of the Fermi function is bounded by
$\beta/4$, so its occupation differs from $\epsilon$ by at most
$\beta g/2$ throughout that band. Narrowing $g$ controls this mismatch
but lengthens transport. Equation~\eqref{eq:outlet-survival} decays as
$t^{-3/2}$ with returns and zeros; energy export is not monotone.
Finite leads eventually recur. For a chain with sites $0,\ldots,N$,
the finite and infinite wavevectors agree through Taylor order $N$.
For $N+1>2gt$, a norm bound on their difference is
\[
 \frac{2(2gt)^{N+1}}{(N+1)!}
       \left(1-\frac{2gt}{N+2}\right)^{-1}.
\]
This bounds propagation, not a global distance between finite and
infinite thermal states.

For the two output cells with populations $.4,.26$, use two identical
leads with $\omega=.1$, $g=.01$, $\beta=\ln9/.1$.
Their total initial excess is $.046$. At the first surface-amplitude
zero, $t=191.5853$, this excess is outside the cells. A finite calculation
with 256 bath sites per lead gives $n_{\rm base}=.10142988$.
At $t=2000$, excess energy $.04599817$ is in the leads and $.04422294$
is beyond site 16. Each bath initially contains energy $2.64854457$.
These values are transport after circuit completion; they do not discard
the output correlations or identify a lattice energy profile with a null
CFT stress tensor.

Reduced thermalization also differs from the exact SWAP of
Appendix~\ref{sec:conserved-core}. A two-site passive exchange at its
perfect-transfer time is an iSWAP. Acting on $|+\rangle\langle+|\otimes
\operatorname{diag}(.9,.1)$, it gives the cell exactly that Gibbs marginal,
but the joint half trace distance from the product of its two marginals
is $.18$. The bath cannot be called fresh for later cell--bath encounters.

\subsection{The outlet can remain on during computation}

Let $H_{\rm hist}$ be Eq.~\eqref{eq:positive-history-clock}, with only its
first $r$ edges active, and let $H_B\ge0$ be a finite bath Hamiltonian.
For a local outlet interaction $V_o$ with $\|V_o\|\le v$, take
$H_{\rm tot}=H_{\rm hist}+H_B+V_o+vI$.
It is nonnegative and time independent. The offset $v$ is part of its
mass and energy account.

\begin{proposition}[Always-on outlet bound]
\label{prop:autonomous-outlet}
Put $T=\pi/(2J)$, choose $0<t_0<T$, and define
\[
 P_0=\Pr\{\operatorname{Bin}(m,\sin^2 Jt_0)\le r\},\qquad
 A_r=J\sum_{j=1}^r\sqrt{j(m-j+1)}.
\]
Compare the actual joint state at $t\in[t_0,T]$ with completed lifted
data, including free evolution to $t_0$, and the free clock wavefunction
at $t_0$, subsequently evolved with the outlet and identity clock edges.
Keeping every reference, their half trace distance is at most
\begin{equation}\label{eq:autonomous-outlet-bound}
 \min\{1,\ vt_0+2\sqrt{P_0}+4A_r(t-t_0)\sqrt{P_0}\}.
\end{equation}
\end{proposition}
\begin{proof}
Removing the always-on outlet only up to $t_0$ costs at most $vt_0$ by
Duhamel's formula. Replacing all prefixes by the completed data costs
at most $2\sqrt{P_0}$ in purification norm. After $t_0$, compare with
the same outlet and free clock; only the first $r$ history edges differ.
The free-clock probability of positions $\le r$ is at most $P_0$ through
$T$. The norm of the differing edges on that state is at most
$4A_r\sqrt{P_0}$. A second Duhamel bound gives the last term.
No independence of the output cells or later fresh-bath composition is
used. The index $\le r$ includes the last active edge.
\end{proof}

For any finite transport window the bound can be made small by increasing
the idle clock and its energy. For example $r=10$, $m=10^{10}$, $J=.001$,
$\sin^2 Jt_0=128/m$ and two outlets with $v=.02$ give
$t_0=.1131371$, $T=1570.7963$ and the sufficient joint bound $.002262742$.
This uses $10^{10}+1$ clock states and clock mean energy $10^7$;
it proves finite compatibility, not an efficient clock. A direct smaller
Hamiltonian with 17 clock states, five battery levels and outlet norm
$10^{-5}$ gives joint purification distance $.00015225$ below its
certificate $.00025067$. Thus spatial export need not be externally
switched, but a fixed-size indefinitely operating exhaust is not obtained.

\section{Thermal sources and their preparation}
\label{sec:thermal-contact}

A conformal fermion defect gives a local relativistic alternative to the
lattice outlet. The real-time defect construction and its thermal
current are established results~\cite[Sections~2, 3 and 5]{BernardDoyonViti2015}.
Take two independent copies of the Majorana construction, so each port
is a complex chiral fermion with $c=1$. In the folded description, a
constant contact angle acts on the two fields at the same contact time:
\begin{equation}\label{eq:thermal-contact}
 \begin{pmatrix}\psi_1^{\rm out}\\\psi_2^{\rm out}\end{pmatrix}
 =S_\theta\begin{pmatrix}\psi_1^{\rm in}\\\psi_2^{\rm in}\end{pmatrix},
 \qquad S_\theta=\begin{pmatrix}
 \cos\theta&-i\sin\theta\\-i\sin\theta&\cos\theta
 \end{pmatrix},\quad \chi=\sin^2\theta.
\end{equation}
The inputs are jointly independent, zero-charge quasifree thermal
streams in one specified time coordinate, at $T_H<T_o$. This assumes a
prepared reservoir, rather than an already prepared outgoing field.

\begin{proposition}[Thermal transfer and its retained correlations]
\label{prop:thermal-contact}
For $\omega>0$ put $n_T=(1+e^{\omega/T})^{-1}$. The source-port
occupation and transferred energy current are
\begin{equation}\label{eq:thermal-contact-output}
 n_\chi=(1-\chi)n_{T_H}+\chi n_{T_o},\qquad
 q_\chi=\frac{\pi\chi}{12}(T_o^2-T_H^2).
\end{equation}
Its cross covariance with the return port has modulus
$\sqrt{\chi(1-\chi)}(n_{T_o}-n_{T_H})$. For $0<\chi<1$ the source
is not a KMS state at any single temperature; at $\chi=1$ it is exactly
the supplied hot state. Defining its stationary marginal entropy-current
excess, with $h(x)=-x\ln x-(1-x)\ln(1-x)$, by
\[
 s_\chi=\frac1\pi\int_0^\infty
       [h(n_\chi)-h(n_{T_H})]d\omega,
\]
one has
\begin{equation}\label{eq:contact-entropy-floor}
 q_\chi\ge T_Hs_\chi+\frac{3s_\chi^2}{\pi},
\end{equation}
strictly for $0<\chi<1$.
\end{proposition}
\begin{proof}
Rotate the full input covariance with $S_\theta$. The factor $1/\pi$
in the integrals includes particles and holes of a complex chiral
fermion; $\pi^{-1}\int\omega n_Td\omega=\pi T^2/12$.
At high frequency $n_\chi\sim\chi e^{-\omega/T_o}$ for $\chi>0$,
whereas a single-temperature Fermi occupation has tail coefficient one.
For fixed energy the Gibbs variational principle maximizes entropy at
$T_{\rm eff}=\sqrt{T_H^2+12q_\chi/\pi}$. Hence
$s_\chi\le\pi(T_{\rm eff}-T_H)/6$, with equality only for that thermal
state. Eliminating $T_{\rm eff}$ gives the bound and its strictness.
\end{proof}

The return occupation is $\chi n_{T_H}+(1-\chi)n_{T_o}$. Both outputs
must be kept when accounting for entropy and correlations. This is not
thermal replacement of an arbitrary qubit against every retained
reference. If the inputs have a cross covariance, even their thermal
marginals do not imply Eq.~\eqref{eq:thermal-contact-output}.
A finite matched-frequency check uses the positive Hamiltonian
$H=\omega(n_1+n_2)+g(c_1^\dagger c_2+c_2^\dagger c_1)$, $0<g<\omega$.
It commutes with the free energy and transfers equal energy out of the
reservoir and into the source. This mode check does not establish a
continuum limit for a microscopic moving contact.

Changing the contact has a cost in the same output fields. For a smooth
prescribed $\theta(t)$ the full unequal-time covariance is
$C_{\rm out}(t,t')=S(t)\operatorname{diag}(C_H,C_o)(t-t')S(t')^\dagger$.
The short-distance kernel is
$C_T(s)=iT/[2\sinh(\pi Ts)]$.
At $t\pm s/2$, its common vacuum term is multiplied by
$\cos[\theta(t+s/2)-\theta(t-s/2)]$.
Point splitting therefore gives an added current $\dot\theta^2/(4\pi)$
in each output, so the total work supplied to the fields is
\begin{equation}\label{eq:contact-pump}
 W_{\rm pump}=\frac1{2\pi}\int\dot\theta^2dt
       \ge\frac{(\Delta\theta)^2}{2\pi\Delta t}.
\end{equation}
This agrees with the energy-shift dissipation form in adiabatic
scattering~\cite[Eq.~(8)]{Avron2001}; here point splitting derives it
exactly for the specified energy-independent contact.
The inequality is Cauchy--Schwarz. A cosine ramp gives
$W_{\rm pump}=\pi(\Delta\theta)^2/(16\Delta t)$.
Both ports' pump energy is charged; it is separate from the static net
thermal transfer. A small angle error alone is no source certificate:
$\theta=\theta_0+\alpha\sin\Omega t$ has mean total pump current
$\alpha^2\Omega^2/(4\pi)$.
Instantaneous thermal entropy formulas do not apply during this switch.
The controller that supplies this work and the contact's material stress
still require their own realization. The input preparation has the
following finite-window alternative.

\subsection{Thermal correlations from a finite-energy preparation}
\label{sec:finite-thermal-window}

Fix a bounded interval $I$ of length $L$ and $T>0$. The geometric KMS
construction is locally a vacuum pullback by the exponential map
\cite[Section~2.8, Proposition~2.4 and Theorem~2.5]{Camassa2012KMS}.
Its compact completion permits a finite global preparation. Use the
state convention
\[
 C_\rho(x,y)=\frac{i\sqrt{\rho'(x)\rho'(y)}}
                   {2\pi[\rho(x)-\rho(y)-i0]},\qquad
 \langle T(x)\rangle_\rho=-\frac{c}{24\pi}\{\rho,x\}.
\]
Here $c=1$ for each complex-fermion stream, and the inverse
diffeomorphism implementer acts on the vacuum.

\begin{proposition}[A thermal restriction with finite global energy]
\label{prop:finite-thermal-window}
There is a smooth increasing $\rho:\mathbb R\to\mathbb R$, equal to
the identity outside a compact interval, whose vacuum-orbit state has
the complete zero-charge thermal restriction on $I$. Put
$h=\log\rho'$ and $g(y)=\rho'(\rho^{-1}(y))$. Its energy and variance are
\begin{align}
 E_\rho&=\frac{c}{48\pi}\int h'(x)^2\,dx,\nonumber\\
 \operatorname{Var}_\rho(H_f)&=
  \frac{c}{48\pi^2}\int_0^\infty
     \omega^3|\widehat{g-1}(\omega)|^2\,d\omega<\infty,
                 \label{eq:thermal-window-moments}
\end{align}
where $\widehat k(\omega)=\int e^{-i\omega y}k(y)dy$.
The global state is pure, not globally KMS. Its outside correlations
purify the thermal restriction.
\end{proposition}
\begin{proof}
Set $\alpha=2\pi T$ and prescribe $\rho=e^{\alpha x}$ on a
neighborhood of $I$. Extend $\log\rho'$ smoothly with compact support.
A disjoint bump can adjust $\int(e^h-1)$ to zero: a positive bump
has unbounded positive integral, while a negative bump on sufficiently
large support has negative range exceeding any fixed surplus.
Then $\rho(x)=x+\int_{-\infty}^x(e^h-1)$ is a compact completion,
up to an irrelevant translation on $I$.
Its covariance on $I$ is $iT/[2\sinh\pi T(x-y-i0)]$.
Wick's rule identifies all local fermion correlations; the cited
geometric construction gives the corresponding observable-net restriction.
Integrating the Schwarzian gives the first moment in
Eq.~\eqref{eq:thermal-window-moments}. Conjugating $H_f$ into the
vacuum gives $T(g)+E_\rho$. Since $H_f\Omega=0$, the variance is the
vacuum second moment of $T(g-1)$, with the normalization of
\cite[Eq.~(4.11)]{FewsterHollands2019}.
The test $g-1$ is smooth and compact, so that moment is finite.
\end{proof}

\begin{proposition}[Preparation cost in the vacuum orbit]
\label{prop:thermal-window-cost}
Every compact-diffeomorphism vacuum preparation with this thermal
restriction obeys
\begin{equation}\label{eq:thermal-window-cost}
 E_\rho\ge\frac{\pi cT^2L}{12}
       -\frac{cT}{6}\tanh\frac{\pi TL}{2}.
\end{equation}
The right side is the infimum in this class when outside support is
unrestricted, and is not attained for $T,L>0$.
\end{proposition}
\begin{proof}
Thermal stress gives $\{\rho,x\}=-(2\pi T)^2/2$ on $I$.
Thus $w=h'$ satisfies $w'=(w^2-\alpha^2)/2$.
The branches $w=\pm\alpha$ and the nonsingular coth branches have
$w^2\ge\alpha^2$. The remaining branches are
$-\alpha\tanh[\alpha(x-s)/2]$; their derivative-square integral on
$I$ is minimized when $s$ is its midpoint. That integral is
$\alpha^2L-4\alpha\tanh(\alpha L/4)$.
For its compact completions, flatten the endpoint derivatives in
collars of width $\epsilon$, at cost $O(\epsilon)$, and join the
endpoint $h$-values to zero over length $D$, at cost $O(D^{-1})$.
A disjoint bump of scale $D$ corrects the displacement integral with
bounded amplitude and cost $O(D^{-1})$. Taking $\epsilon\to0$ and
then $D\to\infty$ reaches the infimum. Zero exterior cost would
force $h'=0$ outside $I$, incompatible with its endpoint derivatives.
\end{proof}

This is a preparation-class bound, not a minimum energy for all
states or complete devices. Its asymptotic slope is the thermal
current. Local negative stress outside $I$ is included; counting only
the positive plateau gives a different quantity. For $L=8$ and
$T=1/(2\pi),.31894305$, explicit smooth completions cost
$.11257242,.60868371$, above the respective infima
$.02748002,.15992956$. Increasing support or sharpening a boundary
can approach the infimum without keeping the drive or rms resources
uniformly bounded.

For a regular increasing clock $\eta(v)$ on a finite interval, complete
$\rho(v)=\exp[2\pi T\eta(v)]$ instead. The entire local covariance
is the thermal $\eta$ covariance pulled back by that clock, and
\begin{equation}\label{eq:prepared-clock-stress}
 \langle T(v)\rangle_\rho
 =\eta'(v)^2\frac{\pi cT^2}{12}
       -\frac{c}{24\pi}\{\eta,v\}.
\end{equation}
Two independently prepared copies therefore supply the two-port
input correlations on the operating interval, including the output
cross covariance and switching stress. Their global purity and
inside--outside correlations remain in the ledger. Thermal local
correlations alone do not establish a finite-window erasure bound.
The cost in Eq.~\eqref{eq:thermal-window-moments} uses the actual
physical coordinate; Eq.~\eqref{eq:thermal-window-cost} cannot be
applied to a rapidity interval without its clock conversion.
The local diffeomorphism has a driven stress-unitary preparation
path. Its autonomous bounded approximation is developed in
Appendix~\ref{sec:local-thermal-preparation}.

\subsection{A record-selected source is not its average covariance}
\label{sec:programmed-source}

Suppose a retained classical program $R$ has probabilities $p_r$, is
jointly independent of the two prepared incoming streams, and selects
$\theta_r$. Conditional on $r$, the preceding contact law holds.
For a finite frequency regulator with independent input mode pairs,
write $\rho_{r,N}=\bigotimes_{j=1}^N\rho_{r,j}$ for the reduced source
and $\bar\rho_N=\sum_rp_r\rho_{r,N}$ for its unconditional state.
The return stream and program are retained in the joint process.

\begin{proposition}[The source carries the program's relations]
\label{prop:programmed-source}
For distinct frequencies $j,k$, with $\Delta n_j=n_o(\omega_j)-n_H(\omega_j)$,
\begin{equation}\label{eq:programmed-source-covariance}
 \operatorname{Cov}(n_j,n_k)=
 \operatorname{Var}_p(\chi_r)\,\Delta n_j\Delta n_k.
\end{equation}
Furthermore,
\begin{equation}\label{eq:programmed-source-entropy}
 S(\bar\rho_N)=\sum_rp_r S(\rho_{r,N})+\mathcal I_N,
 \qquad 0\le\mathcal I_N\le H(R).
\end{equation}
Here $\mathcal I_N=I(R:\text{source})$ in the classical--quantum
record state. The product state with the same one-mode occupations
has entropy $\sum_jh(\sum_rp_r n_{r,j})$, generally larger.
If $H(R)=o(\mathcal T)$ on a growing observation window and the
conditional stationary source entropy rates exist, the unconditional
excess rate is $\bar s=\sum_rp_r s_{\chi_r}$. Its mean current obeys
\[
 \bar q\ge T_H\bar s+\frac{3\bar s^2}{\pi}.
\]
This limit assumes the same program is retained across the window,
not a fresh independently randomized choice at each mode or time.
\end{proposition}
\begin{proof}
Conditional occupations are independent and have means
$n_{r,j}=n_H(\omega_j)+\chi_r\Delta n_j$. Conditional covariance
therefore vanishes; covariance of their means gives
Eq.~\eqref{eq:programmed-source-covariance}.
The entropy identity is the classical--quantum mutual information;
its upper bound is the classical program entropy. Subadditivity
gives the product-marginal upper bound. Divide the identity by the
window length. Averaging Eq.~\eqref{eq:contact-entropy-floor} and
using $\sum_rp_rs_r^2\ge\bar s^2$ proves the current bound.
\end{proof}

A single persistent program bit can thus create correlations across
an arbitrarily long source without contributing an extensive new
entropy of its own. Replacing that state by an independent Gaussian
state at its averaged covariance destroys those relations. If a new
program is acquired at every interval, its entire retained sequence
belongs in $R$; $H(R)$ can then be extensive. Angle switches have
the separate two-port work of Eq.~\eqref{eq:contact-pump}, and neither
entropy identity makes those switches free. These distinctions apply
to a gravitational source selected by internal records: mean stress
alone does not specify the source state.

There is an exact finite autonomous realization of the static
record-controlled exchange. For matched frequencies let
$K_j=c_{H,j}^\dagger c_{o,j}+c_{o,j}^\dagger c_{H,j}$ and use
\[
 V=\sum_r|r\rangle\langle r|\otimes
        \prod_j e^{-i\theta_r K_j},\qquad
 H_0=\sum_j\omega_j(n_{H,j}+n_{o,j}).
\]
It commutes with $H_0$. Put $V$ on one edge of the positive history
clock and identities on the remaining edges. The full Hamiltonian
$H_0+H_{\rm hist}$ is nonnegative, time independent and conserves
$H_0$. Its endpoint evolution executes the exchange uniformly over
all retained references. This is a finite-mode control realization;
its clock initialization and finite thermal preparations are inputs.
It does not implement the moving conformal contact's material stress
or remove its continuum switching work.

\begin{proposition}[Two reservoirs and the complete controller ledger]
\label{prop:programmed-ledger}
Let finite reservoirs start at Gibbs states $\tau_H,\tau_o$ with
inverse temperatures $\beta_H>\beta_o$, jointly independent of a
complete zero-bare-gap controller $M$. Its internal correlations and
retained references are included in $M$, but not purifiers of the
reservoirs: those would violate the stated initial independence.
Under any joint unitary
conserving reservoir bare energy, put $Q=\Delta E_H=-\Delta E_o$,
$D_i=D(\rho_i'\Vert\tau_i)$ and
$I'=S(\rho_H')+S(\rho_o')+S(\rho_M')-S(\rho_{HoM}')$.
Then
\begin{equation}\label{eq:programmed-ledger}
 (\beta_H-\beta_o)Q=D_H+D_o+I'-\Delta S_M.
\end{equation}
If the complete controller entropy is unchanged, the right side is
$D_H+D_o+I'\ge0$. Returning its marginal does not mean eliminating
its output correlations.
\end{proposition}
\begin{proof}
The Gibbs identity gives $D_i=\beta_i\Delta E_i-\Delta S_i$.
Initial independence and unitary invariance of total entropy give
$\Delta S_H+\Delta S_o+\Delta S_M=I'$.
Add the Gibbs identities and use energy conservation.
This is the multi-reservoir form of the relative-entropy ledger
underlying the sharpened Landauer principle~\cite[Theorem~3]{ReebWolf2014}.
\end{proof}

At frequencies $1,2$, temperatures $.7,1.4$, program probability $.37$
and angles $.2,1.1$, direct joint evolution transfers $.13174972$
energy into the colder source. Its cross-frequency covariance is
$.00249817$. Source entropy is $.86881190$, conditional entropy
$.84854079$, and the independent state at averaged occupations has
entropy $.86900230$. A five-state clock has mean energy $.12$ and
interaction cutoff $.24$; its endpoint matrix and the thermal ledger
agree with direct evolution. Both reservoirs have total bare cutoff
$6$ in these units. None of these cutoffs describes an infinite KMS field.

An operational error also needs an energy qualification before it
certifies a source. For $0\le H_E\le E_{\rm cut}I$,
$|\operatorname{tr}H_E(\rho-\sigma)|\le E_{\rm cut}D(\rho,\sigma)$.
Without a common energy constraint, states at trace distance
$\epsilon$ can differ by fixed or unbounded energy: mix the vacuum
with probability $\epsilon$ in occupation $N=1/\epsilon$.
A joint clock or gate-distance certificate alone therefore does not
certify stress of an unbounded field. The continuum source and its
support require energy-weighted control as well as retained records.

\section{A physical contact: reaction, source and recoil}
\label{sec:inertial-contact}

A prescribed scattering angle charges its work but leaves its actuator
outside the model. Give the angle $\theta$ a momentum and inertia $I>0$,
and couple it to a second coordinate $y$ with inertia $J>0$ and potential
\begin{equation}\label{eq:inertial-potential}
 V(\theta,y)=\frac{k}{2}(\theta-y)^2+\frac{K}{2}y^2,
 \qquad k,K>0.
\end{equation}
The coordinates are noncompact; the potential pins them. In the
$\sigma_x$ basis the two-port contact has charges $q_1=1,q_2=-1$.
For a smooth, nonnegative profile $f_a$ of integral one, use the
current Hamiltonian
\begin{equation}\label{eq:inertial-H}
 H=\frac{p_\theta^2}{2I}+\frac{p_y^2}{2J}+V(\theta,y)
   +\pi\sum_{i=1}^2\int
       :\!\left(j_i(x)+\frac{q_i\theta f_a(x)}{2\pi}\right)^2\!:\,dx.
\end{equation}
The currents obey $[j_i(x),j_l(x')]=-i\delta_{il}\delta'(x-x')/(2\pi)$.
At oscillator cutoffs, normal ordering uses the positive completed-square
Hamiltonian. Its $\theta^2$ counterterm must not be dropped. The
continuum identities below apply on a domain where the current and
stress products and their regulated limits exist. The positive form
constructions below supply both a neutral-current Hamiltonian and
the contact on the full fermion carrier. A static
current displacement can still have infinite particle number; its
dressing and asymptotic state need separate representation control.
An oscillator cutoff is not a certificate
for the full two-temperature fermion channel.
The anomaly-induced reaction is the mechanism of the fermion--rotor
model~\cite[Sections~2.1--2.2]{LoladzeOkuiTong2026}. Here it is applied
to the pinned two-coordinate actuator and its stored energy.

\begin{proposition}[Current reaction with finite contact width]
\label{prop:inertial-reaction}
Set $J_i=j_i+q_i\theta f_a/(2\pi)$. Hamiltonian~\eqref{eq:inertial-H}
gives
\begin{align}\label{eq:inertial-current}
 (\partial_t+\partial_x)J_i&=\frac{q_i}{2\pi}\dot\theta f_a,
 &I\ddot\theta+V_\theta&=-\int f_a\sum_i q_iJ_i\,dx,
 &J\ddot y+V_y&=0.
\end{align}
With the free initial currents defined at $t=0$, elimination of the
field gives
\begin{align}\label{eq:inertial-memory}
 I\ddot\theta+V_\theta+\int_0^t\Gamma_a(t-s)\dot\theta(s)\,ds
   &=-\xi_a(t)-\Gamma_a(t)\theta(0),\\
 \Gamma_a(s)&=\frac1\pi\int f_a(x)f_a(x-s)\,dx,\nonumber\\
 \xi_a(t)&=\int f_a(x)\sum_iq_i j_i(0,x-t)\,dx.\nonumber
\end{align}
The integrated retarded coefficient is $\gamma=1/(2\pi)$.
For a Gaussian profile of standard deviation $a$,
\begin{equation}\label{eq:inertial-Gaussian-kernel}
 \Gamma_a(s)=\frac{\gamma}{\sqrt\pi a}e^{-s^2/(4a^2)},
 \qquad \operatorname{Re}\widetilde\Gamma_a(\Omega)
       =\gamma e^{-a^2\Omega^2}.
\end{equation}
\end{proposition}
\begin{proof}
The current commutator gives the first equation in
\eqref{eq:inertial-current}; canonical commutators give the other two.
Solve the transport equation along its characteristics. The initial
$q_i\theta(0)f_a/(2\pi)$ term gives the initial slip in
\eqref{eq:inertial-memory}, and the source integral gives its convolution.
The autocorrelation of $f_a$ is even and has integral one on the whole
line, hence $\int_0^\infty\Gamma_a(s)\,ds=1/(2\pi)$.
Fourier transformation gives \eqref{eq:inertial-Gaussian-kernel}.
\end{proof}

Thus damping is derived from the current algebra, not appended to an
externally specified pump. For a velocity history extended into the
past with bounded derivative,
\begin{equation}\label{eq:inertial-memory-bound}
 |(\Gamma_a*v)(t)-\gamma v(t)|
 \le\|v'\|_\infty\int_0^\infty s\Gamma_a(s)\,ds.
\end{equation}
The last integral is $2\gamma a/\sqrt\pi$ for the Gaussian. If $f_a$
is supported in $[-a,a]$, it is at most $2\gamma a$. This controls
smooth histories, not the ultraviolet error of an entire quantum
channel. The initial slip is also physical: setting the mean angle
to zero does not set its operator fluctuations to zero.

For independent zero-charge thermal inputs with Fermi occupations
$n_H,n_o$, Wick contraction in the original port basis gives the
unsymmetrized incoming force spectrum
\begin{equation}\label{eq:inertial-noise}
 S_\xi(\Omega)=\frac{|\widehat f_a(\Omega)|^2}{2\pi}
 \int_{\mathbb R}\left[n_H(E)(1-n_o(E+\Omega))
                  +n_o(E)(1-n_H(E+\Omega))\right]dE.
\end{equation}
Changing variables in $S_\xi(-\Omega)$ yields
\begin{equation}\label{eq:inertial-noise-commutator}
 S_\xi(\Omega)-S_\xi(-\Omega)
       =2\gamma\Omega|\widehat f_a(\Omega)|^2.
\end{equation}
Indeed each occupation integrates
$n(E)-n(E+\Omega)$ to $\Omega$ over the filled sea.
The retarded response is temperature independent; its noise is not.
Only at equal temperature $T$ does the spectrum reduce to
$\Omega|\widehat f_a|^2/[\pi(1-e^{-\Omega/T})]$ and obey the single
KMS fluctuation--dissipation relation. Different-temperature inputs
are generally correlated in the charge basis and cannot be replaced
by two independent thermal boson baths.

In a point-contact limit with well-defined renormalized products and
the symmetric contact prescription, the total field power and
mechanical energy satisfy
\begin{equation}\label{eq:inertial-power}
 P_{\rm field}=\tfrac12\{\dot\theta,\xi\}+\gamma\dot\theta^2,
 \qquad \dot E_{\rm mech}=-P_{\rm field}.
\end{equation}
This follows by substituting the incoming--outgoing current jump into
the Sugawara stress and using the canonical mechanical equation.
At finite width the full Hamiltonian supplies the energy identity,
including retained current energy. The positive prescribed-angle pump
term alone is not the power of a quantum actuator. Its velocity
correlates with the incoming current. In a stationary interacting
Gibbs state $\langle\dot\theta^2\rangle$ can be positive while the
net energy transfer is zero.

\begin{proposition}[A mechanical record changes the actual field source]
\label{prop:inertial-record-source}
Take a common initial state with zero first moments of the two
mechanical coordinates, their momenta and the currents, and finite
stress expectations as above. In branch $R$, displace $y$ by $A_R$
and leave the other first moments unchanged. Relative to the
undisplaced evolution under the same Hamiltonian, let $v_R(t)$ be
the deterministic mean angle velocity and define
\begin{equation}\label{eq:inertial-record-stress}
 h_R(t,x)=\int_0^t v_R(s)f_a(x-t+s)\,ds,
 \qquad \Delta\langle T_{\rm field}(t,x)\rangle
       =\gamma h_R(t,x)^2.
\end{equation}
Then
\begin{equation}\label{eq:inertial-record-energy}
 \Delta E_{\rm mech}(t)+\gamma\int h_R(t,x)^2\,dx
             =\frac{k+K}{2}A_R^2.
\end{equation}
The actual excess stress scales as $A_R^2$ for fixed hardware.
\end{proposition}
\begin{proof}
The Heisenberg equations are linear. Displacing the initial $y$ adds
the classical solution to the undisplaced operator solution; it leaves
every centered correlation unchanged. Characteristics give
$\Delta\langle J_i\rangle=q_i h_R/(2\pi)$. The zero first moments
remove the cross term in the stress difference, giving
\eqref{eq:inertial-record-stress}. Differentiating its spatial integral
gives $v_R(\Gamma_a*v_R)$, the negative derivative of the mechanical
displacement energy. The initial displacement cost is
$(k+K)A_R^2/2$.
\end{proof}

This cancellation compares displaced and undisplaced evolutions; it
does not discard the quantum noise term in \eqref{eq:inertial-power}.
The common baseline stress still has to be counted. Energy retained
inside the contact profile is not yet energy that has reached an
external outlet. An acquired record can select $A_R$ through a paid
preparation. That preparation, its storage and its physical copying
are additional operations, not free changes of a target potential.

For $I=10,J=20,k=1,K=.1,a=.2$, a 192-mode positive oscillator
regulator initially displaced by $A_R=1.3$ has excess energy $.9295$.
At $t=80$, $.5103950$ is in its shifted field coordinates and
$.4191050$ remains mechanical. Direct time-domain integration of the
continuum memory equation and separate spatial integration give
$.5103953$ and $.4191052$, with total energy residual
$4.8\times10^{-7}$ at step $.005$. This is a check of the continuum
energy identity, not a cutoff channel certificate. The Gaussian profile
has spatial tails. In the regulator's displaced interacting Gibbs state
at $T=.02$, $t=8$, the two branches $A_R=.5,1.3$ give
$\langle\sin^2\theta\rangle=.426644,.716636$. These are coordinate
observables, not transmission probabilities for the two-temperature
fermion problem. A finite regulated bath eventually recurs.

\subsection{Installing a program on the running carrier}

\begin{proposition}[Autonomous installation on an always-running carrier]
\label{prop:inertial-installation}
Let $0\le H_{\rm phys}\le E_{\rm cut}I$ be a finite actuator--field
regulator, and let a prefix of $r$ gates prepare its mechanical state
from an already acquired record. Use a positive history clock with
$m\ge r$, strength $J$, transfer endpoint $T=\pi/(2J)$ and identity
edges after the prefix. The common, always-on Hamiltonian
\begin{equation}\label{eq:inertial-installation-H}
 H=H_{\rm hist}+I_{\rm clock}\otimes H_{\rm phys}\ge0
\end{equation}
has, for $t_0\le t\le T$, joint distance from the completed-prefix
state freely continued under $H_{\rm phys}$ for $t-t_0$ at most
\begin{equation}\label{eq:inertial-installation-bound}
 d_{\rm joint}\le\min\{1,E_{\rm cut}t_0+2\sqrt{P_0}
                     +4A_r(t-t_0)\sqrt{P_0}\},
\end{equation}
where $P_0=\Pr\{\operatorname{Bin}(m,\sin^2Jt_0)\le r\}$ and
$A_r=J\sum_{j=1}^r\sqrt{j(m-j+1)}$. All prior references are retained.
For each finite continuation window and specified error, finite
sufficiently large clock resources attain that error.
\end{proposition}
\begin{proof}
During the prefix, Duhamel comparison with $H_{\rm hist}$ gives
$E_{\rm cut}t_0$. Completing the prefix on its clock tail costs
$2\sqrt{P_0}$. Thereafter compare with the identity-edge clock and
the same $H_{\rm phys}$. Their reference state factors across clock
and data, so the nonidentity edges have norm at most
$4A_r\sqrt{P_0}$; Duhamel gives the last term. The binomial lower
tail decreases up to $T$. Set $\sin^2Jt_0=\mu/m$ with $\mu$ a
sufficiently large multiple of $\log m$. Then $t_0\to0$ and the tail
beats the polynomial growth of $A_r$. Choose $T$ to contain the
specified continuation window.
\end{proof}

For a bounded physical source observable $0\le B\le B_{\rm cut}I$,
the expectation error is at most $B_{\rm cut}d_{\rm joint}$. Thus a
finite-cutoff record can install a mechanical program while the same
carrier and field continue to evolve; no external switch at $t_0$ is
required. A three-level actuator and three-level field, with the old
record reference retained, give joint distance $.00204248$ below the
bound $.04082928$ in one small-clock check. The clock mean energy is
$960$, so that check is not an installation inside the mass-10
gravitational body. Larger clocks improve the guarantee at larger
energy and dimension. The continuum displacement, infrared sector,
spatial transport and matching to a moving action need separate control.

The upper energy cutoff is unnecessary for a class of permitted gates.
The relevant condition concerns how a gate changes energy, rather than
the spectrum of the carrier itself.

\begin{proposition}[Energy-regular return on an unbounded carrier]
\label{prop:energy-regular-return}
Let $H_p\ge0$ be self-adjoint, and let the finite prefix unitaries
$W_j=V_j\cdots V_1$, $W_0=I$, preserve $\operatorname{Dom}H_p$.
Suppose
\[
 B_j=W_j^*H_pW_j-H_p,\qquad
 b=\max_{j\le r}\|B_j\|,\qquad
 c=\max_{j\le r}\|[H_p,B_j]\|
\]
extend to bounded operators, with $B_j$ preserving
$\operatorname{Dom}H_p$. Let $\psi$ be the normalized full input
purification, including prior references, with $E=\|H_p\psi\|<\infty$.
Use the same positive history clock as in
Proposition~\ref{prop:inertial-installation}. For $t_0\le t\le T$, the
joint vectors under $H_{\rm hist}+I\otimes H_p$ and the reference
completed-prefix evolution differ by at most
\begin{equation}\label{eq:energy-regular-vector}
 \varepsilon_m(t)=(E+b)t_0+2\sqrt{P_0}
                  +4A_r(t-t_0)\sqrt{P_0}.
\end{equation}
Their physical rms energies obey
\begin{equation}\label{eq:energy-regular-moments}
 \|H_p\Psi_a(t)\|\le M_a=E+tc+b,\qquad
 \|H_p\Psi_i(t)\|\le M_i=E+b.
\end{equation}
These bounds do not grow with the padded clock. If a self-adjoint
source observable $Q$ has $\operatorname{Dom}H_p\subseteq\operatorname{Dom}Q$
and $\|Q\zeta\|\le a\|H_p\zeta\|+d\|\zeta\|$, then
\begin{equation}\label{eq:energy-regular-source}
 |\langle Q\rangle_a-\langle Q\rangle_i|
 \le\varepsilon_m(t)\,[a(M_a+M_i)+2d].
\end{equation}
In particular, this certifies the unbounded physical energy $Q=H_p$.
For any fixed finite continuation window, $\varepsilon_m\to0$ with
sufficient clock padding.
\end{proposition}
\begin{proof}
The ideal history branches are orthogonal in the clock, and
$\|H_pW_j\psi\|=\|(H_p+B_j)\psi\|\le E+b$. The early Duhamel
comparison therefore costs $(E+b)t_0$; the prefix tail and late
comparison are as in Proposition~\ref{prop:inertial-installation}.
This gives a vector bound on the retained purification, not only a
distance between reduced states.

Put $G=\sum_j|j\rangle\langle j|\otimes W_{\min(j,r)}$ and write the
ungated clock as $K_m$. In the history gauge the total Hamiltonian is
\[
 G^*H_{\rm tot}G=K_m\otimes I+I\otimes H_p+
          \sum_j|j\rangle\langle j|\otimes B_{\min(j,r)}.
\]
The last term and its commutator with $H_p$ are bounded by $b,c$.
The propagator commutator consequently has norm at most $tc$, so
the gauge-state rms energy is at most $E+tc$. Transforming back adds
at most $b$. The same domain-preserving commutator argument applies
to the reference. Finally write the expectation difference as
$\langle\Psi_a-\Psi_i,Q\Psi_a\rangle+
\langle Q\Psi_i,\Psi_a-\Psi_i\rangle$ and apply Cauchy--Schwarz.
The large-clock limit follows from the same binomial tail choice.
\end{proof}

A positive quadratic-form source needs a weaker estimate. If
$0\le Q\le aH_p+dI$ as forms, Cauchy--Schwarz in the form norm gives
\begin{equation}\label{eq:energy-regular-form}
 |\Delta\langle Q\rangle|\le
 (\sqrt{aM_a+d}+\sqrt{aM_i+d})
 \sqrt{a\varepsilon_m(M_a+M_i)+d\varepsilon_m^2}.
\end{equation}
Here $\|H_p^{1/2}(\Psi_a-\Psi_i)\|^2
\le\|\Psi_a-\Psi_i\|\,\|H_p(\Psi_a-\Psi_i)\|$ proves the bound.
A normally ordered local stress is not automatically a positive form
or $H_p$-bounded; its relative estimate must be supplied.

\begin{proposition}[A smooth chiral stress certificate]
\label{prop:smooth-stress-certificate}
Consider a positive-energy chiral CFT satisfying the
Fewster--Hollands hypotheses, with $H_f=\int T(x)\,dx$ and central
charge $c_{\rm ch}$. Let $0\le f\in C_c^\infty$ and
$\Lambda_f>\|f\|_\infty$. Define
\begin{align}
 q_f&=\frac{c_{\rm ch}}{12\pi}\int|(\sqrt f)'|^2\,dx,\nonumber\\
 q_{\rm comp}&=\frac{c_{\rm ch}}{48\pi}
       \int\frac{f'(x)^2}{\Lambda_f-f(x)}\,dx,\qquad
 d_f=q_f+q_{\rm comp}. \label{eq:stress-shift}
\end{align}
Then, as quadratic forms,
\begin{equation}\label{eq:stress-energy-form}
 0\le T(f)+q_f I\le\Lambda_f H_f+d_f I.
\end{equation}
If $H_f\le a_0H_p+b_0I$ as forms for the physical Hamiltonian of
Proposition~\ref{prop:energy-regular-return}, its expectation
certificate \eqref{eq:energy-regular-form} applies to the actual
local stress $T(f)$ with $a=\Lambda_f a_0$ and
$d=d_f+\Lambda_f b_0$. The constant shift cancels between the
two states.
\end{proposition}
\begin{proof}
The chiral quantum energy inequality gives $T(f)\ge-q_fI$
\cite[Theorem 4.1]{FewsterHollands2005}. To obtain the upper bound,
one cannot apply that theorem directly to the non-Schwartz function
$\Lambda_f-f$. Choose $k\in C_c^\infty$, $0\le k\le1$, equal to
one on $[-1,1]$, and set $\chi_L(x)=k(x/L)^2$. For $L$ enclosing
$\operatorname{supp}f$, the admissible test
$g_L=\chi_L(\Lambda_f-f)$ has QEI constant
\[
 q(g_L)=q_{\rm comp}+
       \frac{c_{\rm ch}\Lambda_f}{12\pi L}\int|k'|^2.
\]
Its two derivative supports are disjoint. On the common CFT core,
$T(\chi_L)$ tends in expectation to $H_f$, using the integrable
stress and the translation charge. Thus
$T(g_L)=\Lambda_fT(\chi_L)-T(f)\ge-q(g_L)I$ gives the upper bound.
The positive form then extends continuously to
$\operatorname{Dom}H_f^{1/2}$. Combining it with the stated
physical form inequality and \eqref{eq:energy-regular-form}
proves the certificate.
\end{proof}

For the inertial contact there is a useful conditional transfer of
this field bound to the interacting carrier. In its positive
completed-square regulator, write the shifted field energy as
$\sum_l\omega_l(a_l+s_l\theta)^*(a_l+s_l\theta)$ and put
$d_a=\sum_l\omega_l|s_l|^2
=k_{\rm cur}\int f_a^2/(4\pi)$, where $k_{\rm cur}=2$.
Since $V\ge kK\theta^2/[2(k+K)]$, the elementary square inequality gives
\begin{equation}\label{eq:contact-field-domination}
 H_f\le 2H_{\rm shifted}+2d_a\theta^2
 \le\left[2+\frac{4d_a(k+K)}{kK}\right]H_p.
\end{equation}
In the neutral current Fock representation, that positive Hamiltonian
can be constructed without a static dressing unitary. Set
$g_i(\omega)=q_i\sqrt{\omega}\,\widehat f_a(\omega)/(2\pi)$, so
$\sum_i\int|g_i|^2/\omega=d_a$. On the form domain of
$H_0=H_{\rm mech}+H_f$, define
\begin{equation}\label{eq:current-positive-form}
 \mathfrak q_H[\psi]=\mathfrak q_{\rm mech}[\psi]+
 \sum_i\int_0^\infty
 \left\|\sqrt{\omega}\,a_i(\omega)\psi+
          \frac{g_i(\omega)}{\sqrt{\omega}}\theta\psi\right\|^2d\omega.
\end{equation}
This includes the counterterm. Put
$A_a=2+4d_a(k+K)/(kK)$. The same square inequalities give
$\mathfrak q_0\le(1+A_a)\mathfrak q_H$ and
$\mathfrak q_H\le A_a\mathfrak q_0$. Thus the two form norms,
with $\|\psi\|^2$ added, are equivalent. The closed positive form
defines a self-adjoint Hamiltonian, and
\eqref{eq:contact-field-domination} holds for it.
Coupling cutoffs converging in $L^2(d\omega/\omega)$ converge in
the common form norm and hence in norm resolvent on this field space.
This does not require $g_i/\omega\in L^2$: when
$\widehat f_a(0)\ne0$, the static dressing can fail to be a Fock
unitary even though the Hamiltonian exists.

For a compact contact profile,
a stress test outside its support reads the physical outgoing
current stress itself. A Gaussian profile has tails and is not
covered by that exact exterior assertion.
The gates must be energy-regular for this interacting $H_p$;
free-field regularity alone does not establish that condition when
the gate overlaps the contact. The exterior case admits an exact transfer
below.
This neutral-current construction does not establish a dressed ground
state or asymptotic charged representation. The full fermion carrier
and its prepared thermal inputs can instead be treated directly.
A uniform zero-width limit does not follow from the form inequality.

\paragraph{The contact on the full fermion carrier.}
The two complex fermions need not be replaced by two neutral current
Fock spaces. Work in their full Dirac-sea CAR representation, with
positive translation generator $H_f$. Smooth local gauge implementers
and the current anomaly act on this whole space, including its charged
sectors~\cite[Sections~3.2--3.3 and Appendices~A--B]{CareyLangmann1998}.
In the $\sigma_x$ basis write
$J(f)=\sum_iq_i j_i(f)$, $\kappa_c=\sum_iq_i^2=2$ and
$d_f=\kappa_c\|f\|_2^2/(4\pi)$.

\begin{proposition}[A positive contact on the full CAR carrier]
\label{prop:full-fermion-contact}
For real $f\in C_c^\infty(\mathbb R)$ and the pinned mechanical
potential~\eqref{eq:inertial-potential}, the form
\[
 \mathfrak q_p=\mathfrak q_{\rm mech}+\mathfrak q_f
       +\langle\theta J(f)\rangle+d_f\langle\theta^2\rangle
\]
is closed and nonnegative on the form domain of
$H_0=H_{\rm mech}+H_f$. It defines the single-contact self-adjoint
$H_p$ on the full fermion space. Put
$a_0=kK/[2(k+K)]$ and $L=\max\{2,1+2d_f/a_0\}$.
Then $L^{-1}\mathfrak q_0\le\mathfrak q_p\le L\mathfrak q_0$.
The assertion does not require a static dressing unitary.

For $b\in C_c^\infty((1,2))$ of integral one, set
$b_D(x)=D^{-1}b(x/D)$ and
$f_D=f-(\int f)b_D$, with disjoint supports. The corresponding
compact-gauge contacts converge in norm resolvent to $H_p$.
If $d_b=(\int f)^2\kappa_c\|b\|_2^2/(4\pi D)$ and
$r_D=\sqrt{d_b/a_0}+d_b/a_0$, then, for $Lr_D<1$,
\[
 \|(H_{p,D}+1)^{-1}-(H_p+1)^{-1}\|
       \le\frac{Lr_D}{1-Lr_D}.
\]
\end{proposition}
\begin{proof}
First let $f=F'$ with $F\in C_c^\infty$.
Multiplication by $e^{\pm i\lambda F}$ has a Hilbert--Schmidt
commutator with the one-particle Hardy projection: its squared norm
is $\int\!\int|g(x)-g(y)|^2/(x-y)^2\,dxdy/(4\pi^2)$,
finite by smoothness and compact support of $g-1$.
Thus $W_\lambda=e^{-i\lambda J(F)}$ is implemented on the full CAR
space. On the common smooth core,
$[J(F),H_f]=iJ(F')$ and
$[J(F),J(F')]=i\kappa_c\int F'^2/(2\pi)$ give
\[
 W_\lambda H_f W_\lambda^*
       =H_f+\lambda J(f)+d_f\lambda^2\ge0.
\]
For general $f$, use $f_D$ above. On a finite smooth particle--hole
core, $J(b_D)\to0$ strongly: its vacuum norm is $O(D^{-1})$,
and commuting it past each fermion changes the test function by
multiplication with $b_D$, of supremum norm $O(D^{-1})$.
Meanwhile $d_{f_D}\to d_f$. Passing to the limit and minimizing
over $\lambda$ gives the form inequality
\[
 |\langle\psi,J(f)\psi\rangle|
       \le2\sqrt{d_f}\,\|\psi\|\,\|H_f^{1/2}\psi\|.
\]
Polarization extends the form to $\operatorname{Dom}H_f^{1/2}$.
For each mechanical coordinate the shifted field form
$A_f(\theta)=H_f+\theta J(f)+d_f\theta^2$ is nonnegative and
obeys $A_f\le2H_f+2d_f\theta^2$ and
$H_f\le2A_f+2d_f\theta^2$.
Integrate over the mechanical coordinates and use $V\ge a_0\theta^2$.
The displayed equivalent form norms establish closedness and
positivity.

For the compensator, the same current bound and Cauchy--Schwarz give
\[
 |\langle\theta J((\int f)b_D)\rangle|
       \le\sqrt{d_b/a_0}\,\mathfrak q_0.
\]
The counterterm difference is $d_b\langle\theta^2\rangle$.
Hence
\[
 |\mathfrak q_{p,D}-\mathfrak q_p|\le r_D\mathfrak q_0
       \le Lr_D\mathfrak q_p.
\]
The resolvent identity in the
$(H_p+1)^{1/2}$ form scale gives the stated estimate.
\end{proof}

The finite-rms preparations of
Proposition~\ref{prop:finite-thermal-window} now belong to the same
fermion carrier as the inertial contact. Their outside purifications
remain present. A constant flavor rotation preserves the vacuum but
generally gives off-diagonal covariance for unequal temperatures;
the rotated currents are not two independent KMS baths. The force
law therefore uses the original quasifree input correlations, not
a substitute product of bosonic thermal states. Finite form energy
does not by itself supply the graph regularity required by the
autonomous return theorem. On the stated smooth input core, exterior
even local gates have the same free and interacting commutators by
CAR locality, and the local-source certificate uses $H_f\le LH_p$.

The compensator is a proof and approximation device, not a free heat
outlet. For $\int f=1$, $\operatorname{supp}f\subset[-a,a]$ and
$0\le s\le T<D-a$, its retarded kernel differs from the single
contact's by at most $2d_b$. For every fixed $D$, however,
$\int f_D=0$ makes the integrated kernel vanish. Its dc response is
zero, whereas the single contact's is $1/(2\pi)$.
Norm-resolvent convergence is consequently not a uniform
infinite-time heat-transport statement. The single-contact
Hamiltonian itself only couples in $\operatorname{supp}f$; it does
not require a physical compensator extending to infinity.

These estimates certify smooth, finite-window source charges.
To use a sharp or unbounded gravitational kernel, one must prove
the relevant source support or tail control. An accurate field
gate now has a route to an accurate local source; the source
state, its correlations and its gravitational realization still
have their own inputs.

The gate class is nonempty inside an accessible algebra. Let
$A=\int h(s)e^{isH_p}Ce^{-isH_p}ds$ for a bounded self-adjoint
$C$ and a real $h\in C_c^\infty$. Its first two commutator norms are
bounded by $k_1=\|h'\|_1\|C\|$ and $k_2=\|h''\|_1\|C\|$.
For $V=e^{i\lambda A}$, the energy change $B=V^*H_pV-H_p$ satisfies
\[
 \|B\|\le|\lambda|k_1,\qquad
 \|[H_p,B]\|\le|\lambda|k_2+\lambda^2k_1^2.
\]
Both follow by differentiating
$B=\int_0^\lambda e^{-isA}i[H_p,A]e^{isA}ds$.
Products retain the bounded-commutator property. Consciousness~II's
bounded source preparation and shared-pulse result supply this
smoothed gate class in a selected shared algebra for its stated
translation generator and vacuum~\cite{Liu2026CII}. The exterior
case below transfers it to the interacting current Hamiltonian.
Record-controlled versions
act in the record algebra tensored with that selected algebra.
Thus an acquired record can install a permitted source without
truncating the whole physical spectrum. An exact Weyl displacement
need not satisfy the bounded energy-change condition; it is not
covered by this assertion. Nor does an infinite KMS field have the
finite global rms energy required here.

\begin{proposition}[An exterior source gate for the interacting contact]
\label{prop:exterior-source-gate}
Use the positive current Hamiltonian \eqref{eq:current-positive-form}
or the full fermion contact of Proposition~\ref{prop:full-fermion-contact}
with compact profile $f_a$. Let $C$ be bounded, self-adjoint and localized
in a bounded current interval $I$. Assume it and the free translates
$\alpha_s^f(C)$ for $s\in\operatorname{supp}h$ belong to the selected
accessible algebra, where $h\in C_c^\infty$ is real. Suppose their
localization intervals are disjoint from $\operatorname{supp}f_a$.
For $A=\int h(s)\alpha_s^f(C)ds$ and $V=e^{i\lambda A}$,
\[
 [H_p,A]=[H_f,A],\qquad
 [H_p,[H_p,A]]=[H_f,[H_f,A]].
\]
Both commutators are bounded, with the $k_1,k_2$ estimates above.
The gate is therefore energy-regular for $H_p$, not only for $H_f$.
Its bounded energy change is the same for both Hamiltonians.
Record-controlled versions also meet the autonomous return conditions.
\end{proposition}
\begin{proof}
Write the interaction form as
$H_p-H_f=H_{\rm mech}+\theta\sum_iq_i j_i(f_a)+d_a\theta^2$.
Disjointly supported current smearings have zero symplectic pairing;
their Weyl algebras commute. Thus $A$, and the first two derivative
smearings defining its free commutators, strongly commute with every
$j_i(f_a)$ and act trivially on the mechanical coordinates.
The displayed equalities first hold as polarized forms on the common
form domain. Free translation smoothing gives bounded commutators;
these identities extend through the closed form to
$\operatorname{Dom}H_p$ and imply its invariance under $V$.
Integrating the first identity along $e^{isA}$ gives
$V^*H_pV-H_p=V^*H_fV-H_f$. The second identity gives the required
bounded commutator of this energy change. A finite record control
is a block sum of these same identities with a zero-gap record register.
\end{proof}

This class contains a genuine source preparation. Choose a non-scalar
localized $C$ in a common accessible algebra, with the local vacuum
separating, and a narrow nonnegative $h$ of integral one whose translates
stay exterior. Consciousness~II's bounded-source proposition then gives
$V\Omega$ finite, strictly positive free-field energy for
$0<|\lambda|\|A\|<\pi$. On a finite-rms mechanical state tensored with
the current vacuum, that energy is also the change of $H_p$: the gate
commutes with the contact interaction. With the usual local stress
domain and integrability assumptions, the source perturbation vanishes
outside the translated interval, and a smooth test equal to one there
reads its full positive energy. The autonomous return and QEI bounds
certify this preparation on the actual carrier while retaining the
acquired record and prior references. A non-scalar shared algebra by
itself need not contain such an exterior localized observable; the
localization and permission assumptions are part of this construction.
It supplies a finite source, not a thermal reservoir, a covariant local
history clock or a rule that forms its own permissions.

For an explicit return, take $H_p=N$ on an infinite oscillator and
retain three degenerate bits: a held source bit, its acquired record,
and an output meter. A controlled NOT first acquires the source bit;
that record then rotates oscillator levels $0,1$ by $.2$ or $.9$;
a final controlled NOT reads level $1$ into the output meter.
All higher oscillator levels form a reducing sector. This is an exact
unbounded model, not a cutoff approximation. With initial source
probabilities $.4,.6$ and its old purification retained, a 513-state
clock of mean energy $512$ gives vector error $.118943$ below
$.199530$, and source energy $.3839746172$ against the reference
$.3839484296$. The energy difference is $2.62\times10^{-5}$.
The actual output-record probability is $.3839746172$.
Clock energy decreases by the source gain; total-energy residual is
below $4\times10^{-12}$ in this check. A separate two-level block
calculation agrees with the full-register evolution to
$2.2\times10^{-15}$ for the smaller clocks. Spatial transport,
thermal preparation and the moving support remain distinct from
this acquisition--preparation--reading realization.

\subsection{The contact position can recoil}
\label{sec:moving-quantum-contact}

The contact profile need not follow a trajectory specified from outside.
Let $Y$ be its quantum position, $P_Y=-i\partial_Y$, and $M>0$ its
mass. On the full fermion carrier set
\[
 U_f(a)=e^{-iaH_f},\qquad J_Y=U_f(Y)J(f)U_f(Y)^*.
\]
The mechanical angle and storage coordinate remain internal to the
body. This field translation exists on the full CAR space; it is
not the static infrared dressing of the current.

\begin{proposition}[Finite-mass contact and momentum conservation]
\label{prop:moving-contact}
The form
\[
 \mathfrak q_{\rm mov}=\frac{\|P_Y\Psi\|^2}{2M}
 +\mathfrak q_{\rm mech}+\mathfrak q_f
 +\langle\theta J_Y\rangle+d_f\langle\theta^2\rangle
\]
is closed and nonnegative on the form domain of
$P_Y^2/(2M)+H_{\rm mech}+H_f$. Its self-adjoint Hamiltonian
$H_{\rm mov}$ strongly commutes with $P_{\rm tot}=P_Y+H_f$.
A finite anchor $Z$ of mass $M_0$ can be added with
$P_Z^2/(2M_0)+k_s(Y-Z)^2/2$, $k_s>0$; the conserved momentum
then includes $P_Z$.
\end{proposition}
\begin{proof}
For almost every $Y$, field translation preserves the current form
bound and $H_f$ of Proposition~\ref{prop:full-fermion-contact}.
Integration over $Y$ gives the same equivalent form norms after the
positive position kinetic energy is added. The relative-anchor form
is positive and closed on its natural domain as well.
Simultaneous body and field translation sends $J_Y$ first to
$J_{Y-a}$ and then to $J_Y$. All remaining terms, including the
relative-anchor potential, are invariant. Uniqueness of the closed
form realization implies commutation of $H_{\rm mov}$ with this
translation group and therefore strong commutation with its generator.
\end{proof}

Thus the entire total-momentum distribution is conserved.
These Hamiltonians measure energy above the body and anchor rest
energies; those constant costs must be added in an apparatus ledger.
The translation law is unchanged. On a common
differentiable core the reaction is
\[
 \dot P_Y=-k_s(Y-Z)-\theta\partial_YJ_Y,\qquad
 \dot P_Z=k_s(Y-Z),\qquad \dot H_f=\theta\partial_YJ_Y.
\]
The field-energy change is balanced by the body and anchor. Its
instantaneous sign is not a certificate of exported heat. A positive
record-controlled internal spring can change this dynamics without
breaking translation symmetry; its preparation and gate regularity
remain separate requirements.

The direct-integral unitary $\mathcal W(Y)=e^{-iYH_f}$ gives
$\mathcal W^*P_Y\mathcal W=P_Y-H_f$ and
$\mathcal W^*J_Y\mathcal W=J(f)$. With no anchor, the fixed
total-momentum fiber is
\begin{equation}\label{eq:moving-contact-fiber}
 H_M(p)=H_p+\frac{(p-H_f)^2}{2M}.
\end{equation}
This is the recoil term of a mobile-field canonical transformation,
as in the polaron setting~\cite{LeeLowPines1953}; the current form
here is the full-CAR contact above. For a finite anchor put
$M_t=M+M_0$, $\mu=MM_0/M_t$ and $r=Y-Z$. Its kinetic terms
in a fiber become
\[
 \frac{(p-H_f)^2}{2M_t}
 +\frac{[p_r-(M_0/M_t)H_f]^2}{2\mu}.
\]
Dropping either recoil square changes the model.

For fixed $p$, Eq.~\eqref{eq:moving-contact-fiber} converges in
strong resolvent to $H_p$ as $M\to\infty$. Its positive forms
decrease to $\mathfrak q_p$ on
$\operatorname{Dom}H_p^{1/2}\cap\operatorname{Dom}H_f$.
This intersection contains a core for $\mathfrak q_p$, by equivalence
with the free mechanical/field form. Resolvent variational minimizers
have the closed-form lower bound; approximating a target vector by
that core gives a recovery sequence with recoil penalty tending to
zero. Strict convexity gives norm convergence of the resolvent vectors.
This is not an operator-norm smallness statement. For a static orbit
with the indicated finite fourth energy moment, Duhamel gives
\[
 \|(e^{-itH_M(p)}-e^{-itH_p})\psi\|
 \le\frac1{2M}\int_0^t
       \|(p-H_f)^2e^{-isH_p}\psi\|\,ds.
\]
Finite mean field energy alone does not meet this orbit condition.

For the unbounded harmonic coordinates used here, there is a sharper
limit on a prescribed constant-velocity replacement.
Take $p=Mv$ and remove the free kinetic energy $Mv^2/2$. The fiber
is $H_p-vH_f+H_f^2/(2M)$, not $H_p$.
Put $v_*=a_0/(a_0+d_f)$. For $v<v_*$, the form of $H_p-vH_f$
is closed and nonnegative. For $v_*<v<1$, it is unbounded below
on the smooth core. To see both directions, complete the current
form with coefficient $1-v$:
\[
 H_p-vH_f\ \ge\
 H_{\rm mech,kin}
   +[a_0+d_f-d_f/(1-v)]\theta^2.
\]
For the reverse statement use the compensated compact gauge of
Proposition~\ref{prop:full-fermion-contact}. Its coherent state
has $\langle J(f)\rangle=2d_f\lambda$ and
$\langle H_f\rangle=(d_f+d_b)\lambda^2$. Minimizing over
$\lambda$ gives mechanical leading coefficient
$a_0+d_f-d_f^2/[(1-v)(d_f+d_b)]$. First choose finite $D$ large
enough to make it negative, then send the mechanical packet center
to infinity, with $y=k\theta/(k+K)$. Fixed packet widths and the
conditional gauge derivative have bounded kinetic cost independent
of the center. The energy tends to minus infinity.
At $v=v_*$ the core form is nonnegative, but the previous equivalent
form norm is lost. The full finite-mass $H_{\rm mov}$ remains positive:
its recoil square prevents treating body kinetic energy as an unlimited
constant-velocity source.

Quantum position also changes the locality requirement for the next
thermal preparation. A position wavepacket has tails; its classical
center lying in an interval does not localize the whole protocol.
Choose a smooth position cutoff $\chi$ whose support keeps the full
contact profile in a bounded interval $I$. Let the past interval $A$
contain all free translates of $I$ over the operating time, together
with the other field couplings in the interaction picture. Replacing its interaction
by $\chi\theta J_Y+d_f\chi^2\theta^2$ defines a positive local
comparison Hamiltonian $H_A$. On a common operator-domain trajectory put
\[
 \delta_{\rm loc}=\int_0^t
   \|(H_{\rm mov}-H_A)e^{-isH_{\rm mov}}\Psi\|\,ds.
\]
The retained-reference vector error is at most $\delta_{\rm loc}$.
The thermal comparison in Appendix~\ref{sec:record-thermal-renewal}
with the \emph{actual} record distribution then has trace error at
most $2\delta_{\rm loc}+\sqrt{I_0/2}$: one locality error changes
the joint state, the other changes its record marginal. For source
second-moment norms $M_a,M_l,M_{pl},M_{pa}$ of the actual, local,
local-product and actual-product outputs, respectively, its source
error is at most
\[
 \delta_{\rm loc}(M_a+M_l+M_{pl}+M_{pa})
       +\sqrt{I_0}(M_l+M_{pl}).
\]
All interaction-picture field couplings of the comparison protocol
must lie in $A$;
position-tail probability alone does not bound the generator integral.
The body has genuine quantum recoil, but its Schr\"odinger kinematics
does not supply a relativistic causal support, a clock or a dilaton
stress tensor.

\subsection{Physical locality after unfolding}
\label{sec:physical-locality}

The preceding $H_f$ translates two co-propagating unfolded fields.
Writing physical fields as $\xi_R(x)=\psi_1(x)$ and
$\xi_L(x)=\psi_2(-x)$ changes the bulk momentum to $H_1-H_2$.
The finite-width flavor current before diagonalization becomes
\[
 \int f(x)[\xi_R^\dagger(x)\xi_L(-x)
                    +\xi_L^\dagger(-x)\xi_R(x)]\,dx.
\]
It is a mirror coupling, not a local bulk density. The CAR give
$[J(f),\xi_L(y)]=-f(-y)\xi_R(-y)$, allowing an immediate response
between different positions. Point contact avoids that spatial kernel,
but a profile $f_a(x)=a^{-1}f(x/a)$ at fixed integral has
$d_{f_a}=d_f/a$. The existing energy bounds are not uniform in that
limit. Neither folding nor shrinking supplies a finite local material
body for free.

A local current alternative uses two flavors in each physical direction,
with charges $+1,-1$ after a within-direction flavor rotation. Let
$H_R,H_L\ge0$, $H=H_R+H_L$, $P=H_R-H_L$, and
$d_s=\kappa_s\|f\|_2^2/(4\pi)$, $\kappa_R=\kappa_L=2$.
Both currents $J_R(x),J_L(x)$ now act at the same physical $x$.

\begin{proposition}[Bidirectional current and physical recoil]
\label{prop:bidirectional-current}
The form
\[
 H_p^{\rm bi}=H_{\rm mech}+H+
          \theta[J_R(f)+J_L(f)]+d\theta^2,
 \qquad d=d_R+d_L,
\]
is closed and nonnegative on the free form domain. Translating its
current by $e^{-iYP}$ and adding $P_Y^2/(2M)$ gives a nonnegative
Hamiltonian conserving $P_Y+P$. Its fixed-momentum fiber is
$H_p^{\rm bi}+(p-P)^2/(2M)$.
For the harmonic internal coordinates above and equal $d_R=d_L=d/2$,
$H_p^{\rm bi}-vP$ has a positive closed form for
$|v|<\sqrt{a_0/(a_0+d)}$ and is unbounded below on the smooth core
for $\sqrt{a_0/(a_0+d)}<|v|<1$.
\end{proposition}
\begin{proof}
In physical coordinates the left current has the opposite signs in both
the current commutator and its commutator with $H_L$. For $f=F'$,
$G=J_R(F)-J_L(F)$ therefore satisfies
\[
 e^{-i\theta G}He^{i\theta G}
       =H+\theta(J_R(f)+J_L(f))+d\theta^2.
\]
Compact compensation in each direction handles arbitrary compact $f$,
as in Proposition~\ref{prop:full-fermion-contact}. The same equivalent
form norms and translation proof then apply, with $P$ replacing $H_f$.
Completing the two current forms separately gives
\[
 H_p^{\rm bi}-vP\ge H_{\rm mech,kin}
 +\left[a_0+d-\frac{d_R}{1-v}-\frac{d_L}{1+v}\right]\theta^2.
\]
The bracket is $a_0-dv^2/(1-v^2)$ in the balanced case. For necessity,
use independent compact compensated coherent states in the two
directions. Their minimized leading coefficient is
$a_0+d-d_R^2/[(1-v)(d_R+d_{b,R})]
       -d_L^2/[(1+v)(d_L+d_{b,L})]$.
It is negative for fixed sufficiently spread compensators above the
stated threshold. Harmonic packet centers tending to infinity then
give the negative core sequence, with bounded kinetic cost.
\end{proof}

At equality only core nonnegativity is asserted. The finite-mass
Hamiltonian remains positive with its recoil retained. This alternative
has four complex chiral fields, $c_R=c_L=2$; each actual outlet must
count the fields routed through it. Four input preparations, their
purifiers and the additional control permissions are real resources.
For interval lengths $a,b$ and separation $g>0$ in the product free vacuum,
Appendix~\ref{sec:record-thermal-renewal} uses
$I_0=\frac23\log[1+ab/(g(g+a+b))]$ for these four fields,
twice the two-field value. Both the conditional preparation error
and the source-moment certificate must use this larger $I_0$.
The balanced free carrier and its physical recoil do not yet quantize
a propagating material controller or its complete stress tensor.


\section{Local preparation and correlated-window renewal}
\label{app:thermal-window-proofs}

\subsection{A local preparation can be executed autonomously}
\label{sec:local-thermal-preparation}

The compact preparation unitary $V$ of
Proposition~\ref{prop:finite-thermal-window} is local, but its energy
change need not be bounded. A preparation bit permits a bounded
regularization that controls the field state in energy norm as well
as Hilbert norm. Let $H$ act trivially on that zero-gap bit and set
\[
 C=|1\rangle\langle0|\otimes V+
   |0\rangle\langle1|\otimes V^*,\qquad
 u_0=|0\rangle\Omega,\quad u_1=Cu_0.
\]
Then $C=C^*$, $C^2=I$ and $e^{-i\pi C/2}u_0=-iu_1$.
For a nonnegative even $h\in C_c^\infty(-1,1)$ of integral one,
put $\mu_2=\int s^2h(s)ds$ and
\[
 C_\delta=\int h(s)e^{i\delta sH}Ce^{-i\delta sH}ds,
 \qquad W_\delta=e^{-i\tau C_\delta},\quad\tau=\pi/2.
\]

\begin{proposition}[Local preparation in the energy graph norm]
\label{prop:local-thermal-preparation}
Suppose $F_l(t)=e^{itH}Ce^{-itH}u_l$, $l=0,1$, is twice continuously
differentiable in the $H$ graph norm for $|t|\le\delta_0$.
Let $L_2=\max_{l,t}\|F_l''(t)\|$ and
$G_2=\max_{l,t}\|HF_l''(t)\|$ there. For $0<\delta\le\delta_0$ define
\[
 e_\delta=\mu_2\delta^2L_2/\sqrt2,\qquad
 g_\delta=\mu_2\delta^2G_2/\sqrt2,\qquad
 k_1=\|h'\|_1/\delta,\quad k_2=\|h''\|_1/\delta^2.
\]
The bounded local gate $W_\delta$ has energy-regular constants
$b\le\tau k_1$ and $c\le\tau k_2+\tau^2k_1^2$. Moreover,
\begin{align}
 \|(W_\delta-e^{-i\tau C})u_0\|&\le\tau e_\delta,\nonumber\\
 \|H(W_\delta-e^{-i\tau C})u_0\|
   &\le\tau g_\delta+\tfrac12\tau^2k_1e_\delta.
                      \label{eq:local-preparation-graph}
\end{align}
For a compact diffeomorphism of the vacuum Virasoro module the stated
graph condition holds for free translation $H=H_f$.
For every finite source tolerance, a finite positive autonomous
history clock can therefore execute this preparation while approximating
its local thermal restriction and each specified smooth stress test.
\end{proposition}
\begin{proof}
Time smoothing gives $\|[H,C_\delta]\|\le k_1$ and
$\|[H,[H,C_\delta]]\|\le k_2$, hence the gate constants by
Proposition~\ref{prop:energy-regular-return}.
Evenness of $h$ cancels the first Taylor term. The second-order
remainder bounds $(C_\delta-C)u_l$ by
$\mu_2\delta^2L_2/2$ and its $H$ norm by
$\mu_2\delta^2G_2/2$. The ideal gate path lies in
$\operatorname{span}\{u_0,u_1\}$ with coefficient sum at most
$\sqrt2$. Duhamel gives the first estimate.
For $\zeta\in\operatorname{Dom}H$,
\[
 \|H e^{-itC_\delta}\zeta\|
       \le\|H\zeta\|+t k_1\|\zeta\|.
\]
Inserting this in the same Duhamel integral gives the second estimate,
which tends to zero as $O(\delta)$.
The smooth diffeomorphism representation preserves
$\bigcap_n\operatorname{Dom}L_0^n$, and its action is smooth there
\cite[Sections~5.1--5.2, Proposition~5.1]{FewsterHollands2005}.
M\"obius translation generators act continuously on this domain;
$\Omega$, $V\Omega$ and the required orbit therefore have the stated
graph differentiability. First make the preparation error small by
choosing $\delta$, then, for fixed $\delta$, apply the autonomous
return and smooth-stress certificates. The required clock is finite
at each tolerance, although its resources need not stay bounded as
that tolerance decreases.
\end{proof}

No global vacuum or energy projector enters $C_\delta$. Its support
is the preparation interval enlarged by the short translations.
If this interval stays outside the pinned-current contact, the
commutator equalities above make $W_\delta$ energy regular for the
same interacting $H_p$. Graph convergence also needs its $H_p$-graph
condition on the input; commutator equality alone does not give it.
Schwartz mechanical vectors tensored with the smooth vacuum core
supply an explicit class: the mechanical polynomials, smeared currents
and twice-differentiated exterior orbit act continuously on that core.
Overlapping intervals and more general mechanical states need separate
domain control.

In the free-field execution, field-energy gain is exactly the loss of
the clock-interaction energy $\langle H_{\rm hist}\rangle$; their sum
is conserved. The preparation bit, initial vacuum, allowed local gate,
clock coherence and waveform are supplied resources. Locality here is
locality of a permitted continuum-algebra interaction, not a derived
laboratory actuator or a covariant moving clock. A finite record can
select among these gates while retaining its coherences and references.
Online prediction of its changing thermal clock and physical renewal
remain further requirements.

\subsection{A retained record can select the next thermal window}
\label{sec:record-thermal-renewal}

A later interval of the same field is not an independent reservoir
merely because it has not yet met the contact. The vacuum correlates
disjoint regions; a record made in one can steer the conditional
state in the other. Separation nevertheless gives a quantitative
renewal condition without discarding those correlations.

\begin{proposition}[Record-selected local thermal preparation]
\label{prop:record-thermal-renewal}
Consider $c_f$ independent complex chiral fermions in their vacuum.
Let $A,B$ be split-separated intervals of lengths $a,b$ and gap
$d>0$. The complete past field protocol and acquired record $R$
are a normal local channel on $A$, with initial ancillas independent
of the field. Here $R$ denotes the commuting readout algebra of
the retained register, not an extra physical dephasing.
For each result $r$, let an even $V_r\in\mathcal A(B)$ prepare a
specified thermal restriction $\tau_r$ on $J\subset B$ from the vacuum.
Put
\[
 I_0=\frac{c_f}{6}\log\left(1+\frac{ab}{d(d+a+b)}\right).
\]
After record-controlled preparation,
\begin{equation}\label{eq:record-thermal-distance}
 \frac12\left\|\rho_{RJ}'-
      \sum_r p_r|r\rangle\langle r|\otimes\tau_r\right\|
       \le\sqrt{I_0/2}.
\end{equation}
The norm is that of normal states on the local algebra. The bound
is a probability-weighted conditional error, not a uniform bound
for arbitrarily rare results.

Let $Q_r$ be a smooth source test on $J$ and
$Q=\bigoplus_r V_r^*Q_rV_r$. On the split algebra set
$\alpha=\rho_{RB}$, $\beta=\rho_R\otimes\omega_B$ and suppose
$M_a^2=\alpha(Q^2)$, $M_i^2=\beta(Q^2)$ are finite. Then
\begin{equation}\label{eq:record-source-error}
 |\alpha(Q)-\beta(Q)|\le\sqrt{I_0}(M_a+M_i).
\end{equation}
Neither assertion replaces the actual outside purification by a
product state.
\end{proposition}
\begin{proof}
The past local channel leaves the $B$ marginal equal to $\omega_B$.
Relative-entropy data processing gives
$S(\alpha\Vert\beta)=I(R:B)\le I_{\rm vac}(A:B)$.
For free complex chiral fermions, the multi-interval vacuum entropy
formula gives the displayed $I_0$~\cite[Section~4.1, Eq.~(63)]{CasiniHuerta2009}.
For commuting even observable restrictions it supplies an upper
bound by data processing; the graded fermion fields are not being
replaced by bosonic tensor factors. Measuring an arbitrary binary
effect and applying classical Pinsker gives
$\|\alpha-\beta\|/2\le\sqrt{I_0/2}$.
Controlled preparation and restriction are contractions, proving
Eq.~\eqref{eq:record-thermal-distance}.

For the source bound use the natural-cone representatives
$\xi_\alpha,\xi_\beta$. Spectral Jensen for the relative modular
operator gives~\cite[Section~2, Eqs.~(2.1)--(2.6)]{Araki1976}
\[
 S(\alpha\Vert\beta)\ge
 -2\log\langle\xi_\alpha,\xi_\beta\rangle
 \ge\|\xi_\alpha-\xi_\beta\|^2.
\]
Faithful regularization gives the same bound for finite-entropy
normal states. Expanding the two expectations of $Q$ and applying
Cauchy--Schwarz gives Eq.~\eqref{eq:record-source-error}.
\end{proof}

The estimate can also retain more than the classical record. If the past
protocol outputs a larger controller algebra $\mathcal K$, the same
data-processing argument gives
$S(\rho_{\mathcal KB}\Vert\rho_{\mathcal K}\otimes\omega_B)\le I_0$.
The controller may include its mechanics and references initially
independent of the field. Any subsequent normal joint channel on
$\mathcal K$ and $B$ preserves the local trace bound. Its output
source obeys Eq.~\eqref{eq:record-source-error} with the actual and
comparison output moments. This also controls a later joint contact,
provided its participating algebra is genuinely contained in this
input. It does not include the entire field complement as an
independent reference.

The source moments are explicit inputs. For the record-only estimate, if the past
instrument has local effects $E_r\in\mathcal A(A)$ and
$m_r^2=\omega((V_r^*Q_rV_r)^2)$, then locality and $0\le E_r\le I$
give $M_a^2\le\sum_r m_r^2$ and
$M_i^2=\sum_rp_rm_r^2$. Smooth compact preparations and stress tests
have finite $m_r$. A trace-distance bound alone would not control
an unbounded source.

For a desired local trace error $\epsilon>0$, a sufficient separation is
\[
 d\ge\frac{\sqrt{(a+b)^2+
          4ab/(e^{12\epsilon^2/c_f}-1)}-(a+b)}{2}.
\]
This is the gap required by the vacuum-information certificate,
not a necessary distance for every instrument: an uninformative
record can already be independent at a smaller gap.

For a specified finite sequence, use the actual hull of the past
field operations at each step and a separated preparation support
for the next window. The estimate applies to that actual past,
without resetting it. The local gates can be smoothed and executed
by the positive history clock under the graph-domain conditions
above; its additional gate and source errors are charged separately.
All prior registers and field purifications remain present.
The gap consumes spatial reach and transport time. The record must
reach the preparer through a permitted physical channel; a controlled
unitary does not itself construct that transport. The growing
supports, storage and clock therefore define a finite operating
family, not a fixed-capacity perpetual reservoir.
This local certificate also does not prove that the whole future
field and every old quantum reference are jointly a product state.
Nor does it alone certify the entropy rate of a correlated released field.


\section{Actual diagnostic return and terminal reading}
\label{app:diagnostic-proofs}

\subsection{The reading itself supplies the repair and contact program}
\label{sec:joint-diagnostic-core}

The program need not be a label given to the apparatus. Use three
zero-gap pointer bits for world, sensor and record frame, and three
meter bits. Start the routine meter blank and the other two meters
in independent diagonal states with bit probability $.1$.
Three controlled NOTs copy body parity into the routine meter; two
copy sensor and frame values into their respective noisy meters.
All six bits and the actual held body remain in the core.

The acquired meter word $o$ controls a fixed repair table $a(o)$.
The table is the first two-step decision of Consciousness~III's
continuing-body model~\cite{Liu2026CIII}, with independent coordinate
fault probability $.05$. It depends on the acquired word and known
model, not on the inaccessible true cause. On one history edge use
the unitary
\begin{equation}\label{eq:joint-diagnostic-gate}
 V=\sum_o|o\rangle\langle o|\otimes
       X_{a(o)}\otimes S_{\theta\mathbf1_{a(o)\ne0}},
\end{equation}
where $X_0=I$ and $X_{e_j}=X_j$. The exchange $S$ acts on two
matched-frequency fermion cells initially prepared at temperatures
$.7,1.4$, with $\omega=1$ and $\theta=.7$.
Thus the actual physical repair permission and actual energy transfer
are performed by one program-controlled history edge. Whether the
scorer later calls the repair correct has no effect on the interaction.

Put the five meter gates and this gate on the first six edges of an
eight-edge clock, followed by two identities. Every gate commutes
with reservoir bare energy. The joint positive conserved Hamiltonian
has eight qubit registers, nine clock states and dimension 2304.
At $J=.06$, clock mean energy is $.48$, interaction cutoff $.96$,
and the total spectral cutoff is $2.96$. Initial mean energy is
$1.00197392$. Direct evolution retaining 128 purification columns
agrees with the full gate product to vector error
$2.7\times10^{-14}$; bare energy conservation residual is
$4.5\times10^{-16}$ in this check.
Actual restoration failure is $.02891$, repair probability $.1355$,
and cold-port energy gain is $.00761031$. The latter equals
$\Pr(a\ne0)\sin^2\theta(n_o-n_H)$, not a probability of incorrect
repair supplied as a source command.
The contact angle is a hardware choice: setting it to zero leaves
the logical repair intact. This construction connects actual repair
to an actual source; it does not make that heat transfer a universal
minimum cost of repairing a pointer.
The complete six-bit controller entropy is unchanged. Reservoir
relative entropies are $.0001838782$ and $.0001317734$, and joint
output total correlation is $.0051202854$. They obey
Eq.~\eqref{eq:programmed-ledger} with residual below
$3\times10^{-16}$. Reservoir purification columns are retained for
the vector certificate, not folded into its initially independent
controller in the entropy ledger.

This is a first-interval conserved realization of acquisition returning
to repair and a source. Initial
meter purity, biased noise, reservoir preparation, reference labels
and hardware are supplied and charged separately from the finite
Hamiltonian's energy. The two cells store transferred energy; spatial
release and a covariant supported contact are not implied by this
finite calculation. Its value is that reading, decision, repair and
source no longer need mutually incompatible clocks or an external
selector within the demonstrated interval.

\begin{proposition}[A finite continuing diagnostic return]
\label{prop:finite-continuing-return}
For every specified horizon $L$, the rational pointer model with
$h=1/20$, $\epsilon=1/10$ and the two-step decision admits one finite,
time-independent, nonnegative Hamiltonian. Its completed circuit has
the joint law of all $L$ logical intervals: the same damaged and
repaired body, all actual meter and action records, exact retained
posterior weights, and one record-selected thermal exchange per
interval. No inaccessible fault bit enters the decision circuit.
With a fresh matched-frequency cell pair of gap $\omega$ per interval,
\begin{equation}\label{eq:finite-return-heat}
 Q_L=\omega\sin^2\theta(n_o-n_H)
       \sum_{t=1}^L\Pr(a_t\ne0).
\end{equation}
For a circuit with $m$ padded history edges, clock mean energy is
$Jm$, the total energy cutoff is $2L\omega+2Jm$, and its exact
completion time is $\pi/(2J)$ in units $\hbar=1$.
\end{proposition}
\begin{proof}
Prepare $3L$ independent diagonal fault ancillas and let each control
the corresponding held-body flip. Prepare $3L$ meter bits as above.
The observer's reversible decision uses only the meter words, retained
integer weights and known earlier actions. The rational update has
denominator $800000$ per interval; all weights and comparisons have
finite integer size. A reversible Boolean circuit computes each update
and action, keeping its inputs and garbage. The action controls the
held-body flip and exchange with the next prepared thermal pair.
Every gate commutes with the sum $H_0$ of cell energies, because the
logic is degenerate and each exchange has matched gaps. Put the gates
on a positive history clock. The prefix-unitary conjugation gives
$H_0+JmI+2JJ_x$, proving positivity, bare-energy conservation and
full-reference endpoint completion. Each unused pair is still jointly
independent of the acquired program at its contact. Conditional energy
transfer is zero or $\omega\sin^2\theta(n_o-n_H)$, which proves
Eq.~\eqref{eq:finite-return-heat}.
\end{proof}

For the first three logical intervals, the expected number of repairs is
$.44635607$ and $Q_3=.02506944$ at the stated cell settings.
There are nine readings, nine prepared fault bits and six thermal cells.
Keeping all preceding weight vectors requires at most 984 bits in this
case, before records, arithmetic workspace and reversible garbage.
The general saved-weight cost grows quadratically with $L$.
This construction removes the need for a new externally chosen circuit
at every diagnostic return. Its logical intervals are not equally spaced
proper-time samples of a continuously acting bath: only the completed
joint law is asserted. All ancillas and cells are prepared before the
finite run, all outputs remain in its ledger, and the clock can reverse
after completion. It proves neither a fixed-capacity cyclic apparatus
nor a covariant moving contact. These are distinct renewal and material
implementation questions, not failures of the finite return.

\subsection{An actual terminal record and probe replacement}
\label{sec:terminal-meter}

Let probe $P$ retain any old reference, record $R$ start in $|0\rangle$,
and a paid blank $B$ start in $|+\rangle$. Write
$C_Y=P_-^Y\otimes I_R+P_+^Y\otimes X_R$, then apply the true
$\operatorname{SWAP}_{PB}$. The record-one probability is
$\operatorname{tr}(P_+^Y\rho_P)$. The final probe is $|+\rangle$,
independent of all old references and the record; its old correlations
have moved into $B$ and $R$, not disappeared. Positive generators
$gP_+^Y\otimes(I-X_R)$ and $g(I-\operatorname{SWAP}_{PB})$ implement
these gates at duration $\pi/(2g)$. A fermionic exchange is not a
substitute for this true SWAP.

The physical probe may continue to acquire a phase while the meter
executes. Move the $r$ nonidentity gates of the clock in
Appendix~\ref{sec:conserved-core} to its last $r$ edges. All earlier
edges are identities. Put $T=\pi/(2J)$, choose $t_0<T$ and define
\[
 P_0=\Pr\{\operatorname{Bin}(m,\sin^2Jt_0)\ge m-r\},\qquad
 A=J\sum_{j=m-r+1}^m\sqrt{j(m-j+1)}.
\]

\begin{proposition}[Terminal read under a retained phase coupling]
\label{prop:terminal-meter}
Let the actual probe coupling be $\lambda d(t)Z_P$, with
$|d(t)|\le D$. Relative to writing the probe's actual phase at $t_0$
and replacing it at $T$, the full joint half trace distance obeys
\begin{equation}\label{eq:terminal-meter-bound}
 d_{\rm joint}\le
  2At_0\sqrt{P_0}+2\sqrt{P_0}+\lambda D(T-t_0).
\end{equation}
All old references, the blank, record and clock are included.
For any specified finite trajectory and positive error tolerance,
finite $m$ and $t_0<T$ suffice.
\end{proposition}
\begin{proof}
Before $t_0$, compare with the identity-edge clock and the actual free
probe phase. Only the last $r$ edges differ. Their action on that
reference is bounded by $2A\sqrt{P_0}$, since the binomial upper tail
increases up to $T$. Duhamel gives the first term. At $t_0$, dressing
the reference with the gate prefixes changes its joint state by at
most $2\sqrt{P_0}$. From there to $T$, omit only the bounded phase
coupling, at cost $\lambda D(T-t_0)$; the compiled clock completes
the meter and SWAP exactly. First choose $t_0$ near $T$, then increase
$m$ to suppress the tail.
\end{proof}

The positive clock has mean energy $Jm$, not zero. Add the probe/blank
offset $2\epsilon_P I$, with $2\epsilon_P>\lambda D$, to keep the
complete internal Hamiltonian positive. For $m=10^5$, $r=2$, $T=80$,
$t_0=78$ and $\lambda D=.00032$, a sufficient error is below $.000641$;
clock mean energy is $1963.4954$. This cannot be inserted into a
mass-10 support without changing that support. A direct 17-state clock
check gives joint distance $.0034887$ below its certificate $.0046290$.
The construction renews the probe over a finite window. It consumes a
coherent blank and a record slot and does not reset the entire reader
or its clock. Its phase frame and allowed gate family remain inputs.
HAFF~D also solves a smaller, 129-state meter clock jointly with a
source-matched mean moving support~\cite{Liu2026HAFFD}. Its energy is
included in the initial supported mass, and the actual record depends
on the dilaton trajectory. That mean-field calculation does not include
the contact reservoir's microscopic degrees or stress. For the specified
$C_Y$, SWAP and coherent blank, an antiunitary symmetry makes the mean
probe force vanish throughout the read, although the force operator
and its fluctuations remain nonzero.

\section{Local linear-current control of stress flows}
\label{app:linear-stress}

Consider the neutral current on a circle of length $2\pi$, with
\[
 J_n=\sqrt n\,a_n,\quad J_{-n}=\sqrt n\,a_n^*,\quad J_0=0,
 \qquad H=L_0=\sum_{n\geq1}n a_n^*a_n.
\]
Our Fourier convention is $f_n=(2\pi)^{-1}\int_0^{2\pi}
f(\theta)e^{-in\theta}\,d\theta$. The same convention applies to the
real current smearing $h$. The primitive interaction is $-B j(h)$,
where $j(h)=\sum_{n\geq1}\sqrt n(h_{-n}a_n+h_na_n^*)$ and $B$
is a finite controller coefficient. Fix $h\in C_c^\infty(I,\mathbb R)$
on a proper arc $I$, with $h_n\ne0$ for $n\geq1$.
Such an $h$ exists: in the Fr\'echet space of real smooth functions
supported in a fixed compact subarc, each equation $h_n=0$ defines a
proper closed subspace, so the intersection of their open dense
complements is nonempty. For a prescribed finite cutoff $N$, a
nonnegative sufficiently narrow bump already gives $h_n\ne0$ for
$1\leq n\leq N$; this alternative makes the probe depend on $N$.
No uniform lower bound on $|h_n|$ is assumed.

\begin{proposition}[A neutral-current stress compiler]
\label{prop:current-stress-compiler}
Suppose signed local linear-current coupling, finite degenerate-helper
rotations and record phase gates, calibrated clock preparation and
finite readout are available, with coherent control intervals.
For every real smooth $f$, $t\in\mathbb R$, $E<\infty$ and
$\delta>0$, a finite positive autonomous program implements
$e^{itT(f)}$ within joint half-trace distance $\delta$, uniformly over
inputs with $\operatorname{tr}(H\rho)\leq E$ and arbitrary retained
references. Its helpers and original clock are included in the
comparison with their designated final states. Here
\[
 T(f)=\frac12\sum_{p,q\in\mathbb Z}
                 f_{-p-q}:J_pJ_q:
\]
is the $c=1$ Sugawara stress tensor. The program uses no field reset,
heralding or primitive quadratic coupling.
\end{proposition}

\begin{proof}
We first include the field's free propagation. For a finite quadrature
$Q(v)=\sum_{n\leq N}(v_na_n^*+\overline{v_n}a_n)$ put
\[
 G_n=\frac{v_n}{\sqrt n h_n},\qquad
 g_v(u)=\frac1{2\pi}\sum_{n\leq N}
       (G_ne^{-inu}+\overline{G_n}e^{inu}),\quad 0\leq u\leq2\pi.
\]
The Hamiltonian
\[
 H-sg_v(u)A j(h)+s^2g_v(u)^2 C_h,
 \qquad C_h=\sum_{n\geq1}|h_n|^2,
\]
is positive by completion of the oscillator squares, with shifts
$sg_v(u)Ah_n/\sqrt n$ and $A=X$ or $Y$.
In the interaction picture $[j_u(h),j_w(h)]$ is a scalar.
Thus the Magnus series terminates after its scalar second term, and
$e^{-2\pi iH}=1$ gives the endpoint gate
$e^{i\phi_s}e^{isA Q(v)}$.
The phase $\phi_s$ is computed from the scalar commutator and the
completion integral. Subsequent pulses start at integer periods and
therefore act on the same $Q(v)$. Changing $s$ to $-s$ pays for the
inverse without reversing $H$. A $Y$ pulse may instead use a paid
$Z$ rotation conjugating the $X$ coupling.
For coherent record branches, the known phases are removed by
record phase gates.

The exact trigonometric comparator does not have smooth compact-time
switching. Multiplying it by an all-order-flat endpoint taper
$\chi_\eta$ makes the physical coupling smooth and gives
$\|(1-\chi_\eta)g_v\|_1\leq2\eta\|g_v\|_\infty$.
Writing $C_1=\sum n|h_n|^2$ and $D_h=\sum n^2|h_n|^2$, we have
\[
 \|j(h)\psi\|\leq2\sqrt{C_h}\|H^{1/2}\psi\|
                    +\sqrt{C_1}\|\psi\|,
 \qquad
 \sqrt{\langle H\rangle_u}\leq\sqrt E+
                  \sqrt{D_h}|s|\|g_v\|_1.
\]
The second bound follows directly from the controlled displacement.
It includes the transient excitation of all modes, not just the
selected endpoint modes. Duhamel's formula then makes each taper
error arbitrarily small in joint vector norm, including old references.
For a finite word of such pulses the energy bound is
$(\sqrt E+\sqrt{D_h}L)^2$, where
$L=\sum_{\rm slots}|s|\|g_v\|_1$.

Next use the current-factor cutoff
\[
 T_N(f)=\frac12\sum_{0<|p|,|q|\leq N}
                       f_{-p-q}:J_pJ_q:.
\]
This is a selfadjoint finite oscillator quadratic. Finite occupation
vectors are analytic for a finite quadratic, hence form an essential
selfadjointness core. Let $P_M=1_{[0,M]}(H)$, $M$ an integer. It has
finite rank. For $N>M$, the omitted number-preserving terms have
annihilator index $m\leq M$ and creator index $n>N$. Their norm on
$P_M$ is bounded by $\sqrt{nM}|f_{n-m}|$.
Omitted pair annihilators vanish; pair creators have norm at most
$M+2$. Consequently
\begin{align*}
 \|[T(f)-T_N(f)]P_M\|&\leq b_N(M)\longrightarrow0,\\
 b_N(M)&=\sqrt M\sum_{m=1}^M\sum_{n>N}\sqrt n|f_{n-m}|
       +\frac{M+2}{4}\sum_{k>N}k^2|f_k|.
\end{align*}
For the pair term we used
$\sum_{n=1}^{k-1}\sqrt{n(k-n)}\leq k^2/2$.
These absolutely convergent sums identify the limit on the
finite-energy core as the Sugawara $T(f)$.
Carpi--Weiner's essential selfadjointness theorem for real smooth
stress smearings~\cite[Theorem~4.4]{CarpiWeiner2005} supplies the selfadjoint closure of the limit; it
does not substitute for the current-factor tail estimate.
On this common core, the resolvent identity and density of
$(T(f)-z)D_{\rm fin}$ give strong resolvent convergence.
Hence $e^{itT_N(f)}\to e^{itT(f)}$ strongly, and
\[
 a_N(M,t)=\|[e^{itT_N(f)}-e^{itT(f)}]P_M\|\longrightarrow0.
\]

The uniform reference estimate is obtained only once, at the input
of the whole program. For isometries $V,W$, purification and the
$P_M$ decomposition give
\[
 D_{\rm joint}(V\rho V^*,W\rho W^*)
 \leq\|(V-W)P_M\|+2\sqrt{\frac E{M+1}}.
\]
It is independent of reference dimension. Thus the factor-cutoff
unitaries converge uniformly over the stated mean-energy class.
The core-tail bound is not itself a rate for $a_N(M,t)$, since a
stress-evolved vector need not stay in $P_M$.

Write $r=(q_1,p_1,\ldots,q_N,p_N)$. A real symmetric matrix $K_N$
gives
\[
 T_N=\tfrac12r^TK_Nr-c_N,
 \quad c_N=\tfrac{f_0}{2}\sum_{n\leq N}n,
 \quad T_N=\tfrac12\sum_{\ell\leq2N}\kappa_\ell Q_\ell^2-c_N
\]
after real orthogonal diagonalization. This is an algebraic
decomposition, not a free optical transformation: each $Q_\ell$
is selected directly by the local pulse above.
Finite quadratic Trotter products converge on oscillator Schwartz
vectors and hence strongly, including after restriction to $P_M$.
Each square exponential is compiled with the same retained helper
by four Pauli pulses. For an exponent of magnitude $|\alpha|$,
$n$ blocks and $s=\sqrt{|\alpha|/(2n)}$, the second-moment estimate is
\[
 D_{\rm joint}\leq(13\sqrt2)^{1/3}
       \sqrt{|\alpha|m_2}\,n^{-1/6},\qquad
 m_2=\operatorname{tr}(\rho Q_\ell^2)\leq4\operatorname{tr}(H\rho)+1.
\]
To see the moment bound, on a spectral value $q$ set
$C_\vartheta=e^{i\vartheta X}e^{i\vartheta Y}
 e^{-i\vartheta X}e^{-i\vartheta Y}$.
Taylor's integral remainder gives
$\|C_\vartheta-e^{-2i\vartheta^2Z}\|\le13|\vartheta|^3$:
the third derivative of the four-factor product has norm at most
$64$, and the quadratic exponential contributes at most
$2\vartheta^4$ for $|\vartheta|\le1$; outside that interval the
unitary distance is at most $2$.
Telescoping $n$ blocks yields a spectral error
$\min\{2,A|q|^3\}$, with
$A=13|\alpha|^{3/2}/(2\sqrt{2n})$.
The inequality
$\min\{4,A^2q^6\}\le4^{2/3}A^{2/3}q^2$ and the spectral theorem
give the displayed joint bound, retaining the helper and reference.
The sign is chosen by the Pauli order; all inverse pulses are paid.
Moments for a finite word are evaluated on its ideal prefixes, whose
oscillator moments are finite on $P_M$.
Telescoping the complete word against those prefixes avoids both
fresh-helper resets and a circular block-count estimate.

Finally approximate the smooth finite schedule in $L^1$ by constant
slots with rational history-clock thresholds. Every slot has the
positive field Hamiltonian $H-B_jj(h)+C_hB_j^2$.
Install the finite rotations and record phase gates on the original
clock path; these gates commute with the degenerate data bare energy.
The resulting sum of the positive history clock, data energy and
completed field squares is a positive autonomous Hamiltonian on
the common field-energy domain.
For $G=\max_j\|B_j\|$ and $D_3=\sum n^3|h_n|^2$, its field rms obeys
\[
 D^+\|H\psi_u\|\leq2G\sqrt{D_h}\sqrt{\|H\psi_u\|}
                         +G\sqrt{D_3}.
\]
The bound is independent of clock size. The finite-boundary
history-clock estimate therefore makes the complete endpoint error
arbitrarily small on $P_M$, with the same original clock and all
record branches retained. Choosing successively $M,N$, Trotter
length, square-block counts, taper/slot precision and clock size
allocates any prescribed total $\delta$.
\end{proof}

The resources can grow sharply. With
$\Gamma(v)=\|g_v\|_1\leq2\sum_{n\leq N}|v_n|/(\sqrt n|h_n|)$,
a square exponent of size $|\alpha|$ uses field duration $8\pi n$,
area $2\sqrt{2n|\alpha|}\Gamma(v)$ and positive completion integral
\[
 \frac{2C_h|\alpha|}{\pi}
                  \sum_{n\leq N}\frac{|v_n|^2}{n|h_n|^2}.
\]
Helper rotations, phase gates and readout have their own finite
budgets; rotation slots can last $2\pi$ so the field drift is identity.
For an autonomous clock of duration $T$ and size $m+1$,
$J=\pi/(2T)$, its initial mean energy is $mJ$, norm $2mJ$ and
largest edge coefficient at most $J(m+1)/2$.
Preparation, drive bandwidth and retained storage are additional
resources. These are sufficient finite bounds, not optimal costs or
net-heat identities.

The uniform channel statement assumes only mean energy.
Expectations of unbounded stress require separate graph-norm control:
for smooth $g$, $\|T(g)\psi\|\leq C_g\|(H+1)\psi\|$.
Small joint channel error alone does not certify source convergence
on the whole mean-energy class. The finite program's field moments
also do not include a covariant stress tensor for clock and support.

The target preserves the Virasoro vacuum subspace, although the
individual current pulses and finite $T_N$ approximants can leave it;
the endpoint leakage is part of the error above.
The compiler concerns the neutral current sector. Its controls are
not invariant under the $SO(3)$ action in the $SU(2)_1$ carrier, so
an always-on selecting noise needs a separate control/noise
certificate. The circle, calibrated port, coherent interval and
controller preparation remain physical inputs.

